% Needed for proper arXiv processing, it seems
\pdfoutput=1

\newif\ifdraft\draftfalse % draft = comments
\newif\ifanon\anonfalse   % anon = light double-blind reviewing (POPL)
\newif\iffull\fullfalse    % full = includes things that were cut from
                          % conf. proceedings only because of space
                          % might be turned true for a journal submission
\newif\iflongrefs\longrefsfalse % Long references (e.g. for journal)
\newif\ifbackref\backreffalse % backref option for hyperref
                              % useful for shrinking references
\newif\ifsooner\soonerfalse
\newif\iflater\laterfalse
\newif\ifcamera\cameratrue     % Camera-ready version
\newif\ifcheckpagebudget\checkpagebudgetfalse
\newif\ifconference\conferencetrue
\newif\ifpages\pagesfalse      % Break pages on sections
% !!! PLEASE DON'T CHANGE THESE !!! INSTEAD DEFINE YOUR OWN texdirectives.tex !!!
\newcommand{\xxx}{}
\makeatletter \@input{texdirectives.tex} \makeatother

\documentclass[runningheads]{llncs}

\usepackage{afterpage}
\usepackage{color}
\usepackage{soul}
\usepackage{amssymb,amsmath,amsfonts}
\usepackage{extarrows}
%\usepackage{mathptmx} %-- BROKEN! Makes parens in math mode dissapear!
\usepackage{version}
\usepackage{xspace}
\usepackage{graphicx}
\usepackage{listings}
\usepackage{wrapfig}
\usepackage{ifpdf}
\usepackage{semantic}
\usepackage{mathpartir} % this takes over some things from semantic
\newcommand{\tightlabel}[2]{\makebox[1pt][l]{\raisebox{+1em}{\hspace{0.5em}(#1)}}}
\newcommand{\tightlabelup}[2]{\makebox[1pt][l]{\raisebox{+2.4em}{\hspace{0.5em}(#1)}}}
\usepackage[numbers,sort]{natbib}
\newcommand\citepos[1]{\citeauthor{#1}'s\ \citeyear{#1}}
% \bibpunct();A{},
% \let\cite=\citep
\usepackage{stmaryrd}
\usepackage{subfigure}
\usepackage{multirow}
\usepackage{proof}
\usepackage{framed}
\usepackage{caption}
\usepackage{mathpartir}
\usepackage{enumitem}
\usepackage{tikz}
\usetikzlibrary{positioning}% To get more advances positioning options
\usetikzlibrary{arrows}% To get more arrow heads
\usepackage[T1]{fontenc}
\usepackage[utf8]{inputenc}

% Fixing LNCS nonsense to get a PDF outline
\setcounter{tocdepth}{3}

\newcommand{\kw}[1]{\text{\ls|#1|}}
\newcommand\Goal[3]{\mbox{\ensuremath{#1 \vdash \kw{#2} : \kw{#3}}\xspace}}
\newcommand\VGoal[2]{\mbox{\ensuremath{#1 \vDash \kw{#2}}\xspace}}

% Colors
\definecolor{dkblue}{rgb}{0,0.1,0.5}
\definecolor{dkgreen}{rgb}{0,0.4,0}
\definecolor{dkred}{rgb}{0.6,0,0}
\definecolor{dkpurple}{rgb}{0.7,0,1.0}
\definecolor{purple}{rgb}{0.9,0,1.0}
\definecolor{olive}{rgb}{0.4, 0.4, 0.0}
\definecolor{teal}{rgb}{0.0,0.4,0.4}
\definecolor{azure}{rgb}{0.0, 0.5, 1.0}
\definecolor{gray}{rgb}{0.5, 0.5, 0.5}
\definecolor{dkgrey}{rgb}{0.2, 0.2, 0.2}
\definecolor{lilac}{rgb}{0.70, 0.04, 0.08}
\definecolor{dkyellow}{rgb}{0.5,0.5,0}

\usepackage{hyperref}
\renewcommand{\sectionautorefname}{\S}
\hypersetup{breaklinks=true}
\hypersetup{colorlinks=true}
\hypersetup{linkcolor=black}
\hypersetup{citecolor=black}
\hypersetup{filecolor=black}
\hypersetup{urlcolor=dkblue}
\newcommand\maybecolor[1]{\color{#1}}

\newcommand\mypara[1]{\vskip 0.03in \noindent{\textbf{\itshape{#1}}}\quad}

\newcommand{\aset}[1]{{\ensuremath{\{#1\}}}}

\def\lstlanguagefiles{lstfstar.tex}
\lstset{language=fstar}
\let\ls\lstinline

\def\Snospace~{\S{}}
\def\sectionautorefname{\Snospace}
\def\subsectionautorefname{\Snospace}
\def\subsubsectionautorefname{\Snospace}


\newcommand\fstar{F$^\star$\xspace}
\newcommand\emf{EMF$^\star$\xspace}
\newcommand\metafstar{Meta-F$^\star$\xspace}
\newcommand\lowstar{Low$^\star$\xspace}

% Comments
\newcommand{\comm}[3]{\ifcheckpagebudget\else\ifdraft{{\color{#1}[#2: #3]}}\fi\fi}
\newcommand{\nik}[1]{\comm{dkpurple}{Nik}{#1}}
\newcommand{\ch}[1]{\comm{teal}{CH}{#1}}
\newcommand{\aseem}[1]{\comm{magenta}{Aseem}{#1}}
\newcommand{\guido}[1]{\comm{blue}{Guido}{#1}}
\newcommand{\vd}[1]{\comm{dkgreen}{VD}{#1}}
\newcommand{\tahina}[1]{\comm{dkred}{Tahina}{#1}}
\newcommand{\jonathan}[1]{\comm{pink}{Jonathan}{#1}}
\newcommand{\cpc}[1]{\comm{brown}{Cl\'ement}{#1}}
\newcommand{\zp}[1]{\comm{lilac}{Zoe}{#1}}
\newcommand{\da}[1]{\comm{dkyellow}{Danel}{#1}}

\newcommand*{\EG}{e.g.,\xspace}
\newcommand*{\IE}{i.e.,\xspace}
\newcommand*{\ETAL}{et al.\xspace}
\newcommand*{\ETC}{etc.\xspace}
\begin{document}
%
\def\titlestr{\metafstar: Proof Automation with \\SMT, Tactics, and Metaprograms}
% \def\titlestr{\metafstar: Tactics for Program Verification}
% ??
\title{\titlestr}
%
%\titlerunning{\metafstar}
% If the paper title is too long for the running head, you can set
% an abbreviated paper title here

\ifpages
\let\oldsection\section
\renewcommand\section{\clearpage\oldsection}
\fi

\author{
    Guido Mart\'inez\inst{1,2} \and
    Danel Ahman\inst{3} \and
    Victor Dumitrescu\inst{4} \and
    Nick Giannarakis\inst{5} \and
    Chris Hawblitzel\inst{6} \and
    C\u{a}t\u{a}lin Hri\c{t}cu\inst{2} \and
    Monal Narasimhamurthy\inst{7} \and\\
    Zoe Paraskevopoulou\inst{5} \and
    Cl\'ement Pit-Claudel\inst{8} \and
    Jonathan Protzenko\inst{6} \and\\
    Tahina Ramananandro\inst{6} \and
    Aseem Rastogi\inst{6} \and
    Nikhil Swamy\inst{6}
}
%
\authorrunning{Mart\'inez et al.}
%
\institute{
    $^1$CIFASIS-CONICET \quad
    $^2$Inria Paris \quad
    $^3$University of Ljubljana \\
    $^4$MSR-Inria Joint Centre \quad
    $^5$Princeton University \quad
    $^6$Microsoft Research \quad
    $^7$University of Colorado Boulder \quad
    $^8$MIT CSAIL
}
%
\maketitle              % typeset the header of the contribution
%
\begin{abstract}


%    Program verifiers generally rely on automated provers such as SMT
%    solvers for discharging verification conditions. While this works
%    well fairly often, for facts in theories or logics that provers
%    cannot efficiently handle users need to provide many ad-hoc hints to
%    guide the solver, resulting in verification brittleness and a waste
%    of human effort.
%
%\smallskip
%    This paper introduces \metafstar, a tactics and metaprogramming
%    framework for the \fstar program verifier.
%    The main novelty of \metafstar is allowing to use tactics and 
%    metaprogramming to discharge assertions not
%    solvable by SMT, or to just simplify them into the well-behaved SMT
%    fragments. Plus, \metafstar can be used to
%    generate verified code automatically.

    We introduce \metafstar, a tactics and metaprogramming
    framework for the \fstar program verifier.
    The main novelty of \metafstar is allowing the use of tactics and 
    metaprogramming to discharge assertions not
    solvable by SMT, or to just simplify them into well-behaved SMT
    fragments. Plus, \metafstar can be used to
    generate verified code automatically.

\smallskip
    \metafstar is implemented as an \fstar \emph{effect}, which, given
    the powerful effect system of \fstar{}, heavily increases code reuse
    and even enables the lightweight verification of metaprograms.
    Metaprograms can be either interpreted, or compiled to efficient
    native code that can be dynamically loaded into the \fstar type-checker
    and can interoperate with interpreted code.
    Evaluation on realistic case studies shows that \metafstar
    provides substantial gains in proof development, efficiency, and robustness.


% %% Problem
% Tactics have been used to script proofs in interactive
% theorem provers for decades.
% %
% For nearly as long, semi-automated
% program verifiers have used tools like SMT solvers to improve proof
% automation.
% %
% Despite their strengths, for verification in the large, both
% approaches also have shortcomings:
%     (1) building application-specific tactics requires great effort;
%     % and automated provers are not easily exploitable
%     (2) SMT solvers only efficiently automate a limited class of logics and theories.
% %
% On the theories that are not efficiently SMT-automatable
% (e.g. non-linear arithmetic) users must provide ad-hoc hints,
% resulting in brittleness and a waste of human effort.

%% Solution
% \metafstar is a new metaprogramming
% framework for \fstar which permits to profit from the state-of-the-art
% automation in SMT solvers, while using tactics to stay clear of their
% weaknesses, increasing overall proof robustness.
% %
% Dually, this new style of proof reduces the effort of the tactic
% writer, as tactics need to automate only those parts of a proof that are
% difficult for an SMT solver.
% %
% Beyond tactics, \metafstar can construct and reflect on the
% syntax of \fstar programs, while safely manipulating the internal
% state of its type-checker.
%
% This allows scripting the construction of effectful programs along
% with their proofs of correctness, significantly reducing the need for
% user interaction in some proofs.


%% How the solution works
%     We \emph{embed} \metafstar in \fstar as an \emph{effect}, instead
%     of creating a separate language. This heavily increases reuse and enables interesting
%     new possiblities, such as a (lightweight) verification of metaprograms.
% %
% Unlike programs, metaprograms are meant to be executed at typechecking-time. In
% favour of efficiency, we design and implement a new
% mechanism for extending the \fstar compiler with binary plugins
% compiled from verified source (meta)programs.
% %
% Our experimental evaluation shows that with metaprogram plugins,
% performance and stability of hybrid SMT/metaprogramming proofs
% can greatly exceed
% that of SMT-only proofs.

\ifcamera
\keywords{
Tactics \and
Metaprogramming \and
Program verification \and
Verification conditions \and
SMT solvers \and
Proof assistants}
\fi

\end{abstract}
%
%
%

% \ifdraft{
% \section*{Meta section}
% \guido{Brought back this page for higher-level points, and playing
% devil's advocate with potential critiques that we might want to defend
% against, but I'm not super sure on how to, PLEASE READ. If possible,
% ``reply'' to the critique in the paper body and cross it out here.}
% \begin{enumerate}
% 
%     \item Potential Evil Reviewer \#1: ``you say it's a user defined effect but it only works because of trusted primitive operations''.
% 
%     \item Potential Evil Reviewer \#2: ``you do not compare to SMT tactics \url{https://z3prover.github.io/api/html/ml/Z3.Tactic.html}''
% 
%     \item Potential Evil Reviewer \#3: ``system X also generates typing derivations''
%         \subitem \guido{Is there another metaprogramming system over undecidable typing? Is this new or not?}
% 
% 
%     % \item \ch{Can we try to use Clement's emacs-mode-based coloring,
%     %           or is that as annoying and full of nonsense as the listings package?}
%     %   \cpc{I can port the document fairly easily; but maybe it's better to wait until after the deadline?}
% 
%     % \item \guido{Thinking out loud, maybe we should ``abstract away'' from the
%     %     fact that we use an SMT solver and just call it a ``sink'' or something
%     %     like that? I think our hook and all are VC-related, but not necesarilly
%     %     SMT-related.}\ch{In \fstar{} the sink is always an SMT solver, no?}
%     %     \guido{Yes, but we don't rely on that. If we change it to Lean or whatever
%     %     we can still use these tactics. They're more dependent on the VC generation
%     %     than the actual prover, IMO.}
% \end{enumerate}
% 
% % \ch{Case studies (2018-03-01):
% % - Cannonizer, bit vectors (can use early in the paper to illustrate stuff)
% % - Pattern matcher (but some use of different kind of quotation)
% % - Optimizing VCs (from VALE, optimizing register map updates) [Nik]
% % - Separation Logic (optimizing memory, deep vs shallow) [Aseem,Danel]
% % - Tactics for program transformation:
% %   + loop unrolling, commuting lambda and match (exist already) [Tahina];
% %   + Karthik inspired example: transforming from one monad to another [Danel];
% %   + In-place imperative pretty printer for arith expressions
% %     from pure printer that computes the length [Tahina]
% % - (Maybe) Adding trace to tail-recursive function [Guido]
% % }
% 
% % Arithmetic running example:
% % 1. interactive
% %    - assert_by_tactic inside a simple program, under conditional (factorial?)
% %    - prove goal using no automation, just basic tactics and lemmas
% % 2. try / fail / unification only
% %    - loop blindly trying to apply lemma sequences until something succeeds
% % 3. reflective
% %    - current implementation: using pointwise, and reflecting syntax
% %                              in order to decide which lemmas to apply
% %                              and to do a bit of massaging via pointwise too
% %    - pattern matching
% % 4. deep
% %    - reflect on syntax, check that going back gives the same denotation,
% %      apply "optimization" steps that are proven to preserve the denotation
% \pagebreak
% }\fi

%% %% 2012 ACM Computing Classification System (CSS) concepts
%% %% Generate at 'http://dl.acm.org/ccs/ccs.cfm'.
%% \begin{CCSXML}
%% <ccs2012>
%% <concept>
%% <concept_id>10011007.10011006.10011008</concept_id>
%% <concept_desc>Software and its engineering~General programming languages</concept_desc>
%% <concept_significance>500</concept_significance>
%% </concept>
%% <concept>
%% <concept_id>10003456.10003457.10003521.10003525</concept_id>
%% <concept_desc>Social and professional topics~History of programming languages</concept_desc>
%% <concept_significance>300</concept_significance>
%% </concept>
%% </ccs2012>
%% \end{CCSXML}

%% \ccsdesc[500]{Software and its engineering~General programming languages}
%% \ccsdesc[300]{Social and professional topics~History of programming languages}
%% End of generated code


%% Keywords
%% comma separated list
% \keywords{keyword1, keyword2, keyword3}
%% \keywords are mandatory in final camera-ready submission

%% \maketitle
%% Note: \maketitle command must come after title commands, author
%% commands, abstract environment, Computing Classification System
%% environment and commands, and keywords command.

% \ch{abstract: could consider adding more details: LCF architecture,
%   core proved sound} -- gory detail
% \ch{Also, tactics don't have to produce proof terms for
%   proof-irrelevant positions} -- gory detail

%
% \guido{tentative: We take the novel\ch{Nice try, but what about Mtac,Idris,Lean?
%     They all use monads to represent a computational effect, so at a
%     high enough level they do the same thing. In intro we were more
%     careful at positioning this, and should do the same here. Or at
%     otherwise just drop the novel word.}
% %
% approach of modelling metaprogramming
% as a computational effect within our existing framework of \fstar{}.
% %
% This has a number of benefits,
%    from reusing many language constructs and syntax,
%     to being able to specify metaprograms using the existing WP calculus
% %
%     mechanism\ch{This is for a most part not a feature we want to call
%       out since we're not into the same game as Mtac, and definitely
%       not a feature we want to call out given how much pain it's been
%       causing us even if we {\bf weren't} trying to use this. Things could
%       change if we find a good use case for this in the future (e.g,
%       a new form of effectful proofs by reflection)}
%
    % \guido{tentative: Another particularity of our solution is working
    % on goals generated by a verification condition (VC) generator,
    % which are not explicitly written by the user nor exactly predictable,
    % and we allow the use of tactics to preprocess them before
    % being fed to the SMT solver,
    % without needing to fully proving them.
    % %
    % Our tactic library therefore
    % has a semi-automated flavour, with many ``hammer''-like tactics.
    % }


\section{Introduction}

% \ch{An easier to win competition is against Dafny et al, than against
%   Coq, Isabelle, and Lean. Bringing some of the goodness of Lean to
%   Dafny is already a great feat, why not start and focus on that?}

Scripting proofs using tactics and metaprogramming has a long
tradition in interactive theorem provers (ITPs), starting with Milner's
Edinburgh LCF~\cite{gordon-edinburghLCF:1979}.  In this lineage,
properties of \emph{pure} programs are specified
% and proved -- duplicates "and proofs" below
in expressive higher-order (and often dependently typed) logics,
and proofs are conducted using various imperative programming
languages, starting originally with ML.

Along a different axis, program verifiers like
Dafny~\cite{leino10dafny}, VCC~\cite{CohenMST10},
Why3~\cite{filliatre13why3}, and 
Liquid Haskell~\cite{liquidhaskell} 
target both pure {\em and effectful} programs,
with side-effects ranging from divergence to
concurrency, but provide relatively weak logics for
specification (e.g., first-order
logic with a few selected theories like linear arithmetic).
%
They work primarily by computing verification conditions (VCs)
% \ch{The intro calls them VCs, but can't we call them WPs? (The term VCs
%   seems rather vague to me.) Or at least mention here that these VCs
%   are computed by a WP calculus?}\guido{But that's
%   \fstar-specific, while this is not meant to be}%
from programs, usually relying on annotations such as pre- and postconditions,
and encoding them to automated theorem provers (ATPs) such as
satisfiability modulo theories (SMT) solvers,
often providing excellent automation.


These two sub-fields have influenced one another, though
the situation is somewhat asymmetric.
%
On the one hand, most interactive provers have gained support for
exploiting SMT solvers or other ATPs, providing push-button automation
for certain kinds of assertions~\cite{CzajkaK17, smtcoqCAV2017,
PaulsonB10, KaliszykU14, KaliszykU15a}.
%
% The common design is that a goal in the proof assistant is encoded into
% first order logic and the SMT is queried. If the SMT finds a proof, the
% tools attempt to reconstruct them in the logic of the proof assistant
% and recheck them.
%
On the other hand, recognizing the importance of interactive proofs,
Why3~\cite{filliatre13why3} interfaces with ITPs like Coq.
%
However, working over proof obligations translated from Why3
requires users to be familiar not only with both these systems, but
also with the specifics of the translation.
% , which can be challenging.
% CH: This one seemed a weaker point than the one I kept above:
% proofs and their contexts are manipulated in terms of Coq
% abstractions, rather than source language concepts.
%
%
And beyond Why3 and the tools based on it~\cite{CuoqKKPSY12}, no other
SMT-based program verifiers have full-fledged support for interactive
proving, leading to several downsides:

\paragraph{\bf Limits to expressiveness} The expressiveness of program
verifiers can be limited by the ATP used.
When dealing with theories that are undecidable and
difficult to automate (e.g., non-linear arithmetic or
separation logic), proofs in ATP-based systems %like Dafny~\cite{leino10dafny}
may become impossible or, at best, extremely tedious.
\guido{citations would be nice}

\paragraph{\bf Boilerplate} To work around this lack of automation,
programmers have to construct detailed proofs by hand, often repeating
many tedious yet error-prone steps, so as to provide hints to the
underlying solver to discover the proof. In contrast, ITPs
with metaprogramming facilities excel at expressing
% CH: ad hoc -- negative connotations
%     automated -- ?
domain-specific automation to complete such tedious proofs.

\paragraph{\bf Implicit proof context} In most program verifiers,
the logical context of a proof
is implicit in the program text and depends on the control flow and
the pre- and postconditions of preceding computations. Unlike in
interactive proof assistants, programmers have no explicit access,
neither visual nor programmatic, to this context, making proof
structuring and exploration extremely difficult.
\\

% \da{Are we again pitching \metafstar more to the verifiers community?
% \ch{Yes, and I think it's great to choose one's battles!}
% Because that's how I read the intro at the moment, particularly 
% amplified by making this 1,2,3 list explicit and no similarly 
% explicit list for selling it to interactive theorem proving community.}
% \guido{I think so, yes}
% \guido{Added para below}\ch{I think the new paragraph is fighting a
%   losing battle.  It will dilute the main message and ruin our careful
%   positioning, which are the 3 points above, and you might also get
%   the ire of all the proof assistant people with these dubious claims.}

%\guido{effects?}%
%While interactive theorem provers such as Coq or Lean do not suffer from
%these drawbacks, they require more effort from the user in the
%development of proofs and tactics.\ch{We have very little scientific evidence
%  that this is indeed the case.}
%%
%Automated program verifiers, leveraging state-of-the-art SMT solvers,
%can reduce such effort considerably and, especially for non-experts,
%make verification more efficient.\ch{I'm calling this a
%  claim backed by nothing by hot air.}
% \da{This is indeed hot air, commenting out for time being.}

In direct response to these drawbacks, we seek a system that successfully 
combines the convenience of an
automated program verifier for the common case, while seamlessly
transitioning to an interactive proving experience for those parts of
a proof that are hard to automate. Towards this end, we
propose \metafstar, a tactics and metaprogramming framework for
the \fstar~\cite{mumon,dm4free} program verifier.

\subsection*{Highlights and Contributions of \metafstar}

\fstar has historically been more deeply rooted as an
SMT-based program verifier.
%
Until now, \fstar discharged VCs
exclusively by calling an SMT solver (usually Z3~\cite{MouraB08}), 
%
providing good automation for many common
program verification tasks, but
also exhibiting the drawbacks discussed above.


% However, some theories cannot be efficiently handled by SMT solvers.
% %
% For instance, non-linear arithmetic is well-known to be poorly handled
% by SMT solvers in general, becoming especially unreliable when combined
% with other theories.
% %
% Worse, these solvers provide few facilities for user interaction beyond
% being queried for facts it can prove.
% %
% The lack of fine-grained control over proofs and automation, especially
% when (inherently incomplete) SMT solving times out, is often mentioned
% as a shortcoming of SMT-only proofs.
% \ch{Citation list missing to support last claim; otherwise we can assume responsibility for the claim}
% \guido{Yes, cite would be nice}

\metafstar is a framework that allows \fstar users to manipulate VCs
using \emph{tactics}. More generally, it
supports \emph{metaprogramming}, allowing programmers to script
the construction of programs, by manipulating their
syntax and customizing the way they are type-checked. This allows
programmers to (1) implement custom procedures for manipulating VCs;
(2) eliminate boilerplate in proofs and programs; and (3) to inspect
the proof state visually and to manipulate it programmatically,
addressing the drawbacks discussed above.
%
SMT still plays a central role in \metafstar: a typical usage involves
implementing tactics to transform VCs, so as to bring them into theories
well-supported by SMT, without needing to (re)implement full
decision procedures.
%
Further, the generality of \metafstar allows implementing non-trivial
language extensions (e.g., typeclass resolution) entirely as
metaprogramming libraries, without changes to the \fstar{}
type-checker.

% \da{This makes it sound like no additional compiler-work was needed 
% for getting typeclasses to work, which is not true, right?}
% \guido{Meta-implicits required work, but having that, typeclasses
% are just metaprogrammed (modulo syntactic sugar). Typeclass
% \emph{resolution} is indeed completely metaprogrammed.}
% \ch{Tried to make this claim stronger above}
% \ch{Is Clement's pattern matcher another good example here maybe?}

%
%% \metafstar tactics can be applied to specific assertions, independently
%% of the overall structure of the VC computed by \fstar{}, whose shape,
%% given \fstar's implicit proof context, can be hard for users to
%% predict. \nik{repeated in the itemized list below}
%
% A first technical challenge in doing so is being able to extract
% only a \emph{part} of the VC that is supposed to be processed by a tactic,
% which we solve by a polarity-based criterion described ahead.
%
% Within \metafstar, a call to the SMT solver is just another (end-game)
% tactic.
% \ch{After the recent cuts this paragraph sounds less impressive
%   and dangerously non-novel. Having the SMT solver as just another end
%   game tactic is exactly what the hammers provide. I think we need to
%   explain upfront that our combination between tactics and SMT is
%   different than what the hammers already provide. Also the thus
%   below is not at all justified. One solution would be to simply bring
%   back the cutouts above.}
%   \guido{mentioned above that we improve on the 3 downsides
%   we listed. not great anyway}
%
%% Another key feature is that, at any time during the proof, the
%% facts remaining to be proven can be sent to the SMT solver.
%
%
% We use this approach successfully several times
% (\autoref{sec:experiments}).

% \guido{say something about robustness?}


The technical {\bf contributions} of our work include the following:

% \ch{Not clear how exhaustive the above description was trying to be.
%   Things like generating verified programs wasn't mentioned yet, for
%   instance, but maybe that's okay.}\ch{This only appears in the
%   evaluation, very quickly in a large enumeration but it should be
%   worth mentioning it earlier and more explicitly.}

% This paper proposes a way to allow for fruitful interaction
% between these worlds in
% the context of \fstar{}, a programming language and verification tool
% combining user-defined effects and a full-spectrum dependent type
% system~\cite{mumon, dm4free}.
% %
% Verification in \fstar involves a combination of dependent type-checking
% with the generation of verification conditions in a Hoare logic for
% higher-order programs, discharged by SMT solving.
% %
% However, it previously lacked the ability to script proofs,
% programs, and typing derivations, making it difficult to construct
% large proofs with reliable and efficient automation.
% %
% In this work, we describe and evaluate
% \metafstar, a metaprogramming framework that
% provides users with programmatic access to \fstar's logic and internal
% APIs of its implementation.
% %
% Tactics are a special class of metaprograms that manipulate proof goals,
% and SMT solving is a particular tactic that either fully proves a goal
% or fails.\guido{In honesty we haven't done that yet}

% The novelty of \metafstar lies in the close integration of
% metaprogramming, verification condition generation, and SMT solving,
% enabling new styles of program construction and proof automation.
%

\paragraph*{\bf ``Meta-'' is just an effect (\autoref{sec:tac-effect})}
%
% Rather than defining a separate language, a
% central design principle is that metaprograms are just
% a particular kind of program.
%
\metafstar is implemented using \fstar's extensible effect
system, which keeps programs and
metaprograms properly isolated.
% where required. -- it's not optional, and there is already proper
%
% and the effect system also ensures that ordinary programs are
% isolated from metaprograms, except through calls
% to primitives provided by \metafstar.
% \ch{not the other way around?}\guido{attempted to improve..}
%
Being first-class \fstar programs, metaprograms
are typed, call-by-value, direct-style, higher-order functional
programs, much like the original ML.
% \guido{I find the ``intended'' a bit weird}\da{envisioned?}\guido{Better}
%
Further, metaprograms can be themselves verified (to a degree, see
\autoref{sec:specifying-metaprograms}) and metaprogrammed.
% \ch{Puzzling. To what degree? If this can't be
%   explained here should make sure that this gets explained later, for
%   instance in \autoref{sec:tac-effect}, where this paragraph references.}
%   \guido{Added a forward pointer to the section, intended to mean
%   ``to a degree which will be explained there''}


\paragraph*{\bf Reconciling tactics with VC generation (\autoref{sec:vcgen})}
%
In program verifiers the programmer often guides
the solver towards the proof by supplying intermediate 
assertions. \metafstar retains this style,
but additionally allows assertions to be solved by 
tactics. To this end, a contribution of
our work is extracting, from a VC, a proof state
encompassing all relevant hypotheses, including those implicit
in the program text.
%
% This enables a new verification style in which
% tactic proofs and the SMT solver complement each other well.

\paragraph*{\bf Executing metaprograms efficiently (\autoref{sec:evaluation})}
Metaprograms are executed during type-checking.
%
As a baseline, they can be interpreted using \fstar's existing
(but slow) abstract machine for term normalization, or
a faster normalizer based on normalization by evaluation
(NbE)~\cite{bergerNbe, BoespflugDG11}.
%
% A contribution \guido{not in the summary of contributions.
% but do we even want to say this?}
% of this paper is implementing a new, more efficient normalizer for \fstar based
% on normalization by evaluation (NbE)~\cite{bergerNbe}.
%\guido{or \S\ref{sec:evaluation}}
%
For much faster execution speed, metaprograms can also be run natively.
%
This is achieved by combining the existing extraction mechanism of
\fstar to OCaml with a new framework for safely extending the \fstar
type-checker with such native code.
%
%% At the core of this framework is a calculus for
%% interoperability between \fstar's normalizers and extracted OCaml code,
%% supporting higher-order interactions.
%
%% We examine the tradeoffs of these execution strategies and show how to
%% use them to gain orders-of-magnitude performance improvements relative
%% to the baseline normalizer.
%% \ch{Unclear to what extent this last part
%%   will make it into the paper.}

\paragraph*{\bf Examples (\autoref{sec:by-example}) and evaluation (\autoref{sec:experiments})}
We evaluate \metafstar on several case studies.
%
First, we present a functional correctness proof for the Poly1305
message authentication code (MAC)~\cite{poly1305}, using a novel combination
of proofs by reflection for dealing with non-linear arithmetic and SMT
solving for linear arithmetic. We measure a clear gain in proof
robustness: SMT-only proofs succeed only rarely
(for reasonable timeouts),
whereas our tactic+SMT proof is concise, never fails,
and is faster.
% \ch{This last
%   part is also quite bogus. The fact that the rare succeeding cases
%   are fast seems a strange artifact, not something to compare against.
%   If I tell you that you are as fast as the success case of
%   a bus that fails to arrive at all 99\% of the time, do you take
%   that as a compliment?}
%
Next, we demonstrate an improvement in expressiveness,
by developing a small library for proofs of heap-manipulating programs in
separation logic, which was previously out-of-scope for \fstar.
%
Finally, we illustrate the ability to automatically construct verified
effectful programs, by introducing a library for metaprogramming verified
low-level parsers and serializers with applications to network
programming, where verification is accelerated by processing
the VC with tactics,
and by programmatically tweaking the SMT context.

\medskip
We conclude that tactics and metaprogramming
can be prosperously combined with VC generation and SMT solving
to build verified programs
with better, more scalable, and more robust automation.

\medskip
The full version of this paper, including appendices, can be found online in
{\url{https://www.fstar-lang.org/papers/metafstar}}.

% \da{TODO: double check the above list once s5 is settled}

% \subsection{Summary of contributions and outline}

% This paper describes the design and formalization of a new
% metaprogramming framework suitable for aiding the SMT-based
% verification of effectful programs, together with the implementation
% and evaluation of this framework in \fstar.
% %
% The main contributions of this paper are:

% \begin{enumerate}
% \item The design of the \metafstar effect-based metaprogramming framework
%       and its implementation in the \fstar toolchain, providing a new
%       verification tool that encourages the judicious combination of
%       tactic- and SMT-based automation. \ifconference(\S\ref{sec:design})\fi

% % \item A trust argument for the \metafstar design in the style of LCF, accounting
% %       for metaprogramming in both proof-relevant and proof-irrelevant
% %         contexts. \ifconference(\S\ref{sec:trust})\fi\guido{unsure this is a contribution,
% %         I think it's pretty vanilla, no?}

% \item A way to integrate metaprogramming with verification
%       condition generation in a program verifier for effectful
%       programs. \ifconference(\S\ref{sec:vcgen})\fi

% \item A new multi-language interoperability framework to support
%       integrating compiled plugins with symbolic normalizers at higher
%       order. Our framework is efficient by avoiding redundant
%       embeddings of compiled terms at interface boundaries; and makes
%       novel use of parametricity to allow non-trivial use of
%       computation over open terms in compiled code. \ifconference(\S\ref{sec:evaluation})

% \item An evaluation of \metafstar on several case studies,
%       with qualitative and quantitative assessments of various
%       elements of our work. \ifconference(\S\ref{sec:by-example} and \S\ref{sec:experiments})\fi
% \end{enumerate}

%% \paragraph{Separation logic + SMT for stateful programs}
%% Despite being expressible in its logic, proofs using separation
%% logic~\cite{Reynolds:2002:SLL:645683.664578} have been impractical
%% in \fstar. A key difficulty is in solving for existentially bound heap
%% frames at nearly every step of the proof, which is intractable with an
%% SMT solver. With \metafstar, we programmed a tactic \ls$sl_auto$ to
%% verify stateful programs in a fragment of separation logic. The
%% tactic only solves for heap frames, leaving the rest of the proof to
%% the SMT solver, allowing for the mixture of separation logic
%% with \fstar's regular specifications and theories well-supported in
%% SMT. Here is a small program automatically verified using this approach:

%% \begin{lstlisting}
%% let swap (r1 r2:ref int) : STATE unit (fun post m0 -> exists x y. m0 == r1 |-> x $\bullet$ r2 |-> y /\ post () (r1 |-> y $\bullet$ r2 |-> x))
%%     by sl_auto = let x = !r1 in r1 := !r2; r2 := x
%% \end{lstlisting}

%% \paragraph{Non-linear arithmetic in proofs of cryptography}
%% Non-linear arithmetic poses a serious challenge for predictable and
%% scalable SMT proofs, yet it is crucially needed for the verification
%% of cryptographic primitives~\cite{vale,haclstar}.
%% % , an important use case for \fstar~\cite{everest}
%% % -- CH: let's not make intro overly F*-specific
%% We program a tactic for canonicalizing
%% terms in commutative semirings, and apply it to prove the main lemmas
%% in the correctness proof of an optimized implementation of the
%% Poly1305 MAC algorithm~\cite{poly1305}. Where previous proofs in tools
%% like \fstar and Dafny were painfully detailed and brittle (e.g.,
%% requiring dozens of explicit steps of rewriting of large polynomials
%% to nudge the SMT solver towards a proof), in \metafstar the proof is
%% reduced to a single call to our \ls$canon_semiring$ tactic, leaving
%% the rest of the proof to the efficient, linear arithmetic theory in
%% Z3~\cite{MouraB08}.

%% \ifconference
%% \else
%% Below is a main lemma in the proof of correctness of OpenSSL's
%% Poly1305 Message Authentication Code (MAC)~\cite{poly1305}. A key step
%% in the proof is the final assertion, a identity in non-linear
%% arithmetic that is reduced by the tactic \ls$canon_semiring int_csr$
%% to a problem in linear arithmetic that can easily be discharged by an
%% SMT solver.

%% \begin{lstlisting}
%% let poly_multiply (n p r h r0 r1 h0 h1 h2 s1 d0 d1 d2 h1 h2 hh : int) =
%%   let r14 = r1 / 4 in
%%   let h_r_expand = (h2 * (n * n) + h1 * n + h0) * ((r14 * 4) * n + r0) in
%%   let hh_expand = (h2 * r0) * (n * n) + (h0 * (r14 * 4) + h1 * r0 + h2 * (5 * r14)) * n
%%                 + (h0 * r0 + h1 * (5 * r14)) in
%%   let b = (h2 * n + h1) * r14 in
%%   modulo_addition_lemma hh_expand p b; //(hh_expand + b * p) % p = hh_expand % p
%%   assert (h_r_expand == hh_expand + b * (n * n * 4 + ($-$5))) by (fun _ -> canon_semiring int_csr)
%% \end{lstlisting}
%% \fi
%% \paragraph{Metaprogramming the construction of verified effectful programs}
%% Tactics in \metafstar are a special case of general metaprogramming
%% that supports the inspection and construction of arbitrary program
%% syntax. Our third case study uses \metafstar to implement a
%% \emph{correct-by-construction} transformation of \fstar terms within
%% \fstar itself. Specifically, we generate a low-level formatter for
%% network messages from a high-level, pure specification of the network
%% format, with a proof of their equivalence. The generated low-level
%% code is compatible with the rest of the \fstar toolchain, allowing it
%% to be extracted to efficient C code by % the KreMLin compiler,
%% \fstar's C backend~\cite{lowstar}.


%% \ifconference
%% \else
%% This paper is a short version of a longer report available online at \url{https://arxiv.org/abs/1803.06547}.
%% \metafstar is part of the latest release of \fstar and is under continued active development.
%% The release, including several examples metaprograms, is available
%% at \url{https://github.com/FStarLang/FStar/releases/tag/v0.9.6.0}.
%% \fi
%% \ifconference

% \section{Meet \metafstar; Or, my proofs are so much nicer after I Met-a-\fstar}
\section{\metafstar by Example}
\label{sec:by-example}

\fstar is a general-purpose
programming language aimed at program verification.
%
It puts together the automation of an SMT-backed deductive
verification tool with the expressive power of a language
with full-spectrum dependent types.
%
Briefly, it is a functional, higher-order, effectful, dependently typed
language, with syntax loosely based on OCaml.
%
\fstar supports refinement types and Hoare-style specifications,
computing VCs of computations via a type-level weakest
precondition (WP) calculus packed within \emph{Dijkstra
monads}~\cite{fstar-pldi13}.
%
\fstar's effect system is also user-extensible~\cite{dm4free}.
%
Using it, one can model or embed % higher-order -- C and ASM?
imperative programming in
styles ranging from ML to
C~\cite{lowstar} and
assembly~\cite{valefstar}.
%
After verification, \fstar programs can be extracted to efficient
OCaml or F\# code. A first-order fragment of \fstar, 
called \lowstar, can also be extracted to C via the KreMLin compiler~\cite{lowstar}.
% with manual
%memory management.

This paper introduces \metafstar, a metaprogramming framework for
\fstar that allows users to safely customize and extend \fstar in many
ways.
% \aseem{automatically synthesize verified F* code?}%
%
For instance, \metafstar can be used to preprocess or solve proof
obligations; synthesize \fstar expressions; generate
top-level definitions; and resolve implicit arguments in user-defined
ways, enabling non-trivial extensions.
%
% \nik{revised here ... please check}%
% \aseem{confused, there are only 4 features in the sentence above,
% and there too generate top-level definitions seems very close to
% synthesize F* expressions}%
This paper primarily discusses the first two features.
%
Technically, none of these features deeply increase the expressive power
of \fstar, since one could manually program in \fstar terms that can
now be metaprogrammed. However, as we will see shortly, manually
programming terms and their proofs can be so prohibitively costly as
to be practically infeasible.
% \guido{Is ``none of these'' meant to mean the last two, or all?
% If all of them, I don't agree. E.g. I think the parsers and serializers
% illustrate that, since the metaprogrammed code wouldn't typecheck
% without the special handling of the bijectivity VC. One can probably
% work around that adding assertions, yes, but the program changes.}
% \aseem{One can always write these parsers
% and synthesizers manually, and manually prove the bijection.}
% \guido{Doing it manually implies the program is slightly
% different since it includes the assertions, no? Although indeed
% they extract to the same code. Still I'm not so convinced,
% having tactics we can basically script how \fstar computes and preprocesses VCs
% before encoding. SL is an example.}
%\guido{Bit the bullet}

% and, finally,
% unleash the full power of \fstar's typing judgment (which is undecidable) into the users' hands.

\metafstar is similar to other tactic frameworks, such as
Coq's~\cite{Delahaye00} or Lean's~\cite{EbnerURAM17}, in presenting a
set of goals to the programmer, providing commands to break them down,
allowing to inspect and build abstract syntax, etc.
%
In this paper,
% , mostly for space reasons -- also for "what's interesting" reasons
we mostly detail the characteristics where \metafstar \emph{differs}
from other engines.

% \guido{I would love to mention native plugins above, but there's no \metafstar}
% \guido{After writing this para, I'm wondering why we're only trying to
% sell the ``SMTs are incomplete'' story in the abstract. Opinions?}

This section presents \metafstar informally, displaying its usage
through case studies. We present any necessary \fstar background as needed.

\iffalse
At the core of \fstar is a language of pure total functions based on
an extensional dependent type theory. As such, type-checking in \fstar
is undecidable, but its implementation provides various practical
heuristics to type programs.
%
This core language of pure functions also forms the logic for \fstar's
program specifications and their proofs.

Beyond the core, \fstar provides primitive support for divergent
computations. \fstar also provides an extensible system of monadic
effects, which allows users to program other effects on top of purity
and divergence simply by providing monadic definitions.
%
The effect system includes a weakest-precondition (WP) calculus for
each effect, based on a higher-order adaption of Dijkstra's weakest
precondition transformers~\cite{fstar-pldi13,dm4free}.
%
Programmers can specify precise pre- and postconditions for their
programs using a notation similar to Hoare Type Theory~\cite{nmb08htt}
and \fstar checks that the specifications inferred for these programs
are subsumed by the programmer-given ones.
%
This subsumption check is reduced to a logical validity problem that
was until now discharged using the Z3 SMT solver, and, as of this
work, also metaprogrammed tactics.

%
% \guido{How about: that \fstar discharges through SMT solving.
% In the case the SMT cannot prove the provided formula, the only mean
% the user had to transform it was selectively calling lemmas
% to introduce new facts (and subobligations) into the logical context
% of the program. This work adds a much more expressible mean,
% metaprogrammed tactics.}\cpc{I think this is still too much.  Introducing new lemmas isn't a way to discharge the satisfiability problem, its a way to tweak it.  I agree that it's connected to tactics at a very high level, but I think it makes things less clear to even mention it here.}


\paragraph{Syntax} \fstar syntax is
roughly modeled on OCaml (\ls$val$, \ls$let$, \ls$match$ etc.)
although there are differences to account for additional features.
%
Binding occurrences \ls$b$ of variables take the form \ls$x:t$, declaring
a variable \ls$x$ at type \ls$t$; or \ls$#x:t$ indicating a
binding for an implicit argument.
%
\ls@fun (b$_1$) ... (b$_n$) -> t@ introduces a lambda abstraction, whereas
%
\ls@b$_1$ -> ... -> b$_n$ -> c@ is the shape of a curried function type---we
emphasize the lack of enclosing parentheses on the \ls@b$_i$@\cpc{Is
this needed?}\ch{no!}\nik{yes!}

Function codomains are computation types of the
shape \ls$M t wp$ where \ls$M$ is an effect
label, \ls$t$ a return type, and \ls$wp$ a specification term.
%
Refinement types are written \ls$b{t}$,
e.g., \ls$x:int{x>=0}$ is the type of non-negative integers
(\ls$nat$).
%
As usual, a bound variable is in scope to the right of its binding; we
omit the type in a binding when it can be inferred; and for
non-dependent function types, we omit the variable name.
%
For example, the type of the
pure append function on vectors is written
%
\ls$#a:Type -> #m:nat -> #n:nat -> vec a m -> vec a n -> vec a (m + n)$,
%
with the two unnamed explicit arguments and the return type depending
on the first three implicit arguments.
%
When no effect is specified for an \fstar{} function, the default
primitive effect \ls$Tot$ (for pure, total computations) is used
implicitly.
\fi

% \subsection{Instrumenting and Discharging VCs with Tactics}
% \label{sec:vcgen_example}

% For a first taste of \fstar, let us consider the following implementation of quicksort:
% %
% \begin{lstlisting}
% let rec quicksort (#a:Type) (f:total_order a) (l:list a) : Tot (list a) (decreases (length l)) =
%     match l with
%     | [] -> []
%     | pivot::tl -> let hi, lo = partition (f pivot) tl in quicksort f lo @ pivot :: quicksort f hi
% \end{lstlisting}
% %
% \fstar's syntax resembles OCaml's, but with some differences to
% accommodate its richer type system. The recursive
% function \ls$quicksort$ takes a type argument \ls$a$ (marked as an
% implicit argument by the `\ls$#$' sign); \ls$f$, a total order on \ls$a$;
% and \ls$l$ of type \ls$list a$. The annotation on its result is a \emph{computation type}
% %
% \ls$Tot (list a)$, which states that \ls$quicksort f l$ is pure,
% total function (i.e., terminating on all inputs with no side effects)
% returning a \ls$list a$. To prove termination, \fstar employs a semantic
% check that requires the arguments to each recursive call
% to decrease according to some (user-specified) well-founded
% order. The \ls$decreases (length l)$ annotation here is a programmer-supplied
% hint (a termination measure) that causes \fstar to
% check that, at each recursive call, the length of \ls$quicksort$'s list
% argument decreases.

% Beyond basic type safety and totality checking,
% \fstar also supports stating and proving nearly arbitrary
% program properties through its use of
% dependent types. For example, to specify that \ls$quicksort$ is
% correct, one could annotate it with the following type signature:
% %
% \begin{lstlisting}
% val quicksort (#a:Type) (f:total_order a) (l:list a) : Tot (m:list a{sorted f m /\ perm l m})
% \end{lstlisting}
% %
% The result type shows a use of \fstar's refinement
% types, \ls@x:t{$\phi$}@, the type of terms \ls$x:t$ that satisfy the
% formula $\phi$. The particular refinement used here is the type of
% lists \ls$m:list a$ that are sorted according to the total
% order \ls$f$, and that are permutations of the input list \ls$l$.
% In fact, the type \ls$total_order a$ is defined as the refinement
% of \ls$a -> a -> Tot bool$ satisfying the expected order-theoretic properties.
% \guido{A bit misplaced but can't think of a better presentation.}

% To check that the implementation of \ls$quicksort$ actually has this
% type, \fstar typechecks the implementation, inferring a specification
% for it (via a weakest precondition calculus for higher-order
% programs built into \fstar), and checks that the
% inferred specification is subsumed by the user's annotation. This
% subsumption check is a logical problem that is typically encoded to an
% SMT solver.

% In this particular case, \fstar in fact fails to show that \ls$quicksort$, as
% written above, satisfies the user-provided specification. The proof requires
% a bit of help from the user in the form of some additional facts
% about list concatenation; e.g. that \ls$pivot :: quicksort f hi$ is
% properly sorted.

% The traditional way to complete this proof in \fstar, much like in many other
% Hoare-style verification tools (e.g., Dafny, Why3, etc.), is to decorate
% the program with lemmas and assertions, i.e., intermediate steps that
% guide the SMT solver to the proof, e.g., as in this revised implementation:

% \begin{lstlisting}
% let rec quicksort #a f l = match l with | [] -> []
%     | pivot::tl -> let hi, lo = partition (f pivot) tl in
%                let m = quicksort f lo @ pivot :: quicksort f hi in
%                sorted_app (quicksort f lo) pivot (quicksort f hi); (* lemma call *)
%                perm_app pivot tl (quicksort f lo) (quicksort f hi); (* lemma call *)
%                m
% \end{lstlisting}

% A lemma in \fstar is a pure function, which when applied in a context
% to arguments that validate the formula in the \ls$requires$ clause
% (aka the \emph{precondition}), proves that the formula in
% the \ls$ensures$ clause (aka the \emph{postcondition}) is also
% valid. Lemmas have no computational content (they return the unit value `\ls$()$') and are
% erased before the program executes. The two lemmas from our
% example, \ls$perm_app$ and \ls$sorted_app$, are proven separately, e.g.,
% the signature of \ls$perm_app$ is as follows:
% %(\ls$sorted_app$ is similar,
% %and ensures that \ls$sorted f m$ holds).

% \begin{lstlisting}
% val perm_app #a (hd:a) (tl l_1 l_2:list a) : Lemma (requires perm tl (l_1 @ l_2)) (ensures perm (hd::tl) (l_1 @ hd::l_2))
% \end{lstlisting}

% This verified implementation uses an ``intrinsic proof''. The proof is
% structured by the control flow of the program itself, asserting
% properties at specific program points, avoiding duplication and making
% the overall proof more compact.
% %
% Additionally, the lemmas \ls$perm_app$ and \ls$sorted_app$ capture
% some essence of the inductive invariant of quicksort and the user
% intervention in the proof is limited to only those properties that the
% SMT solver is incapable of proving automatically.  In this case
% explicit lemma invocations like
% %
% \ls$perm_app pivot (quicksort f lo) (quicksort f hi)$ get the job done,
% but in general, finding the precise instantiation to use requires some
% careful thought from the user.\footnote{\fstar also provides a
% facility for expert users to annotate lemmas with patterns that an SMT
% solver can use to automatically find lemma instantiations via
% E-matching. When used judiciously this can be very effective, but poor
% choices of patterns can also drive the SMT solver off the rails.} This
% can be especially hard when the user has no way to directly inspect
% the current set of facts and obligations, and, in some cases,
% the great amount of ``hints'' needed from
% the user may even entirely defeat the purpose of using an automated solver~\cite{why3reflection}.
% \guido{citing a Why3 paper about verifying GMP which
% clearly states that NL-arith is a major pain in the ass. Maybe we want to cite it ahead instead.}

% \paragraph{\metafstar: Associating tactics with assertions}
% Our third attempt at implementing quicksort shows how, using \metafstar, specific
% assertions in a program's VC can be spawned off to a user-specified
% tactic rather than the VC being sent wholesale to an SMT solver:

% \begin{lstlisting}
% let rec quicksort #a f l = match l with | [] -> []
%     | pivot::tl -> let hi, lo = partition (f pivot) tl in
%                let m = quicksort f lo @ pivot :: quicksort f hi in
%                sorted_app (quicksort f lo) pivot (quicksort f hi);  (* retained for illustration *)
%                assert (perm l m) by (qs_tac l); (* new *)
%                m
% \end{lstlisting}
% \label{sec:quicksort}

% Here, the new \ls$assert p by t$ form asserts a formula \ls$p$ to
% to be solved by the \metafstar
% tactic \ls$t$. The context for the goal not only includes hypotheses about
% variables in scope, but also assumptions about the program's control
% flow, such as equations about pattern matching (e.g. \ls$l == pivot::tl$)
% and local definitions (e.g. \ls$m = quicksort f lo @ pivot :: quicksort f hi$),
% as
% well as postconditions of functions evaluated prior to that program
% point (e.g. \ls$sorted f m$, the postcondition of \ls$sorted_app$).

% The tactic \ls$qs_tac$ is a small script meant to tackle the goal
% %
% \ls$perm l (quicksort f lo @ pivot :: quicksort f hi)$:
% %
% \begin{lstlisting}
% let qs_tac l () = with_pat (fun y (h:hyp (squash (l == y))) -> rewrite h); mapply (`perm_app)
% \end{lstlisting}
% %
% Using a library for pattern matching the proofstate (also programmed
% in \metafstar, see \autoref{sec:other}), \ls$qs_tac$
% rewrites the goal using the equation for \ls$l$,
% % (i.e., \ls$l=pivot::tl$, from an equation due to pattern-matching),
% and then solves the goal
% %
% \ls$perm (pivot::tl) m$\guido{was:
% \ls$perm (pivot::tl) (quicksort f lo @ pivot :: quicksort f hi)$, why replace \ls$m$?}
% by unifying it with the postcondition of \ls$perm_app$ and leaving the
% rest to SMT (i.e., the preconditions of \ls$perm_app$).
% \guido{this looks kinda adhoc, why don't we change it to a simplifier
% that also attempts the previous two lemmas? I guess to make it robust
% we'd need a congruence closure.}

% Trapping proof obligations in a program is not only limited to
% assertions. Our final version of quicksort shows that users can
% associate tactics with proof obligations arising from subtyping too:
% %
% \begin{lstlisting}
% let rec quicksort #a f l = match l with | [] -> []
%     | pivot::tl -> let hi, lo = partition (f pivot) tl in
%                let m = quicksort f lo @ pivot :: quicksort f hi in
%                m <: _ by (qs_tac' l)
% \end{lstlisting}
% %
% On the last line, \fstar's type inference concludes that \ls$m$ must
% be subsumed to the annotated result type of the function, and the
% tactic \ls$qs_tac'$ (similar to \ls$qs_tac$) gets to discharge the
% proof of the required refinement formula, \ls$sorted f m /\ perm l m$,
% in the current proof context.

% These small first examples illustrate how programmers can control
% fragments of their proofs using tactics, partially solving goals and
% leaving the rest to SMT.
% %
% At the small scale of quicksort, this hardly makes a difference; but, as
% the next section argues, larger developments may benefit greatly from
% synergy between SMT solving and user-defined tactics.

\subsection{Tactics for Individual Assertions and Partial Canonicalization}
\label{sec:non-linear}

% \guido{All of this was said before.
%     \st{
% Proofs of program correctness involve many concerns, including
% reasoning about a program's memory, its termination behavior, error
% handling, and even security properties like side-channel resistance,
% just to name a few.
% %
% In a setting with tactic-based proof only, the programmer has complete
% control, but is faced with building proofs scripts and ad hoc
% automation for all these concerns.
% %
% On the other hand, SMT solvers can be configured to automatically
% discharge many large proofs, relying on various logic features in
% combination with theories including linear arithmetic, quantifiers based
% on E-matching, uninterpreted functions, datatypes, and boolean
% satisfiability.

% SMT solving, however, is no panacea. Proofs that involve reasoning
% using theories that are not well supported can be inefficient,
% unpredictable, or just plain impossible.
% %
% Proofs about non-linear arithmetic are a case in point.
% %
% }}

Non-linear arithmetic reasoning is crucially needed for the
verification of optimized, low-level cryptographic
primitives~\cite{vale,haclstar}, an important use case
for \fstar~\cite{everest} and other verification frameworks, including
those that rely on SMT solving alone (e.g., Dafny~\cite{leino10dafny})
as well as those that rely exclusively on tactic-based proofs (e.g.,
FiatCrypto~\cite{fiat-crypto}).
%
While both styles have demonstrated significant successes, we make
a case for a middle ground, leveraging the SMT solver for
the parts of a VC where it is effective,
and using tactics only where it is not.

% using their best features
% while steering clear of their weaknesses.



%% \subsection{Non-linear arithmetic in verified cryptography}

%% Proving the correctness of optimized low-level arithmetic code
%% involves many concerns, including reasoning about a program's memory,
%% its termination behavior, error handling, and even side-channel
%% resistance. Amid all these concerns are proofs about non-linear
%% arithmetic.
%% %
%% In a setting with tactic-based proof only, the programmer has complete
%% control, but is faced with building proofs scripts and ad hoc
%% automation for all these concerns.
%% %
%% On the other hand, SMT solvers can be configured to automatically
%% discharge many large proofs, relying on various logic features in
%% combination with theories include linear arithmetic, quantifiers based
%% on E-matching, uninterpreted functions, datatypes, and boolean
%% satisfiability.
%% %
%% But, reasoning about non-linear arithmetic although supported by many
%% solvers, scales poorly and becomes unpredictable, especially in
%% combination with other theories.
%% %
%% In response to this unpredictability of non-linear reasoning many
%% large SMT-based verification efforts simply disable support for
%% non-linear arithmetic~\cite{Ironclad, IronFleet, vale, haclstar, ivy},
%% relying instead on explicitly calling lemmas from libraries of
%% separately proven facts for even the most basic properties, including
%% associativity, commutativity and distributivity---with great
%% determination, this can be made to work, but proofs are painfully
%% verbose and brittle.

%% Our first example of \metafstar illustrates how a tactic can be used
%% to tackle just the one part of a proof that involves non-linear
%% reasoning, automatically reducing several dozen intricate, manual
%% steps of manual reasoning about non-linear arithmetic to problems in
%% linear arithmetic, which can then be efficiently decided by an SMT
%% solver.

% \subsection*{Three Styles of Non-linear Arithmetic Proofs for Poly1305}
% \label{sec:semirings}

We focus on Poly1305~\cite{poly1305}, a widely-used cryptographic MAC
that computes a series of integer multiplications and additions modulo
a large prime number
%
$p = 2^{130} - 5$.  Implementations of the Poly1305 multiplication and
mod operations are carefully hand-optimized to represent 130-bit
numbers in terms of smaller 32-bit or 64-bit registers, using clever
tricks;
%to express mod $p$ operations with addition, multiplication,
%and bitwise operations rather than division.
proving their correctness
requires reasoning about long sequences of
additions and multiplications.

\paragraph{\bf Previously: Guiding SMT solvers by manually
  applying lemmas}

Prior proofs of correctness of Poly1305 and other cryptographic
primitives using SMT-based program verifiers, including
\fstar~\cite{haclstar} and Dafny~\cite{vale},
% Ironclad, IronFleet
% \guido{confused: the last two citations do not seem to mention Poly1305?}, 
use a combination of SMT automation and manual application of lemmas.
%
On the plus side, SMT solvers are excellent at linear arithmetic, so
these proofs delegate all associativity-commutativity (AC) reasoning
about addition to SMT.
%
Non-linear arithmetic in SMT solvers, even just AC-rewriting and
distributivity, are, however, inefficient and unreliable---so much so
that the prior efforts above (and other works too~\cite{Ironclad,
  IronFleet})
% \ch{double check please}\guido{confirmed for Ironclad.  But the
% IronFleet paper and others above doesn't seem to mention that?}
% CH: I think we're good now
simply turn off support for non-linear arithmetic
in the solver, in order not to degrade verification
performance across the board due to poor interaction of theories.
% \guido{we should cite something}.
%
Instead, users need to explicitly invoke lemmas.\footnote{\ls$Lemma (requires
pre) (ensures post)$ is \fstar notation for the type of a %n (erasable)
computation proving \ls$pre ==> post$---we omit \ls$pre$ when it is
trivial. In \fstar's standard library, math lemmas are proved using
SMT with little or no interactions between problematic theory
combinations. These lemmas can then be explicitly invoked in larger
contexts, and are deleted during extraction.}

For instance, here is a statement and proof of a lemma about Poly1305
in \fstar. The property and its proof do not really
matter; the lines marked ``\ls$(* argh! *)$'' do. In this particular proof,
working around the solver's inability to effectively reason about
non-linear arithmetic, the programmer has spelled out basic facts about
distributivity of multiplication and addition, by calling the library
lemma \ls$distributivity_add_right$, in order to guide
the solver towards the proof.
(Below, \ls$p44$ and \ls$p88$ represent $2^{44}$ and $2^{88}$ respectively)
%
\begin{lstlisting}
let lemma_carry_limb_unrolled (a0 a1 a2 : nat) : Lemma (ensures (
   a0 % p44 + p44 * ((a1 + a0 / p44) % p44) + p88 * (a2 + ((a1 + a0 / p44) / p44))
   == a0 + p44 * a1 + p88 * a2)) =
 let z = a0 % p44 + p44 * ((a1 + a0 / p44) % p44)
           + p88 * (a2 + ((a1 + a0 / p44) / p44)) in
 distributivity_add_right p88 a2 ((a1 + a0 / p44) / p44); (* argh! *)
 pow2_plus 44 44;
 lemma_div_mod (a1 + a0 / p44) p44;
 distributivity_add_right p44 ((a1 + a0 / p44) % p44)
           (p44 * ((a1 + a0 / p44) / p44)); (* argh! *)
 assert (p44 * ((a1 + a0 / p44) % p44) + p88 * ((a1 + a0 / p44) / p44)
           == p44 * (a1 + a0 / p44) );
 distributivity_add_right p44 a1 (a0 / p44); (* argh! *)
 lemma_div_mod a0 p44
\end{lstlisting}
%
Even at this relatively small scale, needing to explicitly instantiate
the distributivity lemma is verbose and error prone.
%
Even worse, the user is blind while doing so: the
program text does not display the current set of available facts nor
the final goal.
%
Proofs at this level of abstraction are painfully detailed in some
aspects, yet also heavily reliant on the SMT solver to fill in the aspects of
the proof that are missing.
%
% If the user, say, calls a lemma which does not help the proof, no error
% is raised; in this case, the lemma being proven will likely fail to
% verify, usually without providing much insight on why.
%
% \aseem{F* will catch the mistakes, can we rephrase this?}%
% \guido{It might not: if you call a lemma with the wrong set of arguments
% you get a useless fact, but no error. But wording is poor, yes.}%
% \guido{Attempted to improve.}%

Given enough time, the solver can sometimes find a proof without 
the additional hints, but this is usually rare and dependent on context,
and almost never robust.
%
In this particular example we find by varying Z3's random
seed that, in an isolated setting,
the lemma is proven automatically about 32\% of the time.
%
% Test was performed over F* commit dbc46e902bccb517c9018561a3ecb642012a7242
% and Z3 4.5.1 - 64 bit - build hashcode 1f29cebd4df6
% variying seeds from 0 to 200
%
The numbers are much worse for more complex proofs, and where the
context contains many facts, making this style quickly spiral out of
control.
%
For example, a proof of one of the main lemmas in
Poly1305, \ls$poly_multiply$, requires 41 steps of rewriting for
associativity-commutativity of multiplication, and
distributivity of addition and multiplication---making the proof
much too long to show here.
%
% To reduce this extreme level of manual intervention, prior projects
% have tried to craft E-matching triggers to have the SMT solver find
% the needed instantiations. But good E-matching triggers with AC and
% distributivity rewriting are hard to find; most send the SMT solver
% into a matching loop~\cite{AminLR14, comprehensions-smt}.

% \guido{citation would be nice}
% \ch{+1, this has
%   little credibility without a citation, and it's not clear at all
%   what this is talking about}
% \guido{added citations about matching loops, but they are
%   not related to NL arith}
% \guido{should also cite on ``prior projects'', or simply
% skip this part on triggers, IMO}

\paragraph{\bf SMT and tactics in \metafstar}
% Lightweight -- this means nothing to me

\label{sec:canon_semiring}
\label{sec:poly_multiply}

The listing below shows the statement and proof of \ls$poly_multiply$
in \metafstar, of which the lemma above was previously only a small part.
Again, the specific property proven is not particularly
relevant to our discussion. But, this time, the proof contains just
two steps.
%
\begin{lstlisting}
let poly_multiply (n p r h r0 r1 h0 h1 h2 s1 d0 d1 d2 h1 h2 hh : int) : Lemma
  (requires p > 0 /\ r1 >= 0 /\ n > 0 /\ 4 * (n * n) == p + 5 /\ r == r1 * n + r0 /\
            h == h2 * (n * n) + h1 * n + h0 /\ s1 == r1 + (r1 / 4) /\ r1 % 4 == 0 /\
            d0 == h0 * r0 + h1 * s1 /\ d1 == h0 * r1 + h1 * r0 + h2 * s1 /\
            d2 == h2 * r0 /\ hh == d2 * (n * n) + d1 * n + d0)
  (ensures (h * r) % p == hh % p) =
  let r14 = r1 / 4 in
  let h_r_expand = (h2 * (n * n) + h1 * n + h0) * ((r14 * 4) * n + r0) in
  let hh_expand = (h2 * r0) * (n * n) + (h0 * (r14 * 4) + h1 * r0
                            + h2 * (5 * r14)) * n + (h0 * r0 + h1 * (5 * r14)) in
  let b = (h2 * n + h1) * r14 in
  modulo_addition_lemma hh_expand p b;
  assert (h_r_expand == hh_expand + b * (n * n * 4 + ($-$5)))
      by (canon_semiring int_csr) (* Proof of this step by Meta--F* tactic *)
\end{lstlisting}

First, we call a single lemma about modular addition from \fstar's
standard library.
%
Then, we assert an equality annotated with a tactic (\ls$assert..by$).
%
Instead of encoding the assertion as-is to the SMT solver,
it is preprocessed by the \ls$canon_semiring$ tactic.
%
The tactic is presented with the asserted equality as its
goal, in an environment containing not only all variables in scope
but also hypotheses for the precondition of \ls$poly_multiply$
and the postcondition of the \ls$modulo_addition_lemma$ call
(otherwise, the assertion could not be proven).
%
The tactic will then canonicalize the sides of the equality, but notably
only ``up to'' linear arithmetic conversions.
%
Rather than fully canonicalizing the terms, the tactic just rewrites them
into a sum-of-products canonical form, leaving all the remaining
work to the SMT solver,
which can then easily and robustly discharge the
goal using linear arithmetic only.
% \ch{This is the most important part of the explanation
%   and should be made upfront, before / instead of diving head down into
%   boring details that will lose most readers}%


%
This tactic works over terms in the commutative semiring of integers
(\ls$int_csr$) using proof-by-reflection~\cite{GregoireM05, gonthier2008formal,ChaiebN08,Besson06}.
%
Internally, it is composed of a simpler, also proof-by-reflection based 
tactic \ls$canon_monoid$ that
works over monoids, which is then ``stacked'' on itself to build
\ls$canon_semiring$.
%
The basic idea of proof-by-reflection is to reduce most of the proof
burden to mechanical computation, obtaining much more efficient proofs
compared to repeatedly applying lemmas.
%
For \ls$canon_monoid$, we begin with
a type for monoids,
a small AST representing monoid values, and
a denotation for expressions back into the monoid type.
%
\begin{lstlisting}
type monoid (a:Type) = { unit : a; mult : (a -> a -> a); (* + monoid laws ... *) }
type exp (a:Type) = | Unit : exp a | Var : a -> exp a | Mult : exp a -> exp a -> exp a
(* Note on syntax: $\mathsf{\#a}$ below denotes that $\mathsf{a}$ is an implicit argument *)
let rec denote (#a:Type) (m:monoid a) (e:exp a) : a =
  match e with
  | Unit -> m.unit | Var x -> x | Mult x y -> m.mult (denote m x) (denote m y)
\end{lstlisting}
%
To canonicalize an \ls$exp$, it is first converted to a
list of operands (\ls$flatten$) and then reflected back to the monoid
(\ls$mldenote$).
%
The process is proven correct, in the particular case
of equalities, by the \ls$monoid_reflect$ lemma.
%
% We elide the proof and straightforward definitions.
%
% let rec flatten (#a:Type) (e:exp a) : list a =
%   match e with | Unit -> [] | Var x -> [x] | Mult x y -> flatten x @ flatten y
%
% let rec mldenote (#a:Type) (m:monoid a) (xs:list a) : a =
%   match xs with | [] -> m.unit | [x] -> x | x::xs' -> m.mult x (mldenote m xs')
%
\begin{lstlisting}
val flatten : #a:Type -> exp a -> list a
val mldenote : #a:Type -> monoid a -> list a -> a
let monoid_reflect (#a:Type) (m:monoid a) (e_1 e_2 : exp a)
              : Lemma (requires (mldenote m (flatten e_1) == mldenote m (flatten e_2)))
                       (ensures (denote m e_1 == denote m e_2)) = ...
\end{lstlisting}
%
At this stage, if the goal is \ls$t_1 == t_2$, we require two
monoidal expressions \ls$e_1$ and \ls$e_2$ such that
\ls$t_1 == denote m e_1$ and \ls$t_2 == denote m e_2$.
%
They are constructed by the tactic \ls$canon_monoid$ by inspecting
the \emph{syntax} of the goal, using \metafstar's reflection
capabilities (detailed ahead in~\autoref{sec:quoting}).
%
We have no way to prove once and for all that the expressions built by
\ls$canon_monoid$ correctly denote the terms, but this fact can be
proven automatically at each application of the tactic, by simple unification.
%
The tactic then applies the lemma \ls$monoid_reflect m e_1 e_2$,
and the goal is changed to \ls$mldenote m (flatten e_1) ==$ \ls$mldenote m (flatten e_2)$.
%
Finally, by normalization, each side will be
canonicalized by running \ls$flatten$ and \ls$mldenote$.


%
% Of course, the canonicalized expressions are proven equal to their
% originals by the tactic.
%
The \ls$canon_semiring$ tactic follows a similar approach,
and is similar to existing reflective
tactics for other proof assistants~\cite{ring, GregoireM05},
except that it only canonicalizes up to linear arithmetic,
as explained above.
%
% 
The full VC for \ls$poly_multiply$
% \ch{What function? Now the boring details have been inserted in the middle.}
contains many other facts, e.g.,
that \ls$p$ is non-zero so the division is well-defined
and that the postcondition does indeed hold.
%
These obligations remain in a ``skeleton'' VC that is also easily proven
by Z3.
%
This proof is much easier for the programmer to write and much more
robust, as detailed ahead in \autoref{sec:canonicalization}.
%
The proof of Poly1305's other main lemma, \ls$poly_reduce$, is also
similarly well automated.

\paragraph{\bf Tactic proofs without SMT}

Of course, one can verify
\ls$poly_multiply$ in Coq, following the same conceptual proof used in \metafstar,
but relying on tactics only.
%
Our proof
\iffull
(see Fig.~\ref{fig:coqpolymultiply} in the appendix)
\else
(included in the appendix)
\fi
is 27 lines long, two of which involve the use of Coq's \ls$ring$
tactic (similar to our \ls$canon_semiring$ tactic) and \ls$omega$ tactic
for solving formulas in Presburger arithmetic. 
%
The remaining 25 lines include steps to destruct the
propositional structure of terms, rewrite by equalities, enriching the
context to enable automatic modulo rewriting (Coq does not fully
automatically recognize equality modulo $p$ as an equivalence relation
compatible with arithmetic operators).
%
While a mature proof assistant like Coq has libraries and tools to
ease this kind of manipulation, it can still be
verbose.

In contrast, in \metafstar{} all of these mundane parts of a proof are
simply dispatched to the SMT solver, which decides
linear arithmetic efficiently, beyond the quantifier-free Presburger fragment
supported by tactics like \ls$omega$, handles congruence closure natively, etc.
%
% \guido{I would delete the following sentence, we're repeating it:
% \st{One value proposition of \textbackslash metafstar is that it enables the
% seamless combination of SMT and tactic based proofs, simplifying the
% work of the programmer/tactic author by delegating the many
% SMT-automatable tasks (and only those tasks) to an SMT solver.}}


\iffalse
\subsection{Completing Partial Programs with Metaprograms}
\label{sec:type classes}
\guido{This subsection does not really fit the flow of the subsection, I'd
delete it and redistribute as appropriate}

Tactics in \metafstar are an instance of a more general class of
metaprogramming that supports inspecting, constructing, transforming,
and completing \fstar (meta)programs, i.e., \metafstar
supports metaprogramming at arbitrary rank. In this section, we
illustrate how metaprograms are tightly integrated within
\fstar, enabling entirely in ``user space'' type system features
like type class resolution.

\paragraph{Programs with holes}
%
As in other dependently
typed languages, \fstar programs contain ``holes'' (missing program
fragments) that the type-checker fills using a (higher-order)
unification-based type inference algorithm. A hole
$\Goal{\Gamma}{_}{t}$ is a pair of a typing context $\Gamma$ and a
type $t$, the task of the type-checker being to find a solution \ls$e$
such that the $\Goal{\Gamma}{e}{t}$ judgment holds in \fstar's type
system. Source programs implicitly contain holes, e.g., \ls$id 0$, an
application of the polymorphic identity function
%
\ls$id: #a:Type -> a -> a$ (in \fstar, `\ls$#$' denotes an implicit argument),
contains an implicit hole for
the missing type argument. Programs may also explicitly mention a
hole \ls$_$ when they expect a term to be determined by inference.



\paragraph{Filling holes with metaprograms}
%
\metafstar allows filling holes with terms constructed by arbitrary
metaprograms.
%
Specifically, the user can write \ls$_ by m$ to cause \metafstar to run
the metaprogram \ls$m$ over the hole $\Goal{\Gamma}{\_}{t}$, where \ls$t$
is the expected type at the current program $\Gamma$.
%
For instance, in an \ls$int$-typed context, \ls$_ by exact `(1 + 1)$
fills the hole with the term \ls$1 + 1$, produced by the metaprogram
which solves the hole using a quotation of the syntax of \ls$1 + 1$.



\paragraph{Controlling implicit argument resolution}
%
\metafstar also supports abstracting holes and the metaprograms that
solve them via \emph{meta-implicits}.
%
Function arguments can be annotated with a metaprogram meant to provide
instantiations for the argument, should unification leave it unresolved.
%
The notation \ls$(#[tau] x : t)$ binds the name 
\ls$x$ as an implicit argument with type \ls$t$, and associates
the metaprogram \ls$tau$ with it.
%
If an instantiation for \ls$x$ was not found by unification,
the type-checker runs the metaprogram \ls$tau$ for the hole
$\Goal{\Gamma}{\_}{t}$ on the calling environment $\Gamma$ and uses the
solution as the instantiation for \ls$x$.
%
As a simple example, the function \ls$let diag (x:int) (#[exact (dquote x)] y : int) = (x, y)$,
which uses a dynamic kind of quotation,
effectively has a default value for \ls$y$, namely the value of \ls$x$.
Therefore,
\ls$diag 0$ evaluates to \ls$(0,0)$, while \ls$diag 0 #1$ evaluates
to \ls$(0,1)$.



\paragraph{Generating top-level definitions}
%
Besides solving particular program holes, \metafstar
can generate entire program fragments automatically.
%
The top-level declaration \ls$%splice m$ runs the metaprogram \ls$m$
over the hole \Goal{\emptyset}{_}{(list def)}, that is, with an
empty local scope and the goal to generate a list of
(representations of) top-level definitions, which are then spliced into
the program.
%
Combined with the ability for reflecting over most language constructs
(terms, global and local contexts, top-level declarations, inductive types, ...)
this can serve as a powerful automation tool and eliminate significant boilerplate.
%
A useful first example use is deriving a printer for an inductive type \ls$t$,
which we do via the call \ls$%splice (mk_printer (`t))$.
%
The (straightforward) metaprogram \ls$mk_printer$ receives a static
quotation of a datatype name, looks it up in the global context and
inspects its structure, generating syntax that case-analyzes the value.





% Aside from this, \metafstar can also be used to automatically generate
% boilerplate for functions like \ls$zero$ and \ls$add$ shown above. For
% instance, the following code:
% %
% \begin{lstlisting}
% noeq type monoid a = { zero  : a; plus  : a -> a -> a }
% %splice (mk_class (`monoid))
% \end{lstlisting}
% %
% defines the class of monoids (as a dictionary of its operations) and
% the metaprogram following it
% %
% generates top-level implementations for \ls$zero$ and \ls$add$ at the
% signatures shown previously by simply projecting and applying the
% appropriate fields from the monoid dictionary.
% \guido{Actually we need \ls$zero,add$ to have signatures with meta-implicits,
% was this deleted on purpose? I'll bring it back if not}


%
% For instance, consider the following functions generic over a monoid
% stucture \ls$m$:
% %
% \begin{lstlisting}
% val zero : #a:Type -> (#m : monoid a) -> a
% val add : #a:Type -> (#m : monoid a) -> a -> a -> a
% let sum_list #a (#m : monoid a) (l : list a) = fold_left add zero l
% \end{lstlisting}
% %
% In the body of \ls$sum_list$, the calls to \ls$add$ and \ls$zero$
% require two implicit arguments.
% %
% The first one is readily instantiated by unification to \ls$a$ as
% expected, a consequence of typing constraints.
% %
% For \ls$m:monoid a$, however, there are no constraints, and unification
% cannot help in finding a solution for \ls$m$.
% %
% Nevertheless, in many cases, the user indeed wants to instantiate
% \ls$m$ simply by passing the \ls$m$ taken by \ls$sum_list$.
% %
% Many proof-assistants implement type class systems in order to
% permit this automation.

% With \metafstar, programmers are able to make this kind of decisions by
% using \emph{meta-implicits}.
%

\paragraph{User-defined type class resolution}
%
By combining meta-implicits, \metafstar's reflection capabilities, and
splicing, we implement type classes within \fstar fully as a user-level
feature, except for some syntactic sugar.
%
We use the notation \ls$[|ty|]$ to denote the binder
\ls$(#[tcresolve] _ : ty)$.
%
The \ls$tcresolve$ metaprogram, defined in \fstar's standard library
inspects the local and global contexts, 
and attempts to solve \ls$ty$ by applying values marked with the
(user-defined) \ls$instance$ attribute via a straightforward
backtracking search.
%
The syntax \ls$class ty = { .. }$ is desugared into the two declarations
\ls$type ty = { .. }$ and \ls$%splice (mk_class (`ty))$.
%
The first declaration defines a plain record type.
%
The \ls$%splice$ invocation calls into a metaprogram
creating projectors for the record, but where the
argument is marked as a type class (i.e. \ls$[| ty |]$),
essentially metaprogramming the class methods.
%
Both metaprograms and the changes to the parser account for about
100 lines of code, yet provide a flavor very close to a full type class
experience.


% \iffull
% The \ls$mk_class$ metaprogram also marks instances with an
% (\fstar-specific) attribute \ls$tcnorm$.
% %
% Essentially, \ls$tcnorm$ causes the type-checker to unfold
% the solution before generating verification conditions.
% %
% This is particularly important in a program verifier since,
% for example, addition over integers is monotonic on both
% arguments, but that property is not captured (or even stateable!)
% at the class interface.
% %
% Without it, simple ground-level overloading would hinder
% verification significantly.

% Instances of a class are simply values in this type marked with the
% \ls$instance$ attribute (an opaque alias to \ls$()$), e.g.:
% \begin{lstlisting}
% [@instance] let _ : monoid int = { zero = 0 ; plus = (+) }
% \end{lstlisting}

% The automation need not stop here.
% %
% Metaprograms can also splice instance declarations--although we have not
% done so for the time being.
% %
% Our type class system currently makes no guarantees about the uniqueness
% or order of the resolved instances.
% %
% When it is desirable to, for example, have a particular preference among
% them, a user can easily override the default resolution metaprogram.
% \guido{Well.. but not the sugar}.
% \fi

\fi

\subsection{Tactics for Entire VCs and Separation Logic}
\label{sec:seplogic}

A different way to invoke \metafstar is over an entire VC.
% in a more ``hammer''-like fashion. -- CH: not true, hammers only
% solve goals, which is not what the rest of this para says
%
While the exact shape of VCs is hard to predict, users with some
experience can write tactics that find and solve particular
sub-assertions within a VC, or simply massage them into shapes better
suited for the SMT solver. We illustrate the idea on proofs for
heap-manipulating programs.

\guido{We should be careful on a previous criticism:
are we doing SL or using the standard heap model?
I think we wanted to claim that is advantage to interop
with existing code}

One verification method that has eluded \fstar until now is
separation logic, the main reason being that the pervasive ``frame
rule'' requires instantiating existentially quantified heap variables,
which is a challenge for SMT solvers, and simply too tedious for users.
%
With \metafstar, one can do better.
%
We have written a (proof-of-concept) embedding of separation logic
and a tactic (\ls$sl_auto$) that performs heap frame inference automatically.

The approach we follow consists of designing the WP specifications for
primitive stateful actions so as to make their footprint
syntactically evident.
% in a consistent shape.\ch{What is ``consistent shape'' supposed to mean?
%   syntactically evident?}
%
The tactic then descends through VCs
%transparently
% \ch{What is ``transparently'' supposed to add?}
until it finds an existential for heaps arising from the frame rule.
%
Then, by solving an equality between heap expressions (which requires
canonicalization, for which we use a variant of \ls$canon_monoid$
targeting {\em commutative} monoids)
the tactic finds the frames and instantiates the existentials.
%
\ch{I
  expect that previous tactics for SL in Coq would also do this,
  although I also guess they have to do a lot more. Didn't Danel
  add some references at some point?}%
%
Notably, as opposed to other tactic frameworks for separation
logic~\cite{Appel06, McCreight09, KrebbersTB17, nmb08htt}, this is \emph{all}
our tactic does before dispatching to the SMT solver, which can now be
effective over the instantiated VC.
%
% Conveniently, the tactic need not to do more than instantiating
% frames, since the rest of the VC can easily be proven by the SMT
% solver.


We now provide some detail on the framework.
%
Below,
`\ls$emp$' represents the empty heap,
`$\bullet$' is the separating conjunction and
`\ls$r |-> v$' is the heaplet with the single reference \ls$r$ set to value \ls$v$.%
\footnote{This differs from the usual presentation where these three operators
are heap \emph{predicates} instead of heaps.}
%
Our development distinguishes between a ``heap''
and its ``memory'' for technical reasons, but we will treat
the two as equivalent here.
%
%\guido{Aseem: do you buy it?}
%
Further, \ls$defined$ is a predicate discriminating valid heaps
(as in~\cite{NanevskiVB10}), i.e., 
those built from separating conjunctions of
\emph{actually} disjoint heaps.


We first define the type of WPs and present the WP for the
frame rule:
%
\begin{lstlisting}
let pre = memory -> prop (* predicate on initial heaps *)
let post a = a -> memory -> prop (* predicate on result values and final heaps *)
let wp a = post a -> pre (* transformer from postconditions to preconditions *)

let frame_post (#a:Type) (p:post a) (m_0:memory) : post a =
  fun x m_1 -> defined (m_0 <*> m_1) /\ p x (m_0 <*> m_1)
let frame_wp (#a:Type) (wp:wp a) (post:post a) (m:memory) =
  exists m_0 m_1. defined (m_0 <*> m_1) /\ m == (m_0 <*> m_1) /\ wp (frame_post post m_1) m_0
\end{lstlisting}
%
Intuitively, \ls$frame_post p m_0$ behaves as the postcondition \ls$p$
``framed'' by \ls$m_0$, i.e., 
\ls@frame_post p m$_0$ x m$_1$@ holds when the
two heaps \ls$m_0$ and \ls$m_1$ are disjoint and \ls$p$ holds over the result value \ls$x$ and the conjoined heaps.
%
Then, \ls$frame_wp wp$ takes a postcondition \ls$p$ and initial heap
\ls$m$, and requires that \ls$m$ can be split into disjoint subheaps
\ls$m_0$ (the footprint) and \ls$m_1$ (the frame), such that the
postcondition \ls$p$, when properly framed, holds over the footprint.

In order to provide specifications for primitive
actions we start in small-footprint style.
%
For instance, below is the WP for reading a reference:
%
\begin{lstlisting}
let read_wp (#a:Type) (r:ref a) = fun post m_0 -> exists x. m_0 == r |-> x /\ post x m_0
\end{lstlisting}
%
We then insert framing wrappers around such small-footprint WPs when
exposing the corresponding stateful actions to the programmer, e.g.,
\begin{lstlisting}
val (!) : #a:Type -> r:ref a -> STATE a (fun p m -> frame_wp (read_wp r) p m)
\end{lstlisting}
%
% %
% For instance, this is the specification for reading a reference:
% %
% \begin{lstlisting}
% val (!) #a (r:ref a) : ST a (fun post m0 -> exists (x:a). m0 == (r |> x) /\ post x m0)
% \end{lstlisting}
% %
% Following the ``small footprint'' style, the function states as a
% precondition that the initial heap must contain a single reference,
% \ls$r$, and returns its value without modifying the heap.
% %
% Of course, this function should also be useable when the heap contains
% more than one reference: that is when the frame rule comes in.
%
To verify code written in such style, we annotate the corresponding 
programs to have their VCs processed by \ls$sl_auto$. For instance, 
for the \ls$swap$ function below, the tactic  
successfully finds the frames for the four occurrences of
the frame rule and greatly reduces the solver's work.
%
%\da{used to be 'both' occurrences, but there are four of them, right? 2 x !, 2 x :=?}
%
Even in this simple example, not performing such instantiation would cause
the solver to fail.
%
\begin{lstlisting}
let swap_wp (r_1 r_2 : ref int) =
    fun p m -> exists x y. m == (r_1 |> x <*> r_2 |> y) /\ p () (r_1 |> y <*> r_2 |> x)
let swap (r_1 r_2 : ref int) : ST unit (swap_wp r_1 r_2) by (sl_auto ()) =
    let x = !r_1 in let y = !r_2 in r_1 := y; r_2 := x
\end{lstlisting}
%
%

% The basic mode of operation of the \ls$sl_auto$ tactic is finding, via
% syntax inspection, where the frame rule is applied and solving frame
% equalities over heap expressions.
% %
% For this purpose, we use the \ls$canon_monoid$ tactic, a part of
% \ls$canon_semiring$ shown above, specialized to the monoid of heaps.


% Separation logic provides another case study illustrating the
% benefits of combining SMT-based automation and
% metaprogramming. Here, the main obstacle for
% SMT-only proofs is the inference of heap frames. To address this,
% we have implemented a proof-of-concept encoding of
% separating conjunction and the frame rule in \fstar, and use
% tactics to infer the heap frames, while leaving the rest
% (involving classical logical connectives and even
% inductive predicates) for SMT. Below we explain our development at
% a high-level eliding some implementation details.
%

%% which in its expanded form looks as follows (\ls{m_1 $\bullet$ m_2} is the union of two non-overlapping memories):
%% \begin{lstlisting}
%% exists m_1 x. def (r |-> x $\bullet$ m_1) /\ m_0 == r |-> x $\bullet$ m_1 /\ def (r |-> y $\bullet$ m_1) /\ post () (r |-> y $\bullet$ m_1)
%% \end{lstlisting}

The \ls$sl_auto$ tactic:
(1) uses syntax inspection to unfold and traverse the goal
until it reaches a \ls$frame_wp$---say, the one for \ls$!r_2$;
(2) inspects \ls$frame_wp$'s first explicit argument (here \ls$read_wp r_2$) to
compute the references the current command requires (here \ls$r_2$);
(3) uses unification variables to build a memory expression describing the
required framing of input memory
(here \ls{r_2$\hspace{0.01cm}$ |-> ?u_1 $\bullet$ ?u_2})
and instantiates the existentials of \ls$frame_wp$ with these unification variables;
(4) builds a goal that equates this memory expression with \ls$frame_wp$'s
third argument (here \ls{r_1$\hspace{0.01cm}$ |-> x $\bullet$ r_2$\hspace{0.01cm}$ |-> y});
and (5) uses a commutative monoids tactic
(similar to \autoref{sec:non-linear})
% \da{used to be reference to sec:canonicalization but \ls$canon_monoid$ description is in sec:non-linear}
with the heap algebra
(\ls{emp}, \ls{$\bullet$}) to canonicalize the equality
and sort the heaplets.
%
Next, it can solve for the unification variables component-wise,
instantiating \ls$?u_1$ to \ls$y$ and \ls$?u_2$ to \ls$r_1 |-> x$, 
and then proceed to the next \ls{frame_wp}.
%
%Finally, once the unification variables are solved, the tactic proceeds to the
%next \ls{frame_wp}.

%% combines
%% syntax inspection of the proof goals with a commutative
%% monoid tactic (which closely follows the \ls$canon_semiring$
%% tactic from \autoref{sec:semirings}).
%% Combined with SMT-based automation for the rest, \ls{sl_auto}
%% Here the proof obligation is to show that
%% the annotated WP subsumes the one computed by \metafstar. \ls$sl_auto$
%% achieves this by relying on the framing of the memory of the
%% initial heap being done at the basic commands using \ls$frame_wp$.

In general, after frames are instantiated, the SMT solver can
efficiently prove the remaining assertions, such as the obligations
about heap definedness.
%
Thus, with relatively little effort, \metafstar brings an
(albeit simple version of a) widely used
yet previously out-of-scope program logic (i.e., separation logic) into \fstar.
%
To the best of our knowledge, the ability to \emph{script} separation
logic into an SMT-based program verifier, without any primitive
support, is unique.



% \guido{Removed a para here} CH: okay
\iffalse
We have also written more advanced, but for now less automated proofs for
functional correctness of various operations on mutable singly
linked lists. %% These proofs use tactics
%% to apply the induction principle for lists, work through the VC manually
%% instantiating
%% heap frames, and send all other goals to the SMT.
In these proofs, the heap frame inference is done by simple tactics
(but not yet \ls$sl_auto$),
while reasoning about heap definedness and
(un)folding of inductive predicates is again
%% functions involving the (standard inductive)
%% functional lists used in specifications is also
left to the SMT, thus demonstrating the power of combining metaprogramming and
SMT-based automation.
\fi


%% Generalizing \ls$sl_auto$ to
%% handle branching and (annotated) recursive calls is future
%% work, as is automating the linked lists proofs;
%% for both, we expect the generality of \metafstar to
%% be of great help.






\subsection{Metaprogramming Verified Low-level Parsers and Serializers}
\label{sec:minilowparse}
% \nik{Budget: 2 pages}

% \nik{Tahina: My general advice for this section is to go top-down.
% Start by showing \ls$gen_specs `t$ from the last section without saying how it works.
% Then show how a further tactic transforms the spec into an equivalent low-level
% implementation to be extracted by kremlin. That should set the user's expectations
% for the section: i.e., don't write proofs at all, just generate low-level code equivalent to high-level spec.

% Then, say a bit about how it works: 0. term quotation; 1. syntax
% inspection; 2. syntax transformation; 3. unquotation via \ls$_ by t$.
% All driven by a careful design of a library of combinators.

% It's not yet evident to me how ``custom verification'' / increment
% generation of terms plays into this. But it would be good to bring it out.

% See revised budget. 2 pages.}

% \guido{Not very important, but can we use a concise name for
% both parser/serializer? ``marshalling pair''? ``parsing pair''?}

Above, we used \metafstar to manipulate VCs for user-written code.
%
Here, we focus instead on generating verified code automatically.
%
We loosely refer to the previous setting as using ``tactics'', and to
the current one as ``metaprogramming''.
%
In most ITPs, tactics and
metaprogramming are not distinguished; however in a program verifier
like \fstar, where some proofs are not materialized at
all (\autoref{sec:correctness}),
proving VCs of existing terms is distinct from generating new terms.
%
% \nik{I tried to elaborate on what I thought was the intent here ...
% but not sure if I buy it myself}
% \ch{Why is this?}\guido{``tactics'' do
% not generate any code or transform the program, will make it clearer}
% \guido{how does it look now?}

%% For a function \ls$f$ with a precondition, generating a call to it
%% involves a proof obligation.
%% %
%% By default, when writing the code manually, \fstar builds a single large
%% VC containing all such obligations for the top-level function that is
%% being verified.
%% %
%% With \metafstar, one can instead process those obligations with tactics
%% as they appear, obtaining a finer-grained control of the verification
%% process.

Metaprogramming in \fstar involves programmatically generating a
(potentially effectful) term (e.g., by constructing its syntax and
instructing \fstar how to type-check it) and processing any VCs
that arise via tactics. When applicable (e.g., when working
in a domain-specific language), metaprogramming verified code can
substantially reduce, or even eliminate, the burden of manual proofs.

We illustrate this by automating the generation of parsers and
serializers from a type definition. Of course, this
is a routine task in many mainstream
metaprogramming frameworks (e.g., Template Haskell, camlp4, etc). The
novelty here is that we produce imperative
parsers and serializers extracted to C, with proofs that they are memory safe,
functionally correct, and mutually inverse.
%
This section is slightly simplified, more detail can be found
\iffull
in \autoref{sec:parsers-detail}.
\else
the appendix.
\fi

We proceed in several stages. First, we program a library of pure, high-level parser
and serializer combinators, proven to be (partial) mutual inverses of
each other.
%
A \ls$parser$ for a type \ls$t$
is represented as a function possibly returning a \ls$t$ along with
the amount of input bytes consumed.
%
The type of a \ls$serializer$ for
a given \ls$p:parser t$ contains a refinement%
\footnote{
\fstar syntax for refinements is \ls$x:t\{phi\}$, denoting
the type of all \ls$x$ of type \ls$t$ satisfying \ls$phi$.
%
}
stating that \ls$p$ is an
inverse of the serializer.
%
A \ls$package$ is a dependent record of a
parser and an associated serializer.
%
\begin{lstlisting}
let parser t = seq byte -> option (t * nat)
let serializer #t (p:parser t) = f:(t -> seq byte){forall x. p (f x) == Some (x, length (f x))}
type package t = { p : parser t ; s : serializer p }
\end{lstlisting}
%
Basic combinators in the library include constructs for parsing
and serializing base values and pairs, such as the following:
\begin{lstlisting}
val p_u8 : parse u8
val s_u8 : serializer p_u8
val p_pair : parser t1 -> parser t2 -> parser (t1 * t2)
val s_pair : serializer p1 -> serializer p2 -> serializer (p_pair p1 p2)
\end{lstlisting}
% val p_nlist : n:nat -> parser t -> parser (nlist n t)
% val s_nlist : n:nat -> #p:parser t -> serializer p -> serializer (nlist n t)
%
%
Next, we define low-level versions of these combinators, which work over
mutable arrays instead of byte sequences.
%
These combinators are coded in the \lowstar subset of \fstar (and so can be
extracted to C) and are proven to both be memory-safe
and respect their high-level variants.
%
The type for low-level parsers,
\ls$parser_impl (p:parser t)$, denotes
an imperative function that reads from an array of bytes and
returns a \ls$t$,
% \ch{Isn't the whole idea of low-level
%   parsers that they never materialize the t? That the parse
%   is only ghost? Is this a simplified setting?}
behaving as the specificational
parser \ls$p$.
%
Conversely, a \ls$serializer_impl$ \ls$(s:serializer p)$
writes into an array of bytes, behaving as \ls$s$.

Given such a library, we would like to build verified, mutually
inverse, low-level parsers and serializers for specific data
formats. The task is mechanical, yet overwhelmingly tedious
by hand, with many auxiliary proof obligations of a predictable
structure: a perfect candidate for metaprogramming.

%
%% Here, ``proven-correct'' means that the low-level implementation
%% is memory safe, that it matches the specification, and that, at the
%% specification level, the parser and serializer are inverses of each
%% other.
%% %
%% All of this is proven by a metaprogram, using a combination of SMT,
%% dependent type checking based on higher-order unification, and other
%% \metafstar tactics.\guido{check}
%% %
%% The generated code falls in the \lowstar fragment of \fstar and can
%% hence be compiled to C.
%% %
%% As such, this technique allows to write correct C parsers
%% that are impervious to security vulnerabilities such as
%% CloudBleed~\cite{cloudbleed}.

\paragraph{Deriving specifications from a type definition}

% To do so, the metaprogram inspects the syntax of the verified code
% \guido{``which verified code?''}
% written by the user, and uses a predefined set of combinators
% to generate the final result.
% %
% In some cases, this metaprogram can be run directly from a \fstar data
% type, without further user input.

% \guido{btw, interestingly, the original(?) paper on LCF
% tactics~\cite{milnertactics} also provides an example of proving
% a parser to be the inverse of a printer.}

% \nik{  Forward pointer to section 5 say we evaluate our approach against
%   some alternatives.

% For section 5, how can we evaluate this against other approaches?

% Is the plan we discussed on slack for evaluation still in scope?,
% i.e., comparing the two styles for:

% 3. Parsers/Printers: Tahina
%    a. Monolithic syntax building
%    b. Incremental, custom type-checking
% }

% \paragraph{Overview through an example}

Consider the following \fstar type,
representing lists of exactly $18$ pairs of bytes.
%
\begin{lstlisting}
type sample = nlist 18 (u8 * u8)
\end{lstlisting}
%
% \tahina{TODO: replace with sum types once the implementation is ready}
%
The first component of our metaprogram is \ls+gen_specs+, which
generates parser and serializer specifications from a type definition.
%
%% These specifications
%% are marked as ghost code, i.e. they are not intended to be executed.
%% %
%% They ``operate'' on high-level data structures,
%% e.g., sequences of bytes of arbitrary lengths,
%% and are equipped with correctness
%% proofs, where \ls{serializer_spec} (definition elided) states
%% that the serializer and parser are inverses of each other (\ls$x:a & b$ is
%% syntax for dependent pairs):
%
\begin{lstlisting}
let ps_sample : package sample = _ by (gen_specs (`sample))
\end{lstlisting}
%
The syntax \ls$_ by tau$ is the way to call \metafstar for code
generation. \metafstar will run the metaprogram \ls$tau$ and, if
successful, replace the underscore by the result.
%
In this case, the \ls$gen_specs (`sample)$ inspects the syntax of
the \ls$sample$ type (\autoref{sec:quoting})
and produces the package below (\ls$seq_p$ and \ls$seq_s$ are 
sequencing combinators):
%
\begin{lstlisting}
let ps_sample = { p = p_nlist 18 (p_u8 `seq_p` p_u8)
                 ; s = s_nlist 18 (s_u8 `seq_s` s_u8) }
\end{lstlisting}
%
% \nik{Say something about the proof obligation that arises from making this package and how it is discharged.}
% \guido{Is there one? I thought the bijection was the only non-trivial
% obligation, and that came from translating the enum type to a nat. I'm
% confused by the current thread.}
% \nik{Btw, wouldn't it be easier if instead of building parser and serializer separately, we build the pair compositionally by composing packages for each component of the type?}
%
\paragraph{Deriving low-level implementations that match specifications}

From this pair of specifications, we can automatically generate \lowstar
implementations for them:
%
\begin{lstlisting}
let p_low : parser_impl ps_sample.p = _ by gen_parser_impl
let s_low : serializer_impl ps_sample.s = _ by gen_serializer_impl
\end{lstlisting}
%
which will produce the following low-level implementations:
%
\begin{lstlisting}
let p_low = parse_nlist_impl 18ul (parse_u8_impl `seq_pi` parse_u8_impl)
let s_low = serialize_nlist_impl 18ul (serialize_u8_impl `seq_si` serialize_u8_impl)
\end{lstlisting}
%
For simple types like the one above, the generated code is fairly
simple. However, for more complex types, using the combinator library
comes with non-trivial proof obligations. For example, even for a
simple enumeration, \ls$type color =$ \ls$ Red$ \ls$| Green$,
% just\ch{just? is this code that complex? (looks simple to me)}
the parser specification is as follows:
%
\begin{lstlisting}
parse_synth (parse_bounded_u8 2)
            (fun x2 -> mk_if_t (x2 = 0uy) (fun _ -> Red) (fun _ -> Green))
            (fun x -> match x with | Green -> 1uy | Red -> 0uy)
\end{lstlisting}
%
We represent \ls$Red$ with \ls$0uy$ and \ls$Green$ with \ls$1uy$. The
parser first parses a ``bounded'' byte, with only two
values. The \ls$parse_synth$ combinator then expects functions between
the bounded byte and the datatype being parsed (\ls$color$), which
must be proven to be mutual inverses.
%
This proof is conceptually easy,
but for large enumerations nested deep within the structure of other
types, it is notoriously hard for SMT solvers.
%
Since the proof is
inherently computational, a proof that destructs the inductive type
into its cases and then normalizes is much more natural.
%
With our metaprogram, we can produce the term and then discharge
these proof obligations with a tactic \emph{on the spot},
eliminating them from the final VC.
%
We also explore simply tweaking the SMT context, again via a tactic,
with good results.
%
A quantitative evaluation is provided
in~\autoref{sec:eval-minilowparse}.

% The generated parsers and serializers are correct by construction,
% memory-safe, and mutually inverse.
% %
% On extraction, the combinators are inlined to yield
% \lowstar code that can be translated by the KreMLin
% compiler to efficient C code.

% that is not only memory-safe, but also correct with regards to the
% original specifications,
% \ch{Not proved though for the
%   whole of KreMLin, and only on paper. Would drop
%   ``semantics-preserving''.  In general, I think this last part is
%   overclaiming. The KreMLin tool provides no formal guarantees.}

% %
% These \lowstar combinators operate
% on \emph{buffers}, which are \lowstar mutable data structures to
% represent C pointers and arrays \cite{lowstar}.
% %
% %% As before, the memory-safety and functional correctness of the low-level
% %% combinators are baked into their types, so that \ls+gen_parser_impl+
% %% and \ls+gen_serializer_impl+ generate implementations correct by
% %% construction.
% %% %
% The combinators in these definitions are later inlined by KreMLin~\cite{lowstar}
% to produce executable C code.

% The correctness of the combinators are baked in their types, so
% typechecking the resulting terms generated
% by \ls+gen_specs+, \ls+gen_parser_impl+
% and \ls+gen_serializer_impl+ triggers no verification condition
% other than those preconditions required by the specific combinators%
% \guido{\st{and unification constraints, the latter solved by reflexivity}}.
% %
% In
% particular, neither the correctness of the serializers generated
% by \ls+gen_specs+, nor the memory safety and correctness of the
% implementations generated by \ls+gen_parser_impl+, do appear as such
% as verification conditions when \fstar is type-checking the terms
% generated by those tactics.

\iffalse
\paragraph{How they work}

Both tactics operate by syntax inspection: \ls+gen_specs+ inspects
the syntax of the type, whereas \ls+gen_parser_impl+
(resp. \ls+gen_serializer_impl+) inspects the syntax of the parser
(resp. serializer) specification.
% , taking advantage of the combinators
% in which the specifications are supposed to be written with, and
% generating implementations using the corresponding implementation
% combinators.
These tactics work along 3 successive principles: term quotation,
syntax inspection and syntax construction, all of which
we present more formally and extensively in \autoref{sec:quoting}.

\begin{enumerate}
\item \emph{Term quotation:} when the user calls \ls+gen_specs (`palette)+, they
  \emph{quote} \ls$palette$, hence building a value at type \ls+term+
    representing the corresponding \fstar syntax.

\item \emph{Syntax inspection:} the \metafstar \ls+term+ type representing \fstar terms
  is abstract; to make sense out of it, \metafstar offers a safe \emph{view}
    via \ls$inspect$, a metaprogramming primitive. Consider, for
    instance, the sketch below for generating an implementation by inspecting a
    specification.
\begin{lstlisting}
let gen_parser_impl () : Tac unit =
  let hd, tl = app_head_tail (cur_goal ()) in
  match inspect hd, inspect (`parser_impl), tl with
  | Tv_fvar lid1, Tv_fvar lid2, [ t, ...; p, ... ] when lid1 = lid2 ->
    let p' = gen_parser_impl_aux t p in ... | _ -> ...
\end{lstlisting}
    \ls+gen_parser_impl+, taking no arguments,
    inspects the shape of the goal, then calls a recursive auxiliary tactic,
    \ls+gen_parser_impl_aux+, which takes the parser specification and the
    type of parsed values as \ls+term+s as arguments. The mechanism is
    general, and \metafstar allows inspecting \fstar inductive type definitions, which
    allows \ls+gen_specs+ to work on the definition of enumeration types
    such as \ls+color+.

\item \emph{Syntax construction:}
  symmetrically, one can synthesize \ls+term+s from term views by calling the built-in
    \ls$pack$  tactic, a
    mechanism used by all tactics in this section to assemble applications of
    combinators. \metafstar also allows normalizing terms during their construction,
    which we take advantage of when unfolding definitions. For instance, if the
    parser specification we inspect is an application of a free variable that does not
    correspond to a known combinator, then our tactic unfolds its definition,
    and then performs a normalization step to reduce any $\beta\iota$-redexes
    appearing after such unfolding.

% \begin{lstlisting}
%  match inspect hd with
%    | Tv_FVar v -> let p' = norm_term [iota] (mk_app (unfold_fv v) tl) in gen_parser_impl_aux t p' | ...
% \end{lstlisting}

  % Then, \ls+gen_parser_impl+ assembles combinators by constructing an application, building on the following \ls+mk_app+ application which builds a n-ary application as a cascade of binary applications, using the \metafstar \ls+pack+ tactic that turns a term view into a reified term:
% \begin{lstlisting}
% let rec mk_app (u: term) (l: list (term * aqualv)) : Tac term =
  % match l with | [] -> u | v :: l' -> let u' = pack (Tv_App u v) in mk_app u' l'
% \end{lstlisting}
% Thus, to turn a parser specification \ls+parse_filter p f+ into its
% corresponding implementation, \ls+gen_parser_impl_aux+ only needs to
% turn the parser specification \ls+p+ into its implementation, and
% then assemble the implementation using the \ls+parse_filter_impl+
% combinator. To do that, it first inspects the parser specification
% given as a reified term \ls+p0+, then checks whether it is a parser
% specification based on \ls+parse_filter+ (by comparing the head to
% the reified term obtained by quoting \ls+parse_filter+), and if so,
% builds the actual term application.
% \begin{lstlisting}
% let gen_parser_impl_aux (t0: term) (p0: term) : Tac term =
  % let (hd, tl) = app_head_tail p0 in
  % if term_eq hd (`parse_filter) then match tl with
  % | [(t, Q_Implicit); (p, Q_Explicit); (f, Q_Explicit)] ->
  %   let p' = gen_parser_impl_aux t p in
  %   mk_app (`parse_filter_impl) [(t, Q_Implicit); (p, Q_Implicit); (p', Q_Explicit); (f, Q_Explicit); ]
  % | _ -> fail "Not enough arguments to parse_filter"
  % else ...
% \end{lstlisting}
 % \metafstar also allows normalizing reified terms during their construction,
    % which we take advantage of when unfolding definitions. In the following
    % excerpt of \ls+gen_parser_impl_aux+, if the parser specification is an
    % application of a free variable that does not correspond to a known
    % combinator, then our tactic unfolds its definition, and then performs a
    % normalization step to reduce any $\beta\iota$-redexes appearing after such
    % unfolding:
% \begin{lstlisting}
 % match inspect hd with
  %  | Tv_FVar v -> let p' = norm_term [iota] (mk_app (unfold_fv v) tl) in gen_parser_impl_aux t p' | ...
% \end{lstlisting}

\end{enumerate}
\fi


% Finally, once the abstract syntax for our parsers has been built, they
% are inserted into the \fstar program using the \ls$_ by$ construct, of
% which \ls$_ by gen_specs ...$ is an instance.
  
  % finally, when the user writes:
% \begin{lstlisting}
% let s : parser_impl (dfst ps) = _ by gen_parser_impl
% \end{lstlisting}
% actually, a proof goal is created for \ls+gen_parser_impl+ to ``prove the
% theorem'' \ls+parser_impl (dfst ps)+. As we saw above, \ls+gen_parser_impl+
% calls an auxiliary tactic, \ls+gen_parser_impl_aux+, which generates a reified
% term; so \ls+gen_parser_impl+ just needs to call \ls+exact+ on the result of
% the auxiliary tactic, which solves this goal by turning the reified term
% produced by \ls+gen_parser_impl_aux+ into an actual \fstar term:
% \begin{lstlisting}
% let p' = gen_parser_impl_aux t p in exact p'
% \end{lstlisting}

% Unquotation can also be mixed with syntax construction: if \ls+t+ is a reified
% term, then \ls+`#t+ is its corresponding unquoted \fstar term. So the
% \ls+mk_app+ invocation above can be replaced with:
% \begin{lstlisting}
% let gen_parser_impl_aux (t0: term) (p0: term) : Tac term =
  % let (hd, tl) = app_head_tail p0 in
  % if term_eq hd (`parse_filter) then match tl with
  % | [(t, Q_Implicit); (p, Q_Explicit); (f, Q_Explicit)] ->
  %   let p' = gen_parser_impl_aux t p in
  %   `(parse_filter_impl #(`#t) #(`#p) (`#f))
  % | _ -> fail "Not enough arguments to parse_filter"
  % else ...
% \end{lstlisting}

% \metafstar syntax inspection and quotation constructs also support binders, as illustrated in another example on monadic transformation of printers that we present in the appendix and the supplemental material.


\section{The Design of \metafstar}
\label{sec:design}

Having caught a glimpse of the use cases for \metafstar, we now turn to its
design.
%
As usual in proof assistants (such as Coq, Lean and Idris), \metafstar
tactics work over a set of goals and apply primitive actions to
transform them, possibly solving some goals and generating new goals in
the process.
%
Since this is standard, we will focus the most on describing the aspects
where \metafstar differs from other engines.
%
We first describe how metaprograms are modelled as an effect
(\autoref{sec:tac-effect})
and their runtime model (\autoref{sec:executing}).
%
We then detail some of \metafstar's syntax inspection and building
capabilities (\autoref{sec:quoting}).
%
Finally, we show how to perform some (lightweight) verification of
metaprograms (\autoref{sec:specifying-metaprograms}) within \fstar.


%
% Two \fstar features make particularly strong appearances here.
% %
% First, \fstar can perform verification of effectful programs via a
% proof-irrelevant \emph{weakest-precondition (WP)} calculus, where most
% specifications are automatically inferred by the type-checker, and
% are suitably-shaped for discharging via
% an SMT solver~\cite{fstar-pldi13,mumon}.
%
% Here, we make use of \fstar's user-extensible effect system, by which a
% monadic definition suffices to introduce an effect~\cite{dm4free}.
% %
% We use this feature in order to define a metaprogramming effect,
% \ls$TAC$, and how its computations are interpreted.
%
%
%% We conclude with a brief summary of \metafstar syntax-inspection
%% capabilities. %%  which, since it must support native compilation,
%% take a peculiar flavour.

%  \iffalse
%  \subsection{A short introduction to \fstar's effect system and logic}
%  
%  \nik{This is actually quite long, nearly a full page. Can we trim it
%  to the bare minimum and explain needed background
%  as and when needed in the remaining text. e.g., proof irrelevance is
%  needed until much later}
%  \guido{Agree. Will trim as much as possible. I can introduce most things
%  via \ls$TAC$}
%  
%  At the core of \fstar is a dependently typed language of pure total
%  functions~\cite{mumon}.
%  %
%  This language also forms the logic for \fstar's program specifications
%  and their proofs.
%  %
%  In addition, \fstar also provides primitive support for divergent
%  computations.
%  %
%  Layered on top of this core is an extensible effect system
%  which allows users to program new effects on top of purity and
%  divergence simply by providing monadic definitions.
%  %
%  A type-and-effect system keeps track of the effects of terms, ensuring
%  proper encapsulation.
%  %
%  In particular, it keeps effectful computations separate from the logical
%  core which, unsurprisingly, demands purity for consistency.
%   
%  Besides tracking the type and effect of programs, \fstar also
%  infers and checks arbitrarilly-fine-grained \emph{specifications} of
%  programs, given in the form of weakest-precondition (WP) predicate
%  transformers.
%  %
%  This requires, for each effect, a WP calculus at the type level,
%  known as a \emph{Dijkstra monad}~\cite{fstar-pldi13}.
%  %
%  For a term \ls$e$ to inhabit a computation type \ls$M t wp$, it is
%  required not only that it returns a value in \ls$t$ and its effects are
%  bounded by \ls$M$, but further that \ls$wp$ correctly specifies its runtime
%  behaviour.
%  %
%  The exact interpretation of ``correctly specifies'' varies slightly for
%  each effect \ls$M$.
%  %
%  In particular, when \ls$e : PURE t wp$ holds, \ls$wp$ is of type
%  \ls$(t -> prop) -> prop$ and, given any
%  postcondition \ls$p : (t -> prop)$ such that \ls$wp p$ holds,
%  then \ls$e$ will run to a value in \ls$v:t{p v}$. % such that \ls$p v$ holds.
%  %
%  In general, \ls$wp$ takes a postcondition \ls$post$ and computes a sufficient
%  precondition for \ls$post$ to hold on the computation's final configuration.
%  
%  
%  
%  %This Dijkstra monad is, roughly, a composition of $m$ with the
%  %continuation monad \ls$(a -> prop) -> prop$.
%  
%  %that isgeneric in the effect
%  %\guido{what does ``generic in the effect'' mean?}
%  
%  %and based on an adaptation of Dijkstra's predicate
%  %transformers to higher-order programs~\cite{dm4free, fstar-pldi13}.
%  
%  
%  
%  Whereas the system computes WPs for computations, programmers can can
%  use a notation similar to Hoare Type Theory~\cite{nmb08htt}, which is
%  easily desugared into WPs.
%  %
%  A couple of common abbreviations are shown below.
%  \begin{lstlisting}
%  effect Pure a (requires pre) (ensures post) = PURE a (fun k -> pre /\ forall x. post x ==> k x)
%  effect Lemma (requires pre) (ensures post) = Pure unit (requires pre) (ensures (fun () -> post))
%  \end{lstlisting}
%  %
%  When a user annotates a program with a WP, \fstar checks that the
%  specifications inferred for these programs are subsumed by the
%  programmer-given ones.
%  %
%  This subsumption check is reduced to a logical validity problem
%  that was until now discharged using the Z3 SMT solver,
%  and, as of this work, also metaprogrammed tactics.
%  %
%  % \guido{How about: that \fstar discharges through SMT solving.
%  % In the case the SMT cannot prove the provided formula, the only mean
%  % the user had to transform it was selectively calling lemmas
%  % to introduce new facts (and subobligations) into the logical context
%  % of the program. This work adds a much more expressible mean,
%  % metaprogrammed tactics.}\cpc{I think this is still too much.  Introducing new lemmas isn't a way to discharge the satisfiability problem, its a way to tweak it.  I agree that it's connected to tactics at a very high level, but I think it makes things less clear to even mention it here.}
%  
%  % \citet{dm4free} describe in detail the means by
%  % which \fstar computes VCs for effectful programs, although a full
%  % understanding of \fstar's VC generation algorithm is not necessary for
%  % this paper---we summarize the main concepts below.
%  
%  % Even pure computations in \fstar are given such computation
%  % types. Specifically, the typing judgment
%  % %
%  % \ls@$\Gamma \vdash$ e : PURE a wp@ states that \ls$e$ is pure
%  % computation that evaluates, for any postcondition
%  % \ls$post : a -> Type$ for which \ls$squash (wp post)$ is valid,
%  % to a value of type \ls$x:a{post x}$.  In other words,
%  % \ls$wp : (a -> Type) -> Type$ is
%  % a predicate transformer that computes a proof-irrelevant, sufficient
%  % precondition of \ls$e$ with respect to any \ls$post$.
%  % % \ch{Having a prop
%  % % type would help make the point that this is
%  % % proof-irrelevant. \nik{yep, thought of that too, but didn't want to
%  % % have to introduce that was well} }
%  % %
%  
%  \fi

\subsection{An Effect for Metaprogramming}
\label{sec:tac-effect}

\metafstar tactics are, at their core, programs that transform the
``proof state'', i.e. a set of goals needing to be solved.
%
As in Lean~\cite{EbnerURAM17} and
Idris~\cite{ChristiansenB16}, we define a monad combining exceptions
and stateful computations over a proof state,
along with actions that can access internal components such as the
type-checker.
%
For this we first introduce abstract types for the proof state,
goals, terms, environments, etc., together with functions to
access them, some of them shown below.

% \ch{We should probably credit Lean, Mtac, Idris around here, not
%   only in the related work section.}
% \guido{Did some of that, marking differences too}

% \fstar will then generate both a concrete representation of the effect
% (used to extract and run effectful programs) and a WP calculus for it
% (used to specify and verify them).

% At the heart of \metafstar, as for many other tactic engines, lies the notion
% of a ``proofstate''.
% %
% Informally, it corresponds to all internal state of
% \metafstar that is relevant for executing metaprograms and obtaining
% solutions.
% %
% Its single most important component is the set of goals that are yet to
% be solved, each of which is represented by an environment, a goal type, and
% a unification metavariable for the solution.
% %
% The proofstate also contains the state of each metavariable (in a persistent unionfind graph), configuration
% and debugging flags, and some local state available to metaprograms.
% \guido{unsure we want to mention that state, lset/lget are very hacky}

% The type \ls$proofstate$ is a term-level representation of this internal
% \fstar structure.
% %
% Importantly, it is exposed to the metaprograms as an abstract type, so metaprograms
% cannot use or construct
% \ls$proofstate$ values except through a safe interface---part of it
% shown below.
% %
% For instance, metaprograms can obtain the set of goals of a proofstate.
% %
% Goals, and many other meta-level concepts, follow this same pattern.

\vspace{-5mm}
\noindent
\begin{minipage}[t]{.50\textwidth}
\begin{lstlisting}
type proofstate
type goal
type term
type env
\end{lstlisting}
\end{minipage}
\begin{minipage}[t]{.50\textwidth}
\begin{lstlisting}
val goals_of : proofstate -> list goal
val goal_env  : goal -> env
val goal_type : goal -> term
val goal_solution : goal -> term
\end{lstlisting}
\end{minipage}
%
% \begin{wrapfigure}{r}{0.3\textwidth}
% \begin{lstlisting}
% type proofstate
% type goal
% type env
% type term
% val goals_of : proofstate -> list goal
% val goal_env  : goal -> env
% val goal_type : goal -> term
% val goal_solution : goal -> term
% \end{lstlisting}
% \end{wrapfigure}
%
% \noindent
% \[
% \arraycolsep=1pt
% \begin{array}{l@{\qquad\qquad}lcl}
%     \kw{type proofstate} & \kw{val goals_of}      & \kw{:} & \kw{proofstate -> list goal} \\
%     \kw{type goal}       & \kw{val goal_env}      & \kw{:} & \kw{goal -> env} \\
%     \kw{type env}        & \kw{val goal_type}     & \kw{:} & \kw{goal -> term} \\
%     \kw{type term}       & \kw{val goal_solution} & \kw{:} & \kw{goal -> term} \\
% \end{array}
% \]
%
We can now define our metaprogramming monad: \ls$tac$.
%
It combines \fstar's existing effect for potential divergence (\ls$Div$),
with exceptions and stateful
computations over a \ls$proofstate$.
%
The definition of \ls$tac$, shown below, is straightforward and given in
\fstar's standard library.
%
Then, we use \fstar's effect extension capabilities~\cite{dm4free}
in order to
elevate the \ls$tac$ monad and its actions to an effect, dubbed
\ls$TAC$.
 
\noindent
\begin{lstlisting}
type error = exn * proofstate (* error and proofstate at the time of failure *)
type result a = | Success : a -> proofstate -> result a | Failed  : error -> result a
let tac a = proofstate -> Div (result a)
let t_return #a (x:a) = fun ps -> Success x ps
let t_bind #a #b (m:tac a) (f:a -> tac b) : tac b = fun ps -> ... (* omitted, yet simple *)
let get () : tac proofstate = fun ps -> Success ps ps
let raise #a (e:exn) : tac a = fun ps -> Failed (e, ps)
new_effect { TAC with repr = tac ; return = t_return ; bind = t_bind
                                 ; get = get ; raise = raise }
\end{lstlisting}
%
% \aseem{Definition of \ls{tac a} contains \ls{Div}? Isn't it
% just \ls{M}, a dummy marker, and divergence comes from the fact that
% the effect is not specified \ls{total}?}\guido{Yes... but I guess we
% don't want to go there. Morally we are taking \ls$M$ to be \ls$Div$,
% IMO}
%
The \ls$new_effect$ declaration introduces \emph{computation types} of the
form \ls$TAC t wp$, where \ls$t$ is the return type and \ls$wp$ a specification.
%
However, until \autoref{sec:specifying-metaprograms} we shall only use
the derived form \ls$Tac t$, where the specification is trivial.
%
% \aseem{Have we
% talked about effects being part of a lattice with lifts etc.? If not,
% may be just say add it to the type environment?}\guido{good point, reworded}
%
%
These computation types are distinct from their underlying monadic
representation type \ls$tac t$---users cannot directly access the
proof state except via the actions.
%
The simplest actions stem from the \ls$tac$ monad definition:
\ls$get : unit -> Tac proofstate$ returns the current
proof state and
%\aseem{Return type of \ls{get} should be \ls{proofstate}?}
%\guido{oops :)}
%\nik{rephrased only for formatting reasons ... this was overflowing the right margin}
\ls$raise: exn -> Tac 'a$
fails with the given exception%
\footnote{We use greek letters $\alpha$, $\beta$, ...
to abbreviate universally quantified type variables.}.
Failures can be handled using 
\ls$catch : (unit -> Tac 'a)$ \ls$ -> Tac (either exn 'a)$,
which resets the state on failure, including that of unification
metavariables.
%
We emphasize two points here.
%
First, there is no ``\ls$set$'' action.
%
This is to forbid metaprograms from arbitrarily replacing their
proof state, which would be unsound.
%
% Instead, the proof state can only be modified by primitive actions that
% are guaranteed to preserve soundness (see \autoref{sec:correctness}).
%
Second, the argument to \ls$catch$ must be thunked, since in \fstar{}
impure un-suspended computations are evaluated before they are passed
into functions.\ch{Just that the CBN normalizer doesn't do this, and
  we use the CBN normalizer for tactics, which are super impure
  (because of Div and quoting observing evaluation order). What a ugly hack!}
%
% Since in \fstar impure code is call-by-value,\ch{the usual confusing
%   claims about CBV when our default normalizer is CBN, and we explain
%   later in the paper that we use this CBN normalizer for tactics(!)
%   ... added ``impure code'' but not sure if it's the best fix, since
%   tactics are impure}
%
% All combinators shown subsequently follow this style.

The only aspect differentiating \ls$Tac$ from other
user-defined effects is the existence of effect-specific primitive actions,
which give access to the metaprogramming engine proper.
%
We list here but a few:
%
\begin{lstlisting}
val trivial : unit -> Tac unit $\quad$ val tc : term -> Tac term $\quad$ val dump : string -> Tac unit
\end{lstlisting}
%
All of these are given an interpretation internally by \metafstar.
%
For instance, \ls$trivial$ calls into \fstar's logical simplifier to check
whether the current goal % (i.e. the one at the top of the stack)
is a
trivial proposition and discharges it if so, failing otherwise.
%
The \ls$tc$ primitive queries the type-checker
to infer the type of a given term in the current % goal's
environment
% and return its type (%as in many dependently-typed languages,
(\fstar types are a kind of terms, hence the codomain of \ls$tc$ is also \ls$term$).
%
This does not change the proof state; its only purpose
is to return useful information to the calling metaprograms.
% , who can then inspect it freely.
% (\autoref{sec:quoting}).
%
Finally, \ls$dump$ outputs the current proof state to the user in a
pretty-printed format, in support of user interaction.


Having introduced the \ls$Tac$ effect and some basic actions,
writing metaprograms is as straightforward as writing any other \fstar
code.
%
For instance, here are two metaprogram combinators.
The first one repeatedly
calls its argument until it fails, returning a list of all the
successfully-returned values. The second one behaves similarly,
but folds the results with some provided folding function.
%
\begin{lstlisting}
let rec repeat (tau : unit -> Tac 'a) : Tac (list 'a) =
  match catch tau with   | Inl _ -> []   | Inr x -> x :: repeat tau

let repeat_fold f e tau = fold_left f e (repeat tau)
\end{lstlisting}
 
These two small combinators illustrate a few key points
of \metafstar.
%
As for all other \fstar effects, metaprograms are written in
applicative style, without explicit \ls{return}, \ls{bind}, or \ls{lift}
of computations (which are inserted under the hood).
%
This also works across different effects: \ls$repeat_fold$ can seamlessly combine
the pure \ls$fold_left$ from \fstar's list library with a metaprogram
like \ls$repeat$.
%
Metaprograms are also type- and effect-inferred: while
\ls$repeat_fold$ was not at all annotated, \fstar infers the polymorphic
type
\ls$('b -> 'a -> 'b) -> 'b -> (unit -> Tac 'a) -> Tac 'a$
for it.
%
% In fact, it even infers a specification for it beyond \ls$Tac$.


It should be noted that, if lacking an effect extension feature, one
could embed metaprograms simply via the (properly abstracted) \ls$tac$
monad instead of the \ls$Tac$ effect.\ch{As Lean actually does?}
%
It is just more convenient to use an effect, given we are working within
an effectful program verifier already.
%
% It also has the advantage that users do not need to learn an entirely
% new language to write them. -- CH: off-topic here
%
In what follows, with the exception
of~\autoref{sec:specifying-metaprograms} where we describe
specifications for metaprograms, there is little reliance on
using an effect; so, the same ideas could be
applied in other settings.\ch{just that it's often not our ideas,
  and the rest have already been applied in Lean, Idris, and Mtac}
  \guido{reworded to ``the same ideas''}

\subsection{Executing \metafstar Metaprograms}
\label{sec:executing}

% Metaprograms are run during typechecking time, and have no runtime
% presence.
% %
% This phase separation is enforced by \fstar's effect system: programs
% meant for execution must be in the \ls$ALL$ effect, and there is no lift
% or elimination from \ls$Tac$ into \ls$ALL$, nor can users provide one.
% %
% \guido{I'm not sure this bit is technically accurate. Also unclear
% why we emphasize that they have no runtime presence.}

Running metaprograms involves three steps.
%
First, they are \emph{reified}~\cite{dm4free} into their underlying
\ls$tac$ representation, i.e. as state-passing functions.
%
%
% These are state-passing functions which can be
% safely\ch{safely?} reduced within a
% reduction machine, and even\ch{even?} extracted, e.g. to OCaml.
%
User code cannot reify metaprograms: only \fstar
can do so when about to process a goal.
 
Second, the reified term is applied to an initial proof state, and then
simply evaluated according to \fstar's dynamic semantics, for instance
using \fstar's existing normalizer.
%
For intensive applications,
such as proofs by reflection,
we provide faster alternatives (\autoref{sec:evaluation}).
%x
In order to perform this second step, the proof state, which
up until this moments exists only internally to \fstar,
must be \emph{embedded} as a term, i.e., as abstract syntax.
%
Here is where its abstraction pays off: since metaprograms cannot
interact with a proof state except through a limited interface, it
need not be \emph{deeply} embedded as syntax.
%
By simply wrapping the internal proofstate into a new kind of ``alien''
term, and making the primitives aware of this wrapping, we can readily
run the metaprogram that safely carries its alien proof state around.
%
This wrapping of proof states is a constant-time
operation.
%
% This same principle is applied in \autoref{sec:evaluation}, where we
% describe a novel\guido{check} design for multi-language
% interoperation which greatly increases the performance
% of \textbackslash fstar (and thus \textbackslash metafstar) computations run by the type-checker.


The third step is interpreting the primitives.
%
They are realized by functions of similar types implemented within the
\fstar type-checker, but over an internal \ls$tac$ monad and the concrete
definitions for \ls$term$, \ls$proofstate$, etc.
%
Hence, there is a translation involved on every call and return,
switching between embedded representations and their
concrete variants.
%
Take \ls$dump$, for example, with type \ls$string -> Tac unit$.
%
Its internal implementation, implemented within the \fstar type-checker,
has type
\ls$string -> proofstate ->$
\ls$Div (result unit)$.
%
When interpreting a call to it, the interpreter must \emph{unembed} the
arguments (which are representations of \fstar terms)
into a concrete string and a concrete
proofstate to pass to the internal implementation of \ls$dump$.
%
The situation is symmetric for the return value of
the call, which must be \emph{embedded} as a term.
%
% Fortunately, for the initial and final proof states the embedding and
% unembedding is very cheap, as explained above.

% %
% For efficiency reasons, the term-level \ls$proofstate$
% does not include
% this state, which is taken from a rather large union-find graph within
% \fstar.
% %
% Instead of embedding such graph as a term, we instrument the internal
% implementation of \ls$catch$ with calls to the union-find graph's
% transactional interface to snapshot the graph at entry and roll it back
% on failures.
% %
% By making use of persistent data structures, this only adds an $O(1)$
% overhead in time, provides the expected semantics, and (as \ls$catch$ is
% the only handler) this aspect is not observable by metaprograms.
% %
% (The \ls$fail$ action does not need this instrumentation.)

% \paragraph{\bf Solving goals}

% Ultimately, the objective of a metaprogram is to \emph{solve} its goals
% by applying these primitives.
% %
% A solution to a goal \Goal{\Gamma}{?u}{t} is a fully-defined term \ls$e$
% such that a valid typing derivation \Goal{\Gamma}{e}{t} exists.
% %
% The primitives are implemented so they only discharge goals when
% they can be solved in this manner, except perhaps deferring parts
% of the term as new goals.
% %
% % When executing a metaprogram, the set of unsolved goals is kept
% % in the proofstate, reachable only by the primitives.
% %
% \metafstar enforces that all goals are solved when
% execution finishes, raising an error otherwise.

% %\aseem{Might be useful to remind the reader the shape of the
% %goals: \ls{G |- _ : t}}\guido{good now?}
% %
% % To allow metaprograms to effectively solve goals, we expose most \fstar
% % typing rules as metaprogramming primitives.
% %
% The set of primitives exposed represents a large subset
% of \fstar's typing rules.
% %
% By having direct and near full access to the \fstar typing judgment, along
% with the syntax inspection capabilities of \metafstar, metaprogrammers
% can script the typing rules in whichever way they see fit without being
% able to corrupt the language's semantics (we expand on this ahead\guido{check}).

% %Throughout the execution of a metaprogram, the proofstate carries the set of open goals;
% %
% %and a metaprogram can only successfully finish when this set becomes empty.

% For a concrete example, consider:
% \ls$let f : int -> int = _ by (intro (); exact (`42))$.
% %
% Here, \metafstar will create an initial goal of shape
% \Goal{\kw{Gamma}}{?u_1}{int -> int}, pack it into a proofstate, and
% begin executing the metaprogram.
% %
% By first applying the \ls$intro$ primitive, which implements \fstar's typing
% rule for $\lambda$-abstractions, \metafstar creates a new goal of shape
% \Goal{\kw{Gamma, x:int}}{?u_2}{int}, and solves \ls$?u_1$ by setting
% it to \ls$fun x -> ?u_2$.\footnote{In fact, \ls$?u_2$ is also renamed
% accounting for a locally-nameless representation, but we largely omit
% such details.}
% %
% %(and not very relevant at this level of detail).\guido{just forget
% %about it? I'm a bit scared of getting comments about this, such
% %as the questions on how we handled uvars}
% %
% The proofstate is also modified to ``shift focus'' to the new \ls$?u_2$ goal.
% %
% Next, running \ls$exact$ solves \ls$?u_2$ with the integer literal 42
% (denoted by a static quotation, which are standard and whose details we omit),
% after checking that 42 indeed has type \ls{int}.
% %
% After checking that no unification metavariable remains unsolved
% (otherwise an error would have been raised),
% \metafstar the entire invocation solves the goal with
% \ls$fun x -> 42$, correctly typed at its expected
% type of \ls$int -> int$.
% %

% All these primitive actions do is apply typing rules in backwards
% order, i.e. ``conclusion to premise'', creating new goals
% for each premise.
% %
% Of course, if the hole's type had not been an arrow at the time of
% applying \ls$intro$, the primitive would fail with an error, since the
% corresponding typing rule would not apply.
% %
% Similarly, using \ls$(`true)$ instead of \ls$(`42)$ would have failed.

\subsection{Syntax Inspection, Generation, and Quotation}
\label{sec:quoting}

If metaprograms are to be reusable over different kinds of goals,
they must be able to reflect on the goals they are invoked to solve.
%
Like any metaprogramming system, \metafstar offers a way to inspect and
construct the syntax of \fstar terms.
%
Our representation of terms as an inductive type, and the variants
of quotations, are inspired by the ones
in Idris~\cite{ChristiansenB16} and Lean~\cite{EbnerURAM17}.
%
\ch{The section doesn't talk about efficiency:
``As this inspection happens very often, care must be taken to make it
efficient.''}
%
% A particular requirement of \metafstar is the interoperation with
% other languages (\autoref{sec:evaluation}).
% Supporting this feature requires making some design choices so that
% inspecting and constructing abstract syntax works efficiently
% and consistenly on them.\guido{Not sure what I mean here..
% some operations are not allowed in native code, but the view is designed
% in order to safely compile an link into \fstar.}
%
%
% Our design, especially the quotation mechanism, is inspired by
% Lean's~\cite{EbnerURAM17}.\guido{Not sure we can, or want, to state this. Lean
% has an elaborator and a very well-defined notion of pre-expression.
% AFAICT, it's very different.}

% As we've seen, terms are converted to their representation, and
% vice-versa, on many calls to primitives; hence, this operation
% must be made efficient.
%
%

% Here we present the methods available to metaprograms for
% inspecting and constructing abstract syntax.
% %
% Other metaprogramming systems expose very expressive quotation
% mechanisms

%
% whether or not
% they are in the goal, would be more direct.
% \cpc{Well, we've just seen an example of with\_pat. I'd drop the comment about mapply}
%

\paragraph*{\bf Inspecting syntax} Internally, \fstar uses
a locally-nameless representation~\cite{ChargueraudLocallyNameless}
with explicit, delayed
substitutions.
% \guido{among other non-semantic choices?}.
%
To shield metaprograms from some of this internal bureaucracy, we expose
a simplified view~\cite{WadlerViews} of terms.
%
Below we present a few constructors from the \ls$term_view$
type:

\vspace{-5mm}
\noindent
\begin{minipage}[t]{.42\textwidth}
\begin{lstlisting}
val inspect : term -> Tac term_view
val pack : term_view -> term
\end{lstlisting}
\end{minipage}
\begin{minipage}[t]{.58\textwidth}
\begin{lstlisting}
type term_view =
  | Tv_BVar   : v:dbvar -> term_view
  | Tv_Var    : v:name  -> term_view
  | Tv_FVar   : v:qname -> term_view
  | Tv_Abs    : bv:binder -> body:term -> term_view
  | Tv_App    : hd:term -> arg:term -> term_view
  ...
\end{lstlisting}
\end{minipage}
%
% \guido{Also, in the library  \ls$Tv_BVar$ and \ls$Tv_Var$ both take
% a \ls$bv$, which represents both in the same type. I don't mind this presentation,
% but maybe we need to be faithful}
%
% The actual \ls$term$ type is kept opaque.
%
The \ls$term_view$ type provides the ``one-level-deep'' structure of a
term: metaprograms must call \ls$inspect$ to reveal the structure of the
term, one constructor at a time.
%
The view exposes three kinds of variables: bound variables, \ls$Tv_BVar$;
named local variables \ls$Tv_Var$; and top-level fully qualified names,
\ls$Tv_FVar$.
%
Bound variables and local variables are distinguished since the internal
abstract syntax is locally nameless.
%
For metaprogramming, it is usually simpler to use a fully-named
representation, so we provide \ls$inspect$ and \ls$pack$ functions
that open and close binders appropriately to maintain this invariant.
%
Since opening binders requires freshness, \ls$inspect$ has effect \ls$Tac$.%
\footnote{We also provide functions \ls$inspect_ln$, \ls$pack_ln$
which stay in a locally-nameless representation and are thus pure, total
functions.}
%
As generating large pieces of syntax via the view easily becomes
tedious, we also provide some ways of \emph{quoting} terms:

\iffalse
% In order to avoid the verbosity the view introduces, we instrumented
% \fstar's type-checker to inject calls to \ls$inspect$ and \ls$pack$
% automatically whenever needed, including in pattern matches.
% %
% As a result, most metaprograms do not need to explicitly coerce one to
% the other, allowing to build and match on terms as if the view was
% fully deep.
% %
% \guido{example? also quite green}%
% \nik{How does this view compare with what Lean does? Would be good to compare a bit. IIRC, they don't have a view and%
% do opening/closing explicitly, but that was a while back. Maybe they've acquired a view since?}%
% \guido{Their expr type is deep, so they don't use a view (LN or not). I think they can%
% only do this if (1) there are translations in-place even for compiled code (2) their expr%
% is exactly their internal AST, which would be surprising/}%
% %

% Note that we could have defined the type \ls$term$ of quoted terms
% as a deep one, instead of this shallow view, and allow to obviate
% the \ls$inspect$ing and \ls$pack$ing.
% %
% However, this would hinder performance as terms are passed between
% the type-checker and metaprograms constantly during tactic execution.
% %
% Importantly, while these coertions could be eliminated by rigging
% \fstar's normalizer to perform term as needed,

% While \ls$inspect$
% %\guido{Was: ensures that the invariants of well-scoped terms are maintained,}
% allows for a convenient analysis and construction of terms
% from metaprograms,
% sometimes it is convenient to directly inspect
% locally nameless terms: so, we provide a function, \ls$inspect_ln$, to
% do that, but we set it aside for the rest of this
% paper.
% %
% \guido{Really set it aside? I was hoping to say that this is one is
% indeed pure and can be used in proof contexts. \cpc{If we don't discuss it further, I'd suggest to move this to a footnote}}
% %
% Dually, \ls$pack$ packages the term view back into a term.
% %
% This function does not need to generate any fresh names and so is not effectful.

% \cpc{Can the results be meaningfully different though? I mean, can
% tactic programs realize that there's a difference? If not, or if it's
% only a technical issue, I'd drop everything after --- \guido{Yes they can. And even besides
% tactics you would prove False anyway if the thing is pure.}
% and also because it fails if it encounters an ill-scoped term

% Metaprograms can then inspect the internal structure of the opaque
% \ls$term$s by calling \ls$inspect$, and to construct them via
% \ls$pack$.


% Besides simply exposing the head constructor,
% %
% This is the reason \ls$inspect$ is marked in the \ls$Tac$ effect -- if
% marked \ls$Pure$ it would destroy soundness by returning two slightly
% different views on the same term.

% The \ls$inspect$ function implicitly
% opens binders with fresh names, switching terms to a fully named
% representation lazily---it must be marked effectful because of the fresh names it
% generates, which means that calling it twice can produce different results.
% % \cpc{Can the results be meaningfully different though? I mean, can
% % tactic programs realize that there's a difference? If not, or if it's
% % only a technical issue, I'd drop everything after --- \guido{Yes they can. And even besides
% % tactics you would prove False anyway if the thing is pure.}
% and also because it fails if it encounters an ill-scoped term
\fi

% \ch{The structure around here is strange. We introduce quotations only
%   to return to the original question and solve it without any
%   quotations. So the order should change.
% \guido{I tried to make it more about syntax in general, without
% relying a lot on that example.}}
% \ch{Also quotations are not so easy to understand, so they need good
%   examples too and could maybe be pushed to the end of this subsection?
% \guido{I could reorganize but maybe the cool pattern matcher example should be last?}}

% \nik{We should call out better here what the novel elements
% of what we do are. e.g., if we are just reusing Lean's ideas for
% quotation, then we should say so and provide less detail. Overall, we
% should try to make this sound less like the description of a system
% implementation}
% \guido{Good point, some of our choices are only so because of needing
% to compile. Going over it later.}
% \guido{Still unsure on what to keep...}
% \ch{I think at the very least proper credit should be given for
%   these quotations}
%   \guido{Brief mention above}


\paragraph*{\bf Static quotations}
A static quotation \ls$`e$ is just a shorthand for statically calling
the \fstar{} parser to convert \ls$e$ into the abstract syntax of
\fstar terms above.
%
For instance, \ls$`(f 1 2)$ is equivalent to the following,
%
\begin{lstlisting}
pack (Tv_App (pack (Tv_App (pack (Tv_FVar "f"))
                              (pack (Tv_Const (C_Int 1)))))
               (pack (Tv_Const (C_Int 2))))
\end{lstlisting}

\paragraph*{\bf Dynamic quotations}
A second form of quotation is
\ls$dquote: #a:Type -> a ->$ \ls$Tac term$, an effectful
operation that is interpreted by \fstar's normalizer during
metaprogram evaluation.
%
It returns the syntax of its argument at the time
\ls$dquote e$ is evaluated.
%
Evaluating
\ls$dquote e$ substitutes all the free variables in
\ls$e$ with their current values in the execution environment,
suspends further evaluation, and returns the abstract syntax of the
resulting term. For instance, evaluating
%
\ls$(fun x -> dquote (x + 1)) 16$ produces the abstract
syntax of \ls$16 + 1$.
% \nik{picked some arbitrary syntax below ... discuss and fix impl}\cpc{We should clarify that one is compile-time, and the other is run(well, interpretation)-time.}
%
% Since information from the execution environment is needed to implement
% these quotations, they cannot be made to run natively.\guido{maybe delete}




\paragraph*{\bf Anti-quotations}
Static quotations are useful for building big chunks of syntax concisely,
but they are of limited use if we cannot combine them with existing bits of syntax.
%
Subterms of a quotation are allowed to ``escape'' and be substituted
by arbitrary expressions.
%
We use the syntax \ls$`#t$ to denote an antiquoted \ls$t$, where \ls$t$
must be an expression of type \ls$term$ in order for the quotation to
be well-typed.
%
For example, \ls$`(1 + `#e)$ creates syntax for an addition
where one operand is the integer constant \ls$1$ and the
other is the term represented by \ls$e$.
%
% One can also combine this with dynamic quotes.
% %
% For instance, in the term \ls$let z = 25 in (fun x -> `(y + `#(dquote x))) z$,
% the quotation produces a representation of \ls$y + _$, where the
% underscore is spliced in with a dynamic quotation of \ls$x$, and will
% reduce to a representation of \ls$y + 25$.

\paragraph*{\bf Unquotation}
Finally, we provide an effectful operation,
\ls$unquote: #a:Type ->$ \ls$t:term$ \ls$-> Tac a$, which takes a term
representation \ls$t$
and an expected type for it \ls$a$ (usually inferred
from the context), and calls the \fstar type-checker to check and
elaborate the term representation into a well-typed term.
%
% Note that \ls$unquote #a t$ is similar to \ls$(_ by exact t)$, but
% instead of being triggered by the type-checker, it's the metaprogram
% that causes the transformation.\ch{Didn't get what ``triggered
% by the type-checker'' means.}
%\guido{It's subtle, let's just remove it. First one
%is a term of type \ls$Tac term$, while the second one
%immediately calls into \metafstar to generate the term}


\subsection{Specifying and Verifying Metaprograms}
\label{sec:specifying-metaprograms}

% \ch{Might want to merge this in: -- CH: not needed
% As with all \fstar effects, \ls$TAC$-computations are given \emph{computation
% types} of the form \ls$TAC a wp$, where \ls$a$ is the computation's
% result type and \ls$wp$ is a weakest-precondition transformer, in this
% case of type \ls$tacwp a = proofstate ->$ \ls$(result a -> prop) -> prop$.
% %
% For convenience, we define a \ls$Tac$ effect as a special case of
% \ls$TAC$ with a ``trivial'' WP.
% %
% Until \autoref{sec:specifying-metaprograms}, we only use \ls$Tac$.
% \ch{Why not save this advanced stuff for 3.3 entirely then?}
%   \guido{I wanted to be precise in the use of \ls$Tac$ and \ls$TAC$ here}
% }

Since we model metaprograms as a particular kind of effectful program
within \fstar, which is a program verifier, a natural question to ask is
whether \fstar can specify and verify metaprograms.
%
The answer is ``yes, to a degree''.

%
% The way \fstar does program verification is via a type-level WP-calculus
% embedded in computation types: a \emph{Dijkstra monad}~\cite{fstar-pldi13}.
% %
% For effects that are derived from monadic definitions, such as \ls$TAC$,
% a WP-calculus is also derived automatically by the type-checker, and all
% (meta)programs are annotated with a WP in their types.
% %
% \fstar can infer these WPs automatically, usually in full detail.

To do so, we must use the WP calculus for the \ls$TAC$ effect:
\ls$TAC$-computations are given computation types of the form
\ls$TAC a wp$, where \ls$a$ is the computation's result type and
\ls$wp$ is a weakest-precondition transformer of type
\ls$tacwp a$ = \ls$proofstate -> (result a -> prop) -> prop$.
%
However, since WPs tend to not be very intuitive, we first define two variants
of the \ls$TAC$ effect: \ls$TacH$ in ``Hoare-style'' with pre- and
postconditions and \ls$Tac$ (which we have seen before), which
only specifies the return type, but uses trivial pre- and postconditions.
% \ch{This ``simple
%   form'' was already heavily used above, so good flow requires you to
%   refer to that, not confusingly introduce this as a new concept.}
%
The \ls$requires$ and \ls$ensures$ keywords below
simply aid readability of pre- and postconditions---they are identity functions.
%
\begin{lstlisting}
effect TacH (a:Type) (pre : proofstate -> prop) (post : proofstate -> result a -> prop) =
        TAC a (fun ps post' -> pre ps /\ (forall r. post ps r ==> post' r))
effect Tac (a:Type) = TacH a (requires (fun _ -> True)) (ensures (fun _ _ -> True))
\end{lstlisting}
%
Previously, we only showed the simple type for the \ls$raise$ primitive,
namely \ls$exn ->$ \ls$Tac 'a$.
%
In fact, in full detail and Hoare style, its type/specification is:
\begin{lstlisting}
val raise : e:exn-> TacH 'a $\,\!$ (requires (fun _ -> True))
                          (ensures (fun ps r -> r == Failed (e, ps)))
\end{lstlisting}
expressing that the primitive has no precondition, always fails with
the provided exception, and does not modify the proof state.
%
From the specifications of the primitives,
and the automatically obtained Dijkstra monad,
\fstar can already prove interesting properties about
metaprograms.
%
We show a few simple examples.


The following metaprogram is accepted by \fstar as it can conclude,
from the type of \ls$raise$, that the assertion is unreachable, and
hence \ls$raise_flow$ can have a trivial precondition (as \ls$Tac unit$
implies).
%
\begin{lstlisting}
let raise_flow () : Tac unit = raise SomeExn; assert False
\end{lstlisting}
%
For \ls$cur_goal_safe$ below, \fstar verifies that (given the precondition)
the pattern match is exhaustive.
%
The postcondition is also asserting that the metaprogram always succeeds
without affecting the proof state, returning some unspecified goal.
%
Calls to \ls$cur_goal_safe$ must statically ensure that the goal list is
not empty.
%
\begin{lstlisting}
let cur_goal_safe () : TacH goal (requires (fun ps -> ~(goals_of ps == [])))
                                (ensures (fun ps r -> exists g. r == Success g ps)) =
    match goals_of (get ()) with | g :: _ -> g
\end{lstlisting}
%
Finally, the \ls$divide$ combinator below ``splits'' the goals of a
proof state in two at a given index \ls$n$, and focuses a different
metaprogram on each.
%
It includes a runtime check that the given \ls$n$ is non-negative,
and raises an exception in the \ls$TAC$ effect otherwise.
%
Afterwards, the call to the (pure) \ls$List.splitAt$ function
requires that \ls$n$ be statically known to be non-negative,
a fact which can be proven from the specification for
\ls$raise$ and the effect definition, which defines
the control flow.
%\ch{... and the way bind is defined for the TAC monad?}
%
\begin{lstlisting}
let divide (n:int) (tl : unit -> Tac 'a) (tr : unit -> Tac 'b) : Tac ('a * 'b) =
    if n < 0 then raise NegativeN;
    let gsl, gsr = List.splitAt n (goals ()) in ...
\end{lstlisting}
%
This enables a style of ``lightweight'' verification of metaprograms,
where expressive invariants about their state and control-flow can be
encoded.
%
The programmer can exploit dynamic checks (\ls$n < 0$) and exceptions
(\ls$raise$) or static ones (preconditions), or a mixture of them, as
needed.


Due to type abstraction, though, the specifications of most primitives cannot provide
complete detail about their behavior, and deeper specifications (such as
ensuring a tactic will correctly solve a goal) cannot currently be proven,
nor even stated---to do so would require, at least, an internalization
of the typing judgment of \fstar.
%
While this is an exciting possibility~\cite{AnandBCST18},
we have for now only focused on
verifying basic safety properties of metaprograms, which helps users
detect errors early, and whose proofs the SMT can handle well.
%
Although in principle, one can also write tactics to
discharge the proof obligations of metaprograms.


% In this case, the programmer has annotated that \ls$repeat$
% takes as argument any metaprogram with trivial preconditions
% (i.e. ``total''), and always succeeds, again without any
% meaningful precondition.
% %
% \fstar is able to verify that \ls$repeat$ meets this specification easily.
% %
% Third, they are type-, effect- and specification-inferred: while
% \ls$repeat_fold$ was not at all annotated, \fstar infers
% the polymorphic type
% \ls$('b -> 'a -> 'b) -> 'b -> (unit -> Tac 'a) -> TAC 'a   _$,
% where we omit the (somewhat complex) inferred WP for it.
% %
% Importantly, this WP indeed implies that the program always suceeds,
% i.e. that it is also a \ls$TacS$.
% %
% %
% \fstar also introduces lifts at the WP level, and hence verification
% is also not affected.

% \guido{Maybe drop this following para or dump in related work.}
% \nik{need to downplay this bit about abstraction, since our stuff needs abstraction too for safety
% \guido{how about ``opaque'"? I mean it cannot be *seen*, beyond not being
% able to modify/create values}\guido{how about now}}
% \guido{A bit unsure about MetaCoq}
% \guido{this definitely not here, but maybe somewhere:


% Lean, closer to our design, exposes a very similar monad, but marks
% it with the \texttt{meta} keyword, which is used to distinguish and
% separate meta-level constructs from object-level ones.
% %
% By reusing fstar's effect system, we can avoid the need for such a
% distinction and keep metaprograms more tightly integrated with the
% rest of the language. \guido{``as in..?''}}


% \paragraph{\bf Other approaches}
% %% It is worth noting that the transparent effect definition is what
% %% allows for this partial verification of metaprograms.
% %
% While Idris and MetaCoq also expose a monadic interface, the
% monad's definition is kept opaque.
% %
% In \metafstar, however, exposing the representation to the
% type-checker allows for evaluation of metaprograms
% (by reifying them to their representation), while the
% derived WP calculus enables lightweight verification.
% %% allows to reason about the success and failure of metaprograms---the WPs
% %% are simply an effective way to do so, but not crucial.
 

%% However, \fstar does not try to discharge every written \ls$assert$.
%% %
%% This is for two reasons:
%% (1) it is usually much more efficient to collect many assertions into a single SMT call
%% (2) more crucially, assertions depend on the control flow of the program.
%% %
%% Concretely, it is only because of (2) that \ls$fail_flow$ above succeeds.
%% %
%% The actual VC generated for that program (were it not trivially solved
%% by \fstar internals) involves proving \ls$False$ under the assumption
%% that the \ls$fail$ call succeeded which, as reflected by
%% its specification, is impossible.
%% 
%% More intricate \fstar programs contain many of these assertions in
%% arbitrarily-complex flows.
%% %
%% As a result, the verification conditions that need to be discharged
%% for them are of considerable size, and not easy to understand by
%% programmers.
%% %
%% For instance, the VC for the following implementation of the factorial
%% function:
%% \begin{lstlisting}
%% let rec fact (n:nat) : Tot nat = if n = 0 then 1 else n * fact (n -- 1)
%% \end{lstlisting}
%% %
%% is similar to the following:
%% \begin{lstlisting}
%% forall (k:(nat -> prop)). (forall (x:nat). k x <==> True) ==>
%%      n = 0 ==> (forall (y:int). y == 1 ==> k y)
%%    /\ n $\ne$ 0 ==> b2t (n -- 1 >= 0) /\ n -- 1 < n
%%                  /\ (forall (y:nat{y << n}).  y == n -- 1 ==>
%%                         (forall (fy:nat).  fy == fact y ==>
%%                             (forall (r:int).  r == n * fy ==> b2t (r >= 0) /\ (forall (r':nat). r' == r ==> k r))))
%% \end{lstlisting}
%% which is combining
%% (1) that the function indeed returns a non-negative integer
%% (2) that the recursive calls are made on the domain of the function, \ls$nat$
%% (3) that the recursive calls are made on strictly smaller arguments.
%% %
%% As the reader might intuit, trying to predict the exact shape of the VC
%% for more complex programs quickly becomes untractable..
%% %
%% Even if multiple assertions are combined, ideally the user should be
%% able to apply different tactics to individual assertions, without
%% considering the ambient context.
%% 
%% 



% The exact sequencing of these actions is undetermined:
% metaprogrammers can combine these steps in whatever way they see fit,
% without needing to change the type-checker, making metaprograms akin to
% ``type-checker scripts''.
% \ch{what does compiler scripts mean? (never heard this!)
%   and what does this have to do with compilation?
%   are you again calling the type-checker a compiler?
%   This is super confusing, please stop doing this!}

% we would like
% to not have to trust metaprograms, and instead guarantee that any such
% program can never corrupt the internal state of the type-checker, build
% ill-typed terms, or take proofs of false statements as true.
% %
% In this section
% we describe how we


%% One is having a small, trusted type-checker, and making tactics generate
%% some independently-checkable proof object;
%% sometimes called the \emph{de Bruijn principle}~\cite{Barendregt:2001:PUD:778522.778527}
%% for proof assistants.

\section{\metafstar, Formally}

We now describe the trust assumptions for \metafstar
(\autoref{sec:correctness}) and then how we reconcile tactics within a
program verifier, where the exact shape of VCs is not given, nor known
a priori by the user (\autoref{sec:vcgen}).

% \subsection{view}
% \guido{Possibly better suited for S5}
% The reader may wonder why the view exists, and we simply do not define
% \ls$term$ as a recursive type.
% %
% During metaprogram execution, there are many (un)embeddings between
% internal terms and their \ls$term$ representation at the \metafstar
% level.
% %
% If these conversions needed to be ``deep'' in the term, metaprogram
% execution on large developments would slow down to a crawl.
% %
% When interpreting metaprograms, this can be avoided by a clever
% pun between \ls$term$s and their views; the interpreter can unfold
% the view on pattern-matches and unembeddings can fold it back
% automatically.
% %
% However, this pun is at odds with compiling metaprograms into native
% code, since they link directly with the \fstar compiler and must have
% the exact same representation of terms---lest another translation,
% forcefully deep, must occur.
% %
% It would also prevent using a fully-named representation as
% this unfolding must be seamless, and hence needs to be pure.
% %
% Using a view permits to satisfy all these constraints harmoniously.

\subsection{Correctness and Trusted Computing Base (TCB)}
\label{sec:correctness}

As in any proof assistant,
tactics and metaprogramming would be rather useless if
they allowed to ``prove'' invalid judgments---care must be taken to ensure soundness.
% measures must
% be taken to reasonably ensure that \metafstar does not violate
% in the formal model of \fstar\guido{which doesn't fully exist..}.
%
We begin with a taste of the specifics of \fstar's static
semantics, which influence the trust model for \metafstar, and then
provide more detail on the TCB.

\paragraph{\bf Proof irrelevance in \fstar}
The following two rules for introducing and eliminating refinement types
are key in \fstar, as they form the basis of its
proof irrelevance.

\[
   \inferrule*[lab=T-Refine]
   {
     \Gamma |- e ~:~ t \quad
     \Gamma |= \phi[e/x]
   }
   {
     \Gamma |- e ~:~ x\!:\!t\{\phi\}
   }
   \qquad
   \qquad
   \inferrule*[lab=V-Refine]
   {
     \Gamma |- e ~:~ x\!:\!t\{\phi\}
   }
   {
     \Gamma |= \phi[e/x]
   }
\]

The $\vDash$ symbol represents \fstar's \emph{validity judgment}~\cite{dm4free}
which, at a high-level, defines a proof-irrelevant,
classical, higher-order logic.
%
These validity hypotheses are usually collected by the type-checker, and
then encoded to the SMT solver in bulk.
%
Crucially, the irrelevance of validity is what permits
efficient interaction with SMT solvers, since
reconstructing \fstar terms from SMT proofs is unneeded.

As evidenced in the rules, validity and typing are mutually recursive,
and therefore \metafstar must also construct validity derivations.
%
In the implementation, we model these validity goals as
holes with a ``squash'' type~\citep{Nogin02,AwodeyBauer},
where \ls$squash phi = _:unit{phi}$, i.e., a refinement of \ls$unit$.
%
Concretely, we model \VGoal{\Gamma}{phi} as \Goal{\Gamma}{?u}{squash phi} using 
a unification variable.
%
\metafstar does not construct deep solutions to squashed goals: if they
are proven valid, the variable \ls$?u$ is simply solved by the \ls{unit} value `\ls$()$'.
%
At any point, any such irrelevant goal can be sent to the SMT solver.
%
Relevant goals, on the other hand, cannot be sent to SMT.

% \guido{\st{Interactions between proof relevance and irrelevance occur in both
% directions: proof-relevant rules (e.g., \textsc{T-Refine}) may have
% proof-irrelevant premises, and vice versa. So, when trying to
% type-check a solution to a proof-relevant goal, fstar may generate
% additional proof-irrelevant goals.}}

% Reconstructing the proofs SMT solver produces, if any, back into a
% proper derivation remains a significant challenge (despite recent
% progress~\cite{BohmeW10,smtcoqCAV2017}).

% As a practical matter, we provide metaprogrammers with different ways to
% handle these particular goals:
% %
% We provide users with three options:
% (1) force these conditions to be valid as they arise;
% (2) add them as new goals for further processing;
% (3) attempt to discharge them automatically using SMT.
% %
% This is done is order to support both a lightwi
% %
% Of course, when transforming a proof-irrelevant goal, a proof-relevant
% goal can easily arise (e.g. from an argument to a lemma); proof-relevant
% goals are always added to the proof-state for further processing.


\paragraph{\bf Scripting the typing judgment}
A consequence of validity proofs not being materialized
is that type-checking is undecidable in \fstar.
%
For instance: does the unit value \ls$()$ solve the hole
\Goal{\Gamma}{?u}{squash phi}?
%
Well, only if $\phi$ holds---a condition which no type-checker can effectively
decide.
%
This implies that the type-checker cannot, in general, rely on proof
terms to reconstruct a proof.\ch{Why? I have no clue why this is related
  to decidability. Isn't higher-order unification in Coq already undecidable?
  Does this mean that Coq cannot rely on proof terms?}
%
Hence, the primitives are designed to provide access to the typing
judgment of \fstar directly, instead of building syntax for proof terms.
%
One can think of \fstar's type-checker as
implementing one particular algorithmic heuristic of the typing
and validity judgments---a heuristic which happens to work well in practice.
%
For convenience, this default type-checking heuristic is also available to
metaprograms: this is in fact precisely what the \ls$exact$ primitive
does.
%
Having programmatic access to the typing judgment
also provides the flexibility to tweak VC generation as needed,
instead of leaving it to the default behavior of \fstar.
%
For instance, the \ls$refine_intro$ primitive implements
\textsc{T-Refine}.
%
When applied, it produces two new goals, including that the refinement
actually holds.
%
At that point, a metaprogram can run any arbitrary tactic on it,
instead of letting the \fstar type-checker collect the obligation
and send it to the SMT solver in bulk with others.

% However, situations appear when this ``default'' type-checker is
% inefficient or unsuccessful at solving a particular goal.
% %
% In such situations, the flexibility of scripting typing derivations, as
% opposed to simply syntax, becomes necessary.

% \ch{so far there was
%   only discussion about building derivations for the validity
%   judgment, is that what this is talking about? or something else?}
% \ch{complete?  why is this related?}\guido{removed}
% %
% \ch{Bad transition follows, what's the connection to previous phrase?
%   Switch from \metafstar to \fstar's typing judgment?}%
% \guido{Fixed? not sure what this is}

% \ch{Sadly I didn't understand much of the previous paragraph.
%   In particular, how exactly is \metafstar{} scripting the typing judgment?
%   Any simple examples?}
% \guido{Added something, not sure how convincing it is.}

\paragraph{\bf Trust}
\label{sec:trust}

There are two common approaches for the correctness of tactic engines:
%
(1) the \emph{de Bruijn
criterion}~\cite{Barendregt:2001:PUD:778522.778527}, which
requires constructing full proofs (or proof terms)
% derivations\ch{actually proof terms, no?}
% \guido{usually yes, but in \fstar that would not suffice
% due to the non-syntax-directed rules such as those above.
% e.g. you could just get \ls$()$}%
and checking them at the end, hence reducing trust to
an independent proof-checker; 
%
and (2) the LCF style, which applies backwards reasoning while
constructing validation functions at every
step, reducing trust to primitive, forward-style implementations of the
system's inference rules.

As we wish to make use of SMT solvers
within \fstar, the first approach is not easy.
%
Reconstructing the proofs SMT solvers produce, if any, back into a
proper derivation remains a significant challenge (even despite recent
progress, e.g.~\cite{BohmeW10,smtcoqCAV2017}).
%
Further, the logical encoding from \fstar to
SMT, along with the solver itself, are already part of \fstar's
TCB: shielding \metafstar from them would not significantly increase
safety of the combined system.

Instead, we roughly follow
the LCF approach and implement \fstar's typing rules
as the basic user-facing metaprogramming actions.
%
However, instead of implementing the rules in forward-style and using
them to validate (untrusted) backwards-style tactics, we implement them
directly in backwards-style.
%
That is, they run by breaking down goals into subgoals,
instead of combining proven facts into new proven facts.
% \da{Any hope of concretely illustrating the forwards-backwards
% difference here, would make this paragraph much more
% understandable?} \guido{better?}
%
Using LCF style makes the primitives part of the TCB.
%
However, given the primitives are sound, any combination of them
also is, and any user-provided metaprogram must be safe due to the
abstraction imposed by the \ls$Tac$ effect, as discussed next.
% \guido{very handwavy}
% %
% \ch{This too (maybe just transitioning to next part would help),
%   but also, again no examples are given, so this won't
%   even be understandable to regular readers.}
%   %
% \guido{Pointed to next para now}

% For instance, the \ls$refine_intro$ primitive
% implements \textsc{T-Refine}, produces two subgoals:
% the two hypotheses of \textsc{T-Refine}.
% %
% The first one is a relevant goal which \emph{must} be instantiated, but the
% second one is irrelevant.
%

%
% \begin{tabular}{ll}
% \begin{lstlisting}
% let False$\,$     = squash empty
% let True$\,$      = squash unit
% let ( /\ ) p q  = squash (p * q)
% let ( \/ ) p q  = squash (either p q)
% \end{lstlisting}
%  &
% \begin{lstlisting}
% let ( ==> ) p q = squash (p -> q)       (* written p ==> q *)
% let ~ p  = p ==> False
% let (forall) #a p = squash (x:a -> p x) (* written forall x. p x *)
% let (exists) #a p = squash (x:a & p x)  (* written exists x. p x *)
% \end{lstlisting}
% \end{tabular}

% \ch{Warning: all these definitions might be short-lived.
%   Our prop proposal would change them.}
% \guido{equality is not here}
%
% This plays well for interfacing with SMT solvers (?)
%

%% Following the de Bruijn approach would involve extracting proofs from
%% the SMT solvers and certifying them within \fstar.
%% %
%% While there is work following that route, it requires
%% significant effort and remains an active area of
%% research \cite{BohmeW10,smtcoqCAV2017}.
%% \guido{does this read like ``yeah we didn't feel like doing that''?}
%% %
%% It is unclear it is also feasible in practice since, at the scale
%% of some \fstar developments, these objects can grow to huge
%% sizes.\guido{number?}
%% %
%% Instead, we design \metafstar to closely follow \fstar's typing system.

% By default, during type-checking, metaprograms are reduced
% using \fstar's normalizer, an abstract machine
% %\cpc{Is ``abstract'' machine the right term here? I'd call it a ``virtual'' machine.}
% for computing normal forms
% of
% \emph{pure} terms.
% %, which is needed to perform dependent type-checking\cpc{suggest removing this lart bit of the sentence}.
% %
% However, metaprograms are, a priori, effectful programs which cannot
% be normalized.
% % \cpc{I think ``safely'' is a strange word, here. We say that the normalizer operates on pure programs; it's not clear what it would mean to feed it effectful programs}
% %
% We avoid this restriction by defining our metaprogramming effect
% monadically as proof-state-passing computations, with divergence and
% exceptions, and then have the type-checker automatically convert it into
% a new computational effect, as described by \citet{dm4free}.
% %
% Such effectful metaprograms can then be \emph{reified} into their
% state-passing representations, which are suitable for reduction
% on \fstar's abstract machine.
% % \ch{So the machine can deal with
% %   non-termination just fine, which contradicts your emphasis above
% %   that it only works on {\em pure} terms} -- ok, explained below
% %
% To run a metaprogram, we apply its reified form to the initial proofstate and start
% the abstract machine.
% %
% The metaprogram will then either diverge, or produce a value in the
% \ls$result$ type described in \autoref{sec:tac-effect}.

%% To expand on primitives modifying the proofstate, let us give a more
%% formal description of it.
%% %
% \guido{I'm a bit wary of this entire explanation. On one hand, it's
% very similar to the standard LCF correctness \cite{milnertactics} but
% \textbf{without} the validation procedures, putting (primitive) tactics
% in the TCB. Maybe more similar to \cite{schmidttactics}. On the other
% hand, we MUST say something about it. Maybe our strenght is that the TCB
% is increased only by the primitives, while composition is guaranteed
% to be safe (LCF tactics predate monads for structuring programs by 10
% years). In any case, even with validation, the impls of the typing rules
% are obviously in the TCB.}

%% For instance, suppose the engine has a current goal of shape
%% \Goal{\Gamma}{_}{psi ==> phi}.
%% %
%% While the witness for this goal, if it exists, must be the unit value
%% \ls$()$, the engine will not accept to solve the goal with \ls$()$
%% unless the logical payload (\ls$psi ==> phi$) is proven (or added as a
%% new goal).


\subsubsection{Correct Evolutions of the Proof State}

For soundness, it is imperative that tactics do not
arbitrarily drop goals from the proof state,
and only discharge them when they are solved,
or when they can be solved by other goals tracked in the proof state.
%
For a concrete example, consider the following program:
%
\begin{lstlisting}
let f : int -> int = _ by (intro (); exact (`42))
\end{lstlisting}
%
Here, \metafstar will create an initial proof state with
a single goal of the form
$\left[\Goal{\emptyset}{?u_1}{int -> int}\right]$ and
begin executing the metaprogram.
%
When applying the \ls$intro$ primitive, the proof state
transitions as shown below.
%
% One way a metaprogram could discharge the goal \Goal{\emptyset}{_}{int}
% is by calling \ls$apply `f$, for some \ls$f$ of type \ls$bool -> int$,
% and then solving the new goal.
% %
% The \ls$apply$ primitive will, after checking that the return type of
% \ls$f$ unifies with the goal, replace the original goal and provide
% a new one for each argument of \ls$f$---in this case, one of type \ls$bool$.
% %
% Concretely, \ls$apply `f$ changes the proof state as follows:
%
\[
    \left[ \Goal{\emptyset}{?u_1}{int -> int} \right]
%
    \leadsto
%
    \left[ \Goal{\kw{x:int}}{?u_2}{int} \right]
\]
%
\let\evol\preceq
%
Here, a solution to the original goal
has not yet been built, since it
\emph{depends} on the solution to the goal on the right hand side.
%
When it is solved with, say, \ls$42$, we can solve our original goal
with \ls$fun x -> 42$.
%a
To formalize these dependencies, we say that a proof state $\phi$
\emph{correctly evolves (via $f$) to $\psi$}, denoted $\phi \evol_f \psi$,
when there is a generic transformation $f$, called a \emph{validation}, from solutions to all of $\psi$'s
goals into correct solutions for $\phi$'s goals.
%
When $\phi$ has $n$ goals and $\psi$ has $m$ goals, the
validation $f$ is a function from $\kw{term}^m$ into $\kw{term}^n$.
% \cpc{A solution to a goal is just a term, right? And in the case of proof-irrelevant goals, just unit? One thing that's not clear up to this point in the paper is how you turn a refinement into a goal}
%
Validations may be composed, providing the
transitivity of correct evolution, and if a proof state $\phi$
correctly evolves (in any amount of steps) into a state with no more goals, then we have fully
defined solutions to all of $\phi$'s goals.
%
We emphasize that validations are not constructed
explicitly during the execution of metaprograms.
%
Instead we exploit unification metavariables to instantiate the
solutions automatically.


%
% Some of them, such as the \ls$flip$ action to reorder goals,
% do so trivially by means of a permutation function.
% %
% Primitives for goal management,
% such as \ls$divide : int -> (unit -> Tac 'a) -> (unit -> Tac 'b) -> Tac ('a * 'b)$,
% which divides the current set of goals in two and applies a
% different metaprogram
% to each side, are correct since validations
% are also
% horizontally composable, i.e., if $\phi_1 \evol_{f_1} \psi_1$ and
% $\phi_2 \evol_{f_2} \psi_2$,
% where $f_1 : \kw{term}^{n_1} -> \kw{term}^{m_1}$,
% then their concatenations also correctly evolve to each other, that is:
% $\phi_1@\phi_2 \evol_{(f_1@f_2)} \psi_1@\psi_2$,
% where $(f_1@f_2)(\bar{e}) = f_1 (\kw{take}~n_1~\bar{e}) @ f_2 (\kw{drop}~n_1~\bar{e})$.

% This formulation deviates from the standard LCF one in that\ch{one in that?
%   is the rest talking about standard LCF or what we do?}
% GM: What we do, but it's not that important, chose to remove
% mention of LCF which is different in more ways too

Note that
validations may construct solutions for more than one goal, i.e.,
their codomain is not a single term.
%
This is required in \metafstar, where primitive steps may not only
decompose goals into subgoals, but actually combine goals as well.
%
Currently, the only primitive providing this behavior is \ls$join$,
which finds a maximal common prefix of the environment of two
irrelevant goals, reverts the ``extra'' binders in both goals and
builds their conjunction.
%
Combining goals using \ls$join$ is especially useful for
sending multiple goals to the SMT solver in a single call.
%
When there are common obligations within two goals, \ls$join$ing them
before calling the SMT solver can result in a significantly faster
proof.

\iffull
The situation when this proves useful is in sending the combined goal to
SMT.
%
The bulk of an SMT solver's work is extending its set of facts
by using logical reasoning and triggering quantifiers via E-matching,
the latter being relatively expensive.
%
Once a fact enters the set, it is very efficient to find its
proof\guido{proof? yikes} which essentially means any particular fact
needs to only be proven once.
%
In the where there are common obligations within two goals, \ls$join$ing
them before calling the SMT solver can result in a significantly faster
proof.
%
Joining everything, however, can degrade performance and predictability.
%
Another alternative is, of course, using tactics to cut via the common
obligation, obtaining a separate goal for it and allowing to assume it
elsewhere.
%
While this works, it requires the programmer to identify the formula
and explicitly cut by it, or write an automated procedure to find these
formulas.
%
With \metafstar, we can simply reuse the solver's good support for
identifying repeats.
\fi

We check that every primitive action respects the $\evol$ preorder.
%
This relies on them modeling \fstar's typing rules.
%
For example, and unsurprisingly, the following rule for typing abstractions
is what justifies the \ls$intro$ primitive:
%\guido{omitting the fact that these are computation types; \nik{I think that's okay}}
%
\[
    % \inferrule[T-Abs]
   \inferrule*[lab=T-Fun]
   {
     % \Gamma |- t : Type_i \quad
     \Gamma, x:t |- e : t'
   }
   {
     \Gamma |- \lambda(x:t).e ~:~ (x:t) -> t'
   }
\]
%
Then, for the proof state evolution above,
% \ch{very fuzzy reference, and this is no longer about checking the
% primitive actions, is it?}\nik{hopefully the reference is less fuzzy
% now; and intro is a primitive, so I don't have a problem with this}
the validation function $f$
is the (mathematical, meta-level) function taking a term of type
$int$ (the solution for \ls$?u_2$) and building syntax for its abstraction over $x$.
%
Further, the \ls$intro$ primitive
% \ch{derived primitive? WTF? And I don't
%   even know what ``derived primitive'' you're talking about.} \nik{fixed?}
%   \guido{looks good to me now}
respects the correct-evolution
preorder, by the very typing rule (T-Fun) from which it is defined.
%
In this manner, every typing rule induces a
syntax-building metaprogramming step.
%
Our primitives come from this dual interpretation of typing rules,
which ensures that logical consistency is preserved.


% Should probably come clean about this:
% In practical terms, these rules are new code that is added
% to the TCB.


Since the $\evol$ relation is a preorder, and every metaprogramming
primitive we provide the user evolves the proof state according $\evol$, it is trivially the
case that the final proof state returned by a (successful)
computation is a correct evolution of the initial one.
%
That means that when the metaprogram terminates, one
has indeed broken down the proof obligation correctly,
and is left with a
(hopefully) simpler set of obligations to fulfill.
%
%
%
% \cpc{Could we perhaps mention that another model would be possible for tactics, in which terms produced by tactics are preserved, and the full term is retype-checked?}
% \guido{I don't think that's true, unless we keep typing derivations. \fstar's typing
% is undecidable, and the type-checker is heuristic, that's one of the advantages
% of having \metafstar.}
%
Note that since $\evol$ is a preorder, \ls$Tac$
provides an interesting example of monotonic state~\cite{preorders}.

% \ch{Maybe add that \ls$Tac$ behaves like a monotonic state
%   effect/monad~\cite{preorders}.  This was said in 2 before, but
%   that was way too early.} \da{Added a sentence on this.}
%   \guido{Worried about the ``novel''... I doubt we're the
%   first to notice and we're certainly not the first scenario
%   where that happens.}

\subsection{Extracting Individual Assertions}
\label{sec:vcgen}

As discussed, the logical context of a goal processed by a
tactic is not always syntactically evident in the program.
%, particularly when using the
%\ls$assert p by t$ form shown in \autoref{sec:vcgen_example}. 
% GM: I don't get that remark.
%
And, as shown in the \ls$List.splitAt$ call in \ls$divide$ from
\autoref{sec:specifying-metaprograms}, some
obligations crucially depend on the control-flow of the program.
%
Hence, the proof state must crucially include these assumptions
if proving the assertion is to succeed.
% \ch{succeed doing what?}\nik{fixed?}
%
Below, we describe how \metafstar finds proper contexts in which
to prove the assertions, including control-flow information.
%
Notably, this process is defined over logical formulae and does not
depend at all on \fstar's WP calculus or VC generator: we believe it
should be applicable to any VC generator.


\iffalse
We've shown examples of the \ls$assert phi by tau$ syntax above,
allowing to preprocess \ls$phi$ via the \ls$tau$ tactic.
%
However, a key detail was missing: how are these assertions
handled by the engine?
%
While the initial environment for \ls$synth$ metaprogramming is
clear (the static typing environment) assertions crucially depend
on the program's control flow, which is not reflected in the typing
environment.

For instance, consider the following function to extract a value
from an option type, falling back to a default one when none
is present.
\begin{lstlisting}
val Some?.v : (o:option 'a{isSome? o}) -> Tot 'a   (* previously defined *)
let get_some (default:'a) (opt:option 'a) : 'a = if isSome? opt then Some?.v opt else default
\end{lstlisting}
To type-check this function, \fstar needs to prove that that, in the
\ls$then$ branch, \ls$o$ effectively satisfies the \ls$isSome?$
predicate.
%
This is only possible by considering the program's flow:
there is no way to deduce it from the static typing context.
%
Rather, this only succeeds because the assertion exists in a
\emph{logical} context, determined by the VC of the program and
containing the control-flow assumptions, postconditions of previous
calls, etc.
%
While this example is rather trivial, more intricate \fstar programs
contain many of these assertions in arbitrarily-complex flows.
%
As a result, their verification conditions are of considerable size and
difficult to foresee precisely by programmers.


Instead, \fstar programmers routinely rely on static assertions, which
are easily accommodated in its VC generation procedure.
%
\fi

% \ch{With newly stuff commented above it's at this point unclear why
%   you're starting to talk about the assert keyword. Pointing to an
%   example that uses assert..by in a way that depends on the context
%   would be good.}%
%   \nik{Fixed?}
%   \guido{I think so}
%
As seen in \autoref{sec:non-linear}, the basic mechanism by which
\metafstar attaches a tactic to a specific sub-goal is
\ls$assert phi by tau$. Our encoding of this expression is built
similarly to \fstar's existing \ls$assert$ construct, which is simply
sugar for a pure function \ls$_assert$ of type
\ls$phi:prop -> Lemma (requires phi) (ensures phi)$, which
essentially introduces a cut in the generated VC.
%
That is, the term \ls$(assert phi; e)$ roughly produces the verification condition
\ls$phi /\ (phi ==> VC_e)$, requiring a proof of \ls$phi$ at this point,
and assuming \ls$phi$ in the continuation.
%
% An expression \ls$assert p; e$
% produces a VC of the form \ls$p /\ (p ==> wp_e post)$, where
% \ls$wp_e post$ is the weakest precondition of \ls$e$ with respect
% to some postcondition \ls$post$.
For \metafstar, we aim to keep this style while allowing asserted
formulae to be decorated with user-provided tactics that are tasked with
proving or pre-processing them. We do this in three steps.
% \da{I see 1 (First), 2(Next), but where is 3?}
% \guido{fixed, see: ``Finally, ..''}

First, we define the following ``phantom'' predicate:
%
% \nik{I think we made it abstract to be defensive; e.g., if the preprocessor hook
% didn't correctly remove all by\_tactics from a VC, we wanted the proof
% to fail in Z3 rather than succeeding silently by Z3 unfolding the
% definition. I'm not sure that's relevant to the discussion here and I
% would be find with remove the qualifier for the paper.}
% irreducible // GM: really abstract, but why? can't remember; NS: Change to abstract?
\begin{lstlisting}
let with_tactic (phi : prop) (tau : unit -> Tac unit) = phi
\end{lstlisting}
%
Here \ls$phi `with_tactic` tau$ simply associates the tactic \ls$tau$
with \ls$phi$, and is equivalent to \ls$phi$ by its definition.
%NS: removed, see comment above
%% By marking it \ls$irreducible$, we ensure it is kept as-is
%% throughout VC generation and simplification.
%
Next, we implement the \ls$assert_by_tactic$ lemma, and desugar
\ls$assert phi by tau$ into \ls$assert_by_tactic phi$\! \ls$tau$.
%
This lemma is trivially provable by \fstar. %from the
%definition of \ls$with_tactic$.
%\guido{``does it unfold or not?'' \nik{yes}}
%
\begin{lstlisting}
let assert_by_tactic (phi : prop) (tau : unit -> Tac unit)
                             : Lemma (requires (phi `with_tactic` tau)) (ensures phi) = ()
\end{lstlisting}
%
Given this specification,
the term \ls$(assert phi by tau; e)$ roughly produces the verification
condition
\ls$phi `with_tactic` tau /\ (phi ==> VC_e)$,
with a tagged left sub-goal,
and \ls$phi$ as an hypothesis
in the right one.
%
Importantly, \fstar keeps the \ls$with_tactic$ marker uninterpreted until the VC
needs to be discharged.
%
At that point, it may contain several annotated subformulae.
%
For example, suppose the VC is \ls$VC0$ below, where we distinguish
an ambient context of variables and hypotheses $\Delta$:
\begin{lstlisting}
(VC0)   $\Delta \models$ X ==> (forall (x:t). R `with_tactic` tau$_{\!1}$ /\ (R ==> S))
\end{lstlisting}
%
% \begin{lstlisting}
% $\Delta \models$ P `with_tactic` t_0 ==> (R `with_tactic` t_1
%                              /\ (forall (x:t). (Q x ==> (H x) `with_tactic` t_2)))
% \end{lstlisting}
%
% The third step involves actually running the tactic \ls$t$ to process
% the asserted formula \ls$p$. Whereas previously every VC produced
% by \fstar would directly be encoded to SMT, \metafstar now provides\cpc{``introduces'' rather than ``now provides''} an
% additional step, where prior to the SMT encoding, it traverses the VC
% evaluating any tactics associated with sub-goals in the entire VC.
%
In order to run the \ls$tau$$_{\!1}$ tactic on \ls$R$, it must first be ``split
out''.
%
To do so, all logical information ``visible'' for \ls$tau$$_{\!1}$
(i.e. the set of premises of the implications traversed and
the binders introduced by quantifiers) must be included.
%
% In \fstar, these two cases in fact coincide since implication
% is simply a (non-dependent) universal quantification.
%
As for any program verifier, these hypotheses include the control flow
information, postconditions, and any other logical fact that is known to
be valid at the program point where the corresponding \ls$assert R by$\!
\ls$tau$$_{\!1}$ was called.
%\guido{say why, hard to explain}
%
All of them are collected into $\Delta$ as the term is traversed.
%
In this case, the VC for $R$ is:
\begin{lstlisting}
(VC1)   $\Delta$, _:X, x:t $\models$ R
\end{lstlisting}
%
%
Afterwards, this obligation is removed from the original VC. This is done
by replacing it with \ls$True$, leaving a ``skeleton'' VC with all
remaining facts.
%
\begin{lstlisting}
(VC2)   $\Delta \models$ X ==> (forall (x:t). True$\hspace{0.005cm}$ /\ (R ==> S))
\end{lstlisting}
%
The validity of \ls$VC1$ and \ls$VC2$ implies that of \ls$VC0$.
%
\fstar also recursively descends into \ls$R$ and \ls$S$, in case there are more
\ls$with_tactic$ markers in them.
%
Then, tactics are run on the the split VCs (e.g.,
\ls$tau$$_{\!1}$ on \ls$VC1$) to break them down (or solve them).
All remaining goals, including the skeleton, are sent to the SMT solver.

Note that while the \emph{obligation} to prove \ls$R$, in \ls$VC1$, is
preprocessed by the tactic \ls$tau$$_{\!1}$, the \emph{assumption} \ls$R$ for the
continuation of the code, in \ls$VC2$, is left as-is.
%
This is crucial for tactics such as the canonicalizer from
\autoref{sec:non-linear}: if the skeleton \ls$VC2$
contained an assumption for the canonicalized equality
it would not help the SMT solver show the uncanonicalized postcondition.

However, not all nodes marked with \ls$with_tactic$
are proof obligations.
%
Suppose \ls$X$ in the previous VC was given as \ls$(Y `with_tactic` tau$$_{\!2}$\ls$)$.
%
In this case, one certainly does not want to attempt to prove
\ls$Y$, since it is an hypothesis.
%
While it would be \emph{sound} to prove it and replace
it by \ls$True$, it is useless at best,
and usually irreparably affects the system. Consider asserting the tautology
\ls$(False `with_tactic` tau) ==> False$.
% replacing it with a pair of goals for \ls$False$ and \ls$True ==> False$
% renders it unprovable.

Hence, \fstar splits such obligations only in strictly-positive
positions.\ch{what does strictly positive mean for logical formulas? I
  only know positive and negative. Strictly-positive is for
  inductives.}\nik{it just means, as also in the case of inductives,
  not to the left of any arrow/implication}\ch{Should we explain
  why we need the strict?}
%
On all others, \fstar simply drops the \ls$with_tactic$ marker, e.g., by just unfolding the definition of \ls$with_tactic$.
%
For regular uses of the \ls$assert..by$ construct, however, all occurrences are
strictly-positive.
%
It is only when (expert) users use the \ls$with_tactic$ marker directly
that the above discussion might become relevant.

 
% \begin{figure}
% % \begin{wrapfigure}{l}{0.45\textwidth}
% % \vspace{-12pt}
% \begin{centering}
% \begin{lstlisting}
%                          a      : eqtype
%                          f      : total_order a
%                          l      : list a
%                          h0     : squash (~(Nil? l))
%                          h1     : squash (Cons? l)
%                          pivot  : a
%                          tl     : list a
%                          h2     : squash (l == pivot :: tl)
%                          hi, lo : list a
%                          h3     : squash (partition (f pivot) tl == hi,lo)
%                          -----------------------------------------------------------------------------------------------------------------------------------------------
%                          ?u : squash (perm l (quicksort f lo @ pivot :: quicksort f hi))
% \end{lstlisting}
% \end{centering}
% \caption{Proof state for \ls$qs_tac$ (simplified)}
% \label{fig:proofstate}
% \end{figure}

% For our quicksort example in \autoref{sec:quicksort},
% the proofstate is shown on \autoref{fig:proofstate} (hand-edited for clarity).
% %
% As expected, the context includes binders for the variables in
% the lexical scope.
% %
% But, more interestingly, it also includes hypotheses \ls$h0, h1$,
% representing the result of case analyzing \ls$l$ with the patterns taken
% in order; the hypothesis \ls$h2$ for the destruction of \ls$l$ into its
% components; and \ls$h3$, an equation for the pattern matching on the
% tuple \ls$(hi, lo)$.
% %
% In other words, the proof state summarizes the
% control and data flow assumptions at that program point,
% as expected by programmers using regular assertions in \fstar.
%% This procedure has proved to work well on scale
%%     \guido{and corresponding to what programmers expect?}.

% Building this
% proof state during VC generation in a form that allows us to focus a
% tactic's attention on just a specific part of VC is one of our main
% contributions.

Formally, the soundness of this whole approach is given by the following
metatheorem, which justifies the splitting out of sub-assertions,
and by the correctness of evolution detailed in \autoref{sec:correctness}.
%
The proof of \autoref{thm:split} is straightforward, and included in
\iffull
\autoref{app:correctness}.
\else
the appendix.
\fi
%
We expect an analogous property
to hold in other verifiers as well (in particular, it
holds for first-order logic).

\begin{theorem}
    \label{thm:split}
    Let $E$ be a context with 
    $\Gamma \vdash E : prop \Rightarrow prop$,
    and $\phi$ a squashed proposition such that $\Gamma \vdash \phi : prop$.
    Then the following holds:
    \[
        \infer{\Gamma \vDash E[\phi]}
              {\Gamma \vDash E[\top] & \Gamma, \gamma(E) \vDash \phi}
    \]
    where $\gamma(E)$ is the set of binders $E$ introduces.
    If $E$ is strictly-positive, then the reverse implication holds as well.
\end{theorem}
%


%% solving \ls$squasha (x > 0)$ by applying
%% the lemma \ls$x:pos -> squash (x > 0) -> squash$, p when trying to a
%% %


%% to a prometaprogram solving a

%% %
%% A curious fact must be considered for a metaprogramming system like the
%% one we describe.
%% %
%% We have explained how guards arise as side-conditions
%% from type-checking terms, and which must be later proven.
%% %
%% Since metaprograms are allowed to call type-checker internals
%% somewhat arbitrarily, they can also cause guards to be
%% generated, which must be proven.
%% %
%% The setting is kept in the proofstate and can be modified
%% by metaprograms as needed.


% smt goals, guard_policy

% Once that we have a way to carve out subgoals from automatically-generated VCs,

% \nik{bottom-up evaluation etc.}
% \guido{
% Correctness proof: each action respects a preorder. Relate to
% monotonicity paper, but we do not internalize the proof (we would need
% to internalize typing).

% Detail the preorder: a set of goals s1 into another s2, such that there
% is mapping from all s2-solutions to each s1-solution.

% Exceptions should not bring problems. I was thinking that \textbf{any} of the two
% combinations of state plus exceptions are safe. We discard it just out of convenience
% (to de-unify). Danel pointed me to Update Monads to formalize this and the previous part,
% but I'm still kind of confused.}


%% \paragraph{Initial state}
%% There are several ways for metaprograms to come to life.\ch{What
%%  does marked for execution mean? You mean integrated into an \fstar{} file?}
%% %
%% We have already seen examples of \ls$synth$, which calls the tactic
%% to produce a term during type-checking; \ls$%splice$, which is used to
%% generate top-level definitions; and \ls$assert phi by tau$ for proving
%% assertions.\ch{There is also a way to process VCs for top-level definitions}
%% %
%% The initial goal for \ls$synth$ and \ls$%splice$
%% are clear: a term of the expected type in the current typing environment,
%% and a term of type \ls$def$ for \ls$%splice$.\ch{Was type def introduced
%%   already, even so probably better to not explain in words: definition type?}
%% %
%% However, the environment for \ls$assert phi by tau$ must also
%% contain assumptions for the control flow of the program,
%% otherwise
%% assertions such as:
%% \begin{lstlisting}
%% let abs (i : int) : nat = if i < 0 then (assert (--i > 0) by tau; -- i) else i
%% \end{lstlisting}
%% are simply unprovable.
%% %
%% In the next subsection, we describe how we ensure a proper initial
%% environment is provided for tactics that are used for solving assertions with tactics.


% For instance, during this traversal, when\cpc{Word missing here} encounters the
% sub-goal \ls$p `with_tactic` t$, it unfolds the definition
% of \ls$with_tactic$, turning the goal into \ls$p$. Then, it
% unquotes \ls$t$ into a tactic term of type \ls$unit -> Tac unit$, and
% calls it to work only on the current goal \ls$p$. When \ls$t$ returns,
% the traversal can continue processing the proof state that remains.
% As such, all of this, including the traversal, can be accommodated
% within \metafstar{}'s tactic framework with no additional trust
% requirements. However, since this traversal is very common in
% practice, \metafstar provides primitive support for it out of the box,
% rather than forcing the user to define a tactic to do it.\cpc{As written, this sounds deceptive: we start by saying that it could be user-defined, and then we say that we made it a primitive.  The logical conclusion would be that we put it in the stdlib, not that we made it a primitive, AFAICT.}
% \ch{I find introducing m a super big detour that
% blocks understanding the main ideas. I would suggest to rephrase this
% in terms of a primitive functionality, and if you wish then discuss
% trust and m.}

%% As we've described, VCs are generated automatically by the \fstar{}
%% while type-checking programs and discharged automatically in the
%% background.
%% %
%% For allowing metaprograms to run on these VCs and process them appropriately,
%% we must allow users to specify that they want a certain tactic to process a VC.
%% %
%% However, that is not enough by itself: VCs can quickly become unwieldly
%% when the complexity of programs or of specifications increases.
%% %
%% Consider the following factorial program:
%% %
%% \begin{lstlisting}
%% val fact : nat -> Tot nat
%% let fact n = if n = 0 then 1 else n * fact (n -- 1)
%% \end{lstlisting}
%% %
%% The specification is simply stating that, given a \ls$nat$ (natural number) \ls$n$,
%% \ls$fact n$ returns a \ls$nat$, where \ls$nat$ is defined as the nonnegative
%% refinement of the integers.
%% % .\ch{and natural numbers are
%% % defined as only the non-negative integers!}
%% %
%% Yet the VC for it looks surprisingly complex at a first glance:
%% % \ch{Why surprising?  surprisingly complex maybe?}
%% \begin{lstlisting}
%% $\models$ forall (k:(nat -> prop)). (forall (x:nat). k x <==> True) ==>
%%      n = 0 ==> (forall (y:int). y == 1 ==> k y)
%%    /\ n $\ne$ 0 ==> b2t (n -- 1 >= 0) /\ n -- 1 << n
%%                  /\ (forall (y:nat{y << n}).  y == n -- 1 ==>
%%                         (forall (fy:nat).  fy == fact y ==>
%%                             (forall (r:int).  r == n * fy ==> b2t (r >= 0) /\ (forall (r':nat). r' == r ==> k r))))
%% \end{lstlisting}
%% %
%% There are 3 main subgoals in this VC, which state that:
%% (1) the recursive call to \ls$fact$ is on a natural number argument;
%% (2) \ls$fact$ returns a natural number;
%% (3) every call to \ls$fact$ is made for a smaller number, i.e., that it terminates.
%% %
%% Furthermore, it involves somewhat roundabout (mostly
%% due to the WP calculus from which it arises)
%% and dealing with many details that the programmer should not have to worry
%% about, such as the type of postconditions, \fstar{}'s generic termination
%% metric, and boolean-to-type coercions.

%% This illustrates a key point about the introduction of tactics in
%% \guido{first use of auto-active, probably shouldn't be}
%% auto-active verification systems: the VC has an inpredictable shape,
%% it can be obscure for the layman programmer,
%% and it might include several properties in the same formula.
%% %
%% Furthermore slight changes to the program will impact the VC in possibly
%% nontrivial ways.
%% %
%% Hence, the programmer needs a mean to single-out the properties she
%% wishes to prove via application of a tactic, while being able to safely
%% ignore the rest and leave it to be proven by the usual method.


%% We provide this flexibility via a special marker named \ls$by_tactic$,
%% of type \ls$tactic unit -> prop -> prop$.
%% %
%% This marker is simply defined as projecting its second argument, that
%% is:
%% %
%% \begin{lstlisting}
%% abstract let by_tactic tau phi = phi
%% \end{lstlisting}
%% %
%% The \ls$abstract$ qualifier, however, is key: it prevents the marker
%% from being reduced and disappearing, maintaining the tactic annotations
%% provided by the user.
%% %
%% These annotations then make it into the VCs which, before being sent to
%% the SMT solver, have these marked subformulae first ``split out'' from
%% the main VC.
%% %
%% They can then be preprocessed with their assigned tactics in order to
%% break them down, possibly proving them or generating more subgoals.
%% %
%% The final set of goals is then fed to the SMT solver one at a time.

%% % Logically, however, the marker has no meaning, and we provide a lemma
%% % stating exactly this:
%% % \begin{lstlisting}
%% % val by_tactic_seman : tau:(unit -> Tac unit) -> phi:prop -> Lemma (phi <==> by_tactic tau phi)
%% % \end{lstlisting}



%% An example of these annotations is the following program, which is in fact what
%% the \ls$assert phi by tau$ notation we have already seen desugars to:

%% \guido{Not exactly what we have, no prop, squashed, and no polymorphism}
%% \begin{lstlisting}
%% let assert_by_tactic (phi:prop) (tau:tactic 'a) : Pure unit (requires by_tactic tau phi)
%%                                                  (ensures fun _ -> phi)
%%         = ()
%% \end{lstlisting}
%% %
%% We note that \ls$assert_by_tactic$ is a definition just like any other,
%% and as such it must be proven to meet its specification of ensuring
%% \ls$phi$ when assuming \ls$by_tactic tau phi$.
%% %
%% Since the program is a no-op, this fact is an implication
%% \ls$by_tactic tau phi ==> phi$, which is trivially proven by \fstar{}
%% without using the SMT solver.

%% However, calling \ls$assert_by_tactic phi tau$ generates, at that point,
%% a VC requiring to prove its precondition.
%% %
%% Suppose we have the following call to \ls$assert_by_tactic$
%% (using the sugared notation), marking an equality on integers
%% to be preprocessed with the \ls$canon$ tactic,
%% which canonicalizes arithmetic expressions by distributing and commuting with respect to
%% a syntactic ordering on terms, easing or
%% entirely eliminating the work of the SMT solver:
%% \guido{Maybe we want to point to the cooler canonicalizer.}
%% %
%% \begin{lstlisting}
%% val f1 : unit -> Pure unit (requires psi_1) (ensures fun () -> psi_2)
%% val f2 : int  -> Pure unit (requires  xi_1 n) (ensures fun () ->  True)
%% let h n : unit = f1 (); assert (2 * (n + 3) == 2 * n + 6) by canon; f2 n
%% \end{lstlisting}
%% %
%% The VC for this program is a formula similar to:
%% \begin{lstlisting}
%% forall n. psi_1 /\ (psi_2 ==> [by_tactic canon (2 * (3 + n) == 2 * n + 6)] /\ [2 * (3 + n) == 2 * n + 6 ==> (xi_1 n)])
%% \end{lstlisting}

%% % That much is the user's view, but down in the VC world we need to
%% % slay a different beast. CH: too unscientific
%% %
%% % Implementing this abstraction requires solving the following challenge:
%% % since VCs are generated at the granularity of each top-level
%% % definition\guido{correct?}, they contain much more than a single assertion.
%% % %
%% % Usually, in a program composed of a sequence of let-bindings, the VC is
%% % a chain of the pre- and postconditions of each binding, where any given
%% % precondition appears in a logical position where all of the previous
%% % bindings' postconditions are assumed.

%% \ch{It seems confusing to swap the order of arguments between
%%   by\_tactic and assert\_by\_tactic. Can we swap the arguments in
%%   by\_tactic too? It's anyway not consistently used.}
%% \guido{I wouldn't mind. Should we reflect it in the implementation?}

%% In this VC, the propostion we have asserted appears marked with the
%% tactic in a (logically-)\emph{positive} position, and by itself in a
%% \emph{negative} one.
%% %
%% By positive position, we mean the subexpressions of the formula that are
%% not within the antecedent of any implication, and likewise for negative
%% positions and consequents.
%% %
%% It is clear that subformulae appearing in positive positions are proof
%% obligations, and those in negative positions are assumptions for the
%% proof of the corresponding positive position \guido{make more formal}.
%% %
%% As such, we will only attempt to simplify subformulae which appear both
%% marked with a tactic \emph{and} in a positive position.
%% %
%% For negative positions, we simply remove the marker and splice in the
%% original formula.

%% \guido{this is probably super boring and unneeded}
%% There are also operators of mixed polarity, such as \ls$<==>$.
%% %
%% For them, if one of their operands must have a formula split out,
%% we must unfold the if-and-only-if into a conjunction of implications,
%% which we then treat in the same manner.
%% %
%% This allows to still efficiently use tactics without unfolding
%% this logical operator, something which can have a very noticeable
%% impact on the size of the VC.

%% The first step is splitting the formula out of the rest of the VC.
%% %
%% This process of simply replacing the tactic-annotated formula with
%% \ls$True$ (true), and generating an extra subgoal for \ls$phi$ in an
%% extended \emph{logical environment}, which contains assumptions about
%% the program flow.

%% For our example, the resulting VC then looks like
%% \begin{lstlisting}
%% forall n. psi_1 /\ (psi_2 ==> True /\ [2 * (3 + n) == 2 * n + 6 ==> (xi_1 n)])
%% \end{lstlisting}
%% where the proof obligation for the equality has vanished, and is now in a
%% different goal which is, morally, the following:
%% \begin{lstlisting}
%% forall n. psi_2 ==> [by_tactic canon (2 * (3 + n) == 2 * n + 6)]
%% \end{lstlisting}
%% %
%% This is in fact modelled as a goal of precisely the stated proposition
%% (\ls$2 * (3 + n) == 2 * n + 6$) in an environment extending the original
%% VC one with extra binders for the assumptions, such as \ls$n$
%% and the proof of \ls$psi_2$.


%% It is only at this point that the \ls$canon$ tactic can be run
%% on the goal, obtaining a new set of simplified subgoals
%% to feed to the SMT.
%% %
%% (In this particular case, the \ls$canon$ tactic will prove this equality
%% internally, not needing further subgoals.)


%% % Now that the formula of interest has been isolated, we run the tactic
%% % \ls$canon$ on it to get a list of subgoals, which will be fed to the
%% % solver\ch{SMT? only after canon runs though!}  alongside the first one.

%% The formal backing to this VC transformation
%% is the following \fstar{} metatheorem, which is a form of
%% \guido{propositional + functional, say it?}
%% extensionality:
%% %
%% \begin{theorem}
%%     Let $E$ be a context of type $(\Gamma_1, prop) \rightarrow (\Gamma_2, prop)$
%%     and $\Gamma_1 \vdash \phi : prop$ then
%%     \[
%%         \infer{\Gamma_2 \vDash E[\phi]}
%%               {\Gamma_2 \vDash E[\top] & \Gamma_2, \gamma(E) \vDash \phi}
%%     \]
%%     where $\gamma(E)$ represents the binders introduced by $E$
%%     at the hole's position.
%% \end{theorem}
%% \guido{Prove? Our metatheory is actually lacking a bit for it....}
%% %
%% The $\Gamma_2, \gamma(E)$ environment is precisely the extended context
%% we mentioned before.
%% %
%% Running in this environment is what soundly allows metaprograms to use
%% information coming from the control flow of a program, which can be
%% crucial in proving the needed assertion.\guido{somewhat repeated, could drop}


%% Note that this theorem does not depende at all on the shape
%% of the contexts \ls$E$.


%% \guido{Statement of first theorem from correctness.org. Explain the positivity
%% analysis and the recent special semantics for $\iff$. In a sense, ours is a very
%% practical approach matching on symbols themselves and not their definitions.
%% We are currently crucially relying on them not being unfolded.}

% \guido{Maybe these should not be here?}
% Basic examples
% Pattern matching library
% Arithmetic canonizer
% bvtac


\section{Executing Metaprograms Efficiently}
\label{sec:evaluation}
%\nik{Budget: 2.5 pages}

% This mostly supersedes section 3.3... once this section is fully fleshed out,
% do a pass (TODO) over 3.3 to make sure we haven't lost anything

\fstar provides three complementary mechanisms for running metaprograms.
The first two, \fstar's call-by-name (CBN)
interpreter and a (newly implemented) call-by-value (CBV)
NbE-based evaluator, support strong
reduction---henceforth we refer to these as ``normalizers''.
%% \zp{Couldn't parse the previous sentence,
%% so rephrased a bit; old sentence commented out below}
%The first two mechanisms allow using \fstar's call-by-name (CBN)
%interpreter or its (newly implemented) call-by-value (CBV)
%NbE-based~\cite{bergerNbe} evaluator: both support strong
%
In addition, we design and implement a new \emph{native plugin}
mechanism that allows both normalizers to interface with \metafstar
programs extracted to OCaml, reusing \fstar's existing extraction pipeline for
this purpose.
%
%% Experiments (\autoref{sec:experiments})\guido{check} show that NbE improves on the
%% interpreter by orders of magnitude and native plugins further accelerate
%% NbE.
%
Below we provide
a brief overview of the three mechanisms.

\subsection{CBN and CBV Strong Reductions}

As described in \autoref{sec:tac-effect}, metaprograms, once reified,
are simply \fstar terms of type \ls$proofstate -> Div (result a)$.
%
As such, they can be reduced using \fstar's existing computation
machinery, a CBN interpreter for strong reductions based on the Krivine
abstract machine (KAM)~\cite{cregut2007strongly,Krivine2007}.
%
Although complete and highly configurable, \fstar's KAM interpreter
is slow, designed primarily for converting types during dependent
type-checking and higher-order unification.

Shifting focus to long-running metaprograms, such as
tactics for proofs by reflection,
we implemented an NbE-based strong-reduction evaluator for \fstar
computations. The evaluator is implemented in \fstar and extracted to
OCaml (as is the rest of \fstar{}), thereby inheriting CBV from OCaml.
%
It is similar to~\citepos{BoespflugDG11} NbE-based
strong-reduction for Coq,
although we do not implement their low-level,
OCaml-specific tag-elimination optimizations---nevertheless, 
%
it is already vastly more efficient than the KAM-based interpreter.
%
%% NbE for \fstar is implemented in \fstar and extracted to OCaml, along
%% with the rest of the \fstar compiler.\footnote{This is
%% straightforward inasmuch as \fstar supports ML-style general recursion
%% and non-positive types (safely isolated from its logical core).}
%% %
%% Inheriting the reduction strategy of OCaml, NbE reductions in \fstar
%% use call-by-value.\guido{this para is pretty weird to read, and badly placed}
%% \aseem{Repurposed it in the para above.}
%
%% Unlike Coq, however, since \metafstar unifies programming and
%% metaprogramming,
%
% Both the normalizers  support effectful
% operators for dynamic quotation and unquotation, as well as general
% recursion for potentially divergent metaprograms.
%
%% Since metaprograms are, by design, potentially divergent, both of
%% these reduction machineries support general recursion.
%
% Nevertheless, // GM why "nevertheless?"
%% In what follows, we refer to \fstar's CBN and CBV reducers as
%% ``normalizers''. \aseem{We already called them as normalizers in
%% the first paragraph of the section.}

\ch{Nobody worried that changing between CBN and CBV at effect Div
  does change the semantics?}

\subsection{Native Plugins \& Multi-language Interoperability}

% \guido{Q: multi-language interop is needed only due to open reductions,
% no? metaprograms are closed, even if they handle syntax, so this might
% raise an eyebrow?}

Since \metafstar programs are just \fstar programs, they can also be
extracted to OCaml and natively compiled.
%
Further, they can be dynamically linked into \fstar as ``plugins''.
%
Plugins can be directly called from the type-checker, as is done for the
primitives, which is much more efficient than interpreting them.
% However, unlike regular \fstar programs, extracted metaprograms need to
% be executed during type-checking, and so,
%
%\guido{This ``so'' is not very convincing to me, flow is bad}
% we need to invoke the extracted metaprograms from the type-checker.
%
However, compilation has a cost, and it is not convenient to compile
every single invocation.
%
Instead, \metafstar enables users to choose which metaprograms
are to be plugins
(presumably those expected to be computation-intensive, e.g.
\ls$canon_semiring$).
%
Users can choose their native plugins, while still quickly
scripting their higher-level logic in the interpreter.

This requires (for higher-order metaprograms) a form of multi-language
interoperability, converting between representations of terms used in
the normalizers and in native code.
%
% Concretely, we provide a way to dynamically load compiled metaprograms
% into a running \fstar process, with bindings that allow compiled code to
% be invoked from the normalizers.
%
We designed a small multi-language calculus, with ML-style polymorphism,
to model the interaction between normalizers and plugins and
conversions between terms.
%
\iffull
We outline it in
\autoref{app:native}.
\else
See the appendix for details.
\fi


Beyond the notable efficiency gains of running compiled code vs.
interpreting it, native metaprograms also require fewer embeddings.
% \ch{
%  Isn't ``converting between representations of terms used in the normalizers
%   and in native code'' above also embedding? Which embeddings are you talking
%   about here? The ones from 4.1?}%
%   \guido{Yes, but once compiled there is no barrier
%   between the metaprogram and the primitives, so no more embeddings.
%   Embeddings happen when crossing a boundary. Tried to make it clearer.}%
%
Once compiled, metaprograms work over the internal, \emph{concrete}
types for \ls$proofstate$, \ls$term$, etc., instead of over their
\fstar representations (though still treating them abstractly).
%
Hence, compiled metaprograms can call primitives without needing to
embed their arguments or unembed their results.
%
Further, they can call each other directly as well.
%
Indeed, operationally there is little operational difference
between a primitive and a compiled metaprogram used as a plugin.


Native plugins, however, are not a replacement for the normalizers,
for several reasons.
%
First, the overhead in compilation might not
be justified by the execution speed-up.
%
Second, extraction to OCaml erases types and proofs.
%
As a result, the \fstar \emph{interface} of the native plugins can only
contain types that can also be expressed in OCaml, thereby excluding
full-dependent types---internally, however, they can be dependently typed.
%% We provide a type-directed way to reconstruct a value returned by
%% OCaml as an \fstar term, but this mechanism only supports plugins
%% whose \fstar interface contains types that can also be expressed in
%% OCaml, e.g., the \emph{interface} of a plugin cannot use full
%% dependent types.
%
%\guido{It's now unclear why one needs to reduce open terms}
Third, being OCaml programs, native plugins do not support reducing
open terms, which is often required.
%
However, when the programs treat their open arguments parametrically,
relying on parametric polymorphism, the normalizers can pass such arguments
\emph{as-is}, thereby recovering open reductions in some cases.
This allows us to
use native datastructure implementations (e.g. \ls{List}),
which is much faster than using the normalizers, even for open terms.
%
\iffull
We discuss this briefly in
\autoref{app:native}.
\else
See the appendix for details.
\fi

%For space reasons,
%we delegate the calculus to \autoref{app:native}.

%
% Finally, metaprograms that use dynamic quotations cannot be evaluated
% natively, since it relies on dynamically embedding a value into its
% syntax as a term, which is not possible in a compiled OCaml program,
% e.g., for closures. Likewise, unquotation is not supported, since it
% amounts to reflective code evaluation.

%% between the source normalizer terms and their erased, compiled
%% counterparts.
%% relating our
%% formalization to the actual implementation as we go.
%
%% Reflecting our requirement that plugins are only supported at
%% ML-typeable interfaces, rather than full \fstar, our source language
%% only contains terms at ML types (i.e,. rank 1 polymorphism over simple
%% types).
%% %
%% Conversely, the target language is the untyped lamdba calculus, reflecting
%% the type-less native representation of compiled OCaml code.
%


% BNF symbols 
\newcommand{\bnfalt}{{\bf \,\,\mid\,\,}}
\newcommand{\bnfdef}{{\bf ::=~}}

\newcommand\eqdef{\mathrel{\overset{\makebox[0pt]{\mbox{\normalfont\tiny\sffamily def}}}{=}}}

\newcommand{\lk}[1]{%
  \texttt{%
    \fontdimen2\font=2pt   % Shrink default inter-word space
    #1}}

%% COLORS
\definecolor{dkblue}{rgb}{0,0.1,0.5}
\definecolor{dkgreen}{rgb}{0,0.5,0}
\definecolor{dkred}{rgb}{0.7,0,0}
\definecolor{dkpurple}{rgb}{0.7,0,0.4}

\newcommand{\highlight}[1]{%
  \colorbox{dkpurple}{$\displaystyle#1$}}

\newcommand{\blue}[1]{{\textcolor{dkblue}{#1}}}
\newcommand{\red}[1]{{\textcolor{dkpurple}{#1}}}


%% SOURCE
\newcommand{\source}[1]{\blue{#1}}
\newcommand{\sourcelk}[1]{\blue{\lk{#1}}}
\newcommand{\sexp}{\source{e}}
\newcommand{\sexpp}{\source{e'}}
\newcommand{\sexpi}{\source{e_1}}
\newcommand{\sexpii}{\source{e_2}}
\newcommand{\sexpip}{\source{e_1'}}
\newcommand{\sexpiip}{\source{e_2'}}

\newcommand{\sval}{\source{v}}
\newcommand{\svali}{\source{v_1}}
\newcommand{\svalii}{\source{v_2}}
\newcommand{\sctxeq}{\cong_{ctx}}
%%% TYPES
\newcommand{\stau}{\source{\tau}}
\newcommand{\ssig}{\source{\sigma}}
\newcommand{\staui}{\source{\tau_1}}
\newcommand{\stauii}{\source{\tau_2}}
\newcommand{\staup}{\source{\tau'}}
\newcommand{\sdelta}{\source{\delta}}
\newcommand{\sdeltai}{\source{\delta_1}}
\newcommand{\sdeltaii}{\source{\delta_2}}
\newcommand{\sdeltap}{\source{\delta'}}
\newcommand{\stunit}{\sourcelk{unit}}
\newcommand{\snat}{\sourcelk{nat}}
\newcommand{\sttop}{\source{\top}}
\newcommand{\sarrow}[2]{#1\source{\rightarrow}#2}
\newcommand{\sprod}[2]{#1\source{\times}#2}
\newcommand{\ssum}[2]{#1\source{+}#2}
\newcommand{\sall}[2]{\source{\forall}#1\source{.}#2}


%%% ENV
\newcommand{\senv}{\source{\Gamma}}
\newcommand{\stenv}{\source{\Delta}}
\newcommand{\senvemp}{\source{\cdot}}
\newcommand{\shastyp}[2]{#1 : #2}
\newcommand{\senvext}[3]{#1,\shastyp{#2}{#3}}
\newcommand{\senvtext}[2]{#1,#2}

\newcommand{\tenv}{\target{\Gamma}}
\newcommand{\tenvemp}{\target{\cdot}}
\newcommand{\tenvext}[2]{#1,#2}

%%% SUBST
\newcommand{\sesubst}{\source{\delta}}
\newcommand{\sesubstemp}{\source{\cdot}}
\newcommand{\sesubstsingl}[2]{\source{[}{#1}\source{\mapsto}{#2}\source{]}}
\newcommand{\sesubsttwo}[4]{\source{[}{#1}\source{\mapsto}{#2},{#3}\source{\mapsto}{#4}\source{]}}
\newcommand{\sesubstlist}[1]{\source{[}{#1}\source{]}}
\newcommand{\bind}[2]{{#1}\source{\mapsto}{#2}}
\newcommand{\sesubstext}[3]{#1\source{[}{#2}\source{\mapsto}{#3}\source{]}}

\newcommand{\ssubst}[2]{\source{[}#2\source{/}#1\source{]}}
\newcommand{\tsubst}[2]{\target{[}#2\target{/}#1\target{]}}
\newcommand{\tesubstsingl}[2]{\target{[}{#1}\target{\mapsto}{#2}\target{]}}

%%% EXPRESSIONS
% unit
\newcommand{\sunit}{\sourcelk{()}}
% var
\newcommand{\svar}[1]{\source{#1}}
% abstractions
\newcommand{\sabs}[3]{\source{\lambda} #1 \source{:} #2 \source{.} #3}
\newcommand{\sapp}[2]{#1~#2}
% type abstractions
\newcommand{\stabs}[2]{\source{\Lambda #1\source{.} #2}}
\newcommand{\stapp}[2]{#1~#2}
% pairs
\newcommand{\spair}[2]{\source{(}#1 \source{,} #2\source{)}}
\newcommand{\sproj}[2]{#1\source{.}#2}
\newcommand{\sproji}[1]{#1\source{.i}}
\newcommand{\sfst}[1]{\sproj{#1}{\source{1}}}
\newcommand{\ssnd}[1]{\sproj{#1}{\source{2}}}
% sums
\newcommand{\sinl}[1]{\sourcelk{inl}~#1}
\newcommand{\sinr}[1]{\sourcelk{inr}~#1}
\newcommand{\scasex}[5]{\sourcelk{case}\source{(} #1  \source{,x:#2.} #3  \source{,x:#4.} #5 \source{)}}
\newcommand{\scasey}[5]{\sourcelk{case}\source{(} #1  \source{,y:#2.} #3  \source{,y:#4.} #5 \source{)}}
\newcommand{\scase}[5]{\sourcelk{case}\source{(} #1  \source{,} #2 \source{.} #3  \source{,} #4 \source{.} #5 \source{)}}
% alien
\newcommand{\salien}[3]{\source{\{}#1 \source{\}}_{#2}^{#3}}
\newcommand{\sopaque}[1]{\source{\langle}#1 \source{\rangle}}
% parentheses
\newcommand{\spar}[1]{\source{(}#1\source{)}} 

%% TARGET
\newcommand{\target}[1]{\red{#1}}
\newcommand{\targetlk}[1]{\red{\lk{#1}}}

%%% EXPRESSIONS
\newcommand{\texp}{\target{e}}
\newcommand{\texpp}{\target{e'}}
\newcommand{\tval}{\target{v}}
\newcommand{\tvali}{\target{v_1}}
\newcommand{\tvalii}{\target{v_2}}
\newcommand{\texpi}{\target{e_1}}
\newcommand{\texpii}{\target{e_2}}
\newcommand{\texpip}{\target{e_1'}}
\newcommand{\texpiip}{\target{e_2'}}

% unit
\newcommand{\tunit}{\targetlk{()}}
% var
\newcommand{\tvar}[1]{\target{#1}}
% abstractions
\newcommand{\tabs}[2]{\target{\lambda} #1 \target{.} #2}
\newcommand{\tapp}[2]{#1~#2}
% pairs
\newcommand{\tpair}[2]{\target{(}#1 \target{,} #2\target{)}}
\newcommand{\tproj}[2]{#1\target{.}#2}
\newcommand{\tproji}[1]{#1\target{.i}}
\newcommand{\tfst}[1]{\sproj{#1}{\target{1}}}
\newcommand{\tsnd}[1]{\sproj{#1}{\target{2}}}
% sums
\newcommand{\tinl}[1]{\targetlk{inl}~#1}
\newcommand{\tinr}[1]{\targetlk{inr}~#1}
\newcommand{\tcasex}[3]{\targetlk{case}\target{(}#1  \target{,x.} #2  \target{,x.} #3 \target{)}}
\newcommand{\tcasey}[3]{\targetlk{case}\target{(}#1  \target{,y.} #2  \target{,y.} #3 \target{)}}
\newcommand{\tcase}[6]{\targetlk{case}\target{(}  \target{,} #2 \target{.} #3  \target{,} #4 \target{.} #5 \target{)}}
% alien
\newcommand{\talien}[3]{\target{[}#1 \target{]}_{#2}^{#3}}
\newcommand{\topaque}[1]{\target{\langle}#1 \target{\rangle}}
% parentheses
\newcommand{\tpar}[1]{\target{(}#1\target{)}} 

% TYPING
\newcommand{\iswf}[2]{#1 \vdash #2 ~\lk{wf}}
\newcommand{\hastyp}[4]{#1;~#2 \vdash #3 : #4}
\newcommand{\isok}[3]{#1;~#2 \vdash #3 ~\lk{ok}}

%% COERSIONS
\newcommand{\kap}[4]{\mathcal{K}_{#1}^{#2}(#3,#4)}
\newcommand{\mfun}[2]{\lk{fun}~#1\Rightarrow#2}
\newcommand{\tcoerce}[3]{\mathcal{\target{C}}_{#1}^{#2}(#3)}
\newcommand{\scoerce}[3]{\mathcal{\source{C}}_{#1}^{#2}(#3)}
\newcommand{\prot}[3]{\mathcal{P}_{#2}^{#3}(#1)}

%% SEMANTICS
\newcommand{\seval}{\source{~\rightarrow~}}
\newcommand{\sevals}{\source{~\rightarrow^{*}~}}
\newcommand{\sevalm}[1]{\source{~\rightarrow}^{#1}~}
\newcommand{\teval}{\target{~\rightarrow~}}
\newcommand{\tevals}{\target{~\rightarrow^{*}~}}
\newcommand{\tevalm}[1]{\target{~\rightarrow}^{#1}~}
\newcommand{\sevalr}[1]{\source{~\xrightarrow{\textsc{#1}}~}}
\newcommand{\tevalr}[1]{\target{~\xrightarrow{\textsc{#1}}~}}
\newcommand{\sevalrs}[1]{\source{~\xrightarrow{\textsc{#1}}^{*}~}}
\newcommand{\tevalrs}[1]{\target{~\xrightarrow{\textsc{#1}}^{*}~}}
\newcommand{\isval}[1]{\lk{is\_val}~#1}
% CONTEXT
\newcommand{\scon}[1]{\source{E[}#1\source{]}}
\newcommand{\tcon}[1]{\target{E[}#1\target{]}}

% TRANSLATION
\newcommand{\trans}{~\source{\rightsquigarrow}~}

% FREQUENTLY USED TERMS
%% SOURCE
\newcommand{\sdup}{\stabs{\svar{B}}{\sabs{\svar{x}}{\svar{B}}{\spair{\svar{x}}{\svar{x}}}}}
\newcommand{\sduptyp}{\sall{\svar{B}}{\sarrow{\svar{B}}{\sprod{\svar{B}}{\svar{B}}}}}
\newcommand{\sapptyp}{\sall{\svar{A}}{\sarrow{\spar{\sduptyp}}{\sarrow{\svar{A}}{\sprod{\svar{A}}{\svar{A}}}}}}
\newcommand{\sapptypfree}{\sarrow{\spar{\sduptyp}}{\sarrow{\svar{A}}{\sprod{\svar{A}}{\svar{A}}}}}

\newcommand{\sarith}{\source{(1 + 2 + \dots)}}
\newcommand{\sinctyp}{\sarrow{\snat}{\snat}}

\newcommand{\sapptypalt}
{\sall{\svar{A}}
  {\sall{\svar{B}}
    {\sarrow{\spar{\sarrow{\svar{A}}{\svar{B}}}}{\sarrow{\svar{A}}{\svar{B}}}}}}
\newcommand{\sapptypaltf}
{\sall{\svar{B}}
     {\sarrow{\spar{\sarrow{\svar{A}}{\svar{B}}}}{\sarrow{\svar{A}}{\svar{B}}}}}
\newcommand{\sapptypaltff}
   {\sarrow{\spar{\sarrow{\svar{A}}{\svar{B}}}}{\sarrow{\svar{A}}{\svar{B}}}}

\newcommand{\sidf}[1]{\sabs{\svar{y}}{#1}{\svar{y}}}

\newcommand{\sappid}
{ \stabs
    {\svar{A}}
    {\sabs{\svar{x}}{\svar{A}}
    {\sabs{\svar{f}}
      {\sarrow{\svar{A}}{\svar{A}}}
      {{\sapp{\svar{f}}{\spar{\sapp{\spar{\sidf{\svar{A}}}}{\svar{x}}}}}}
    }}
}
\newcommand{\sappidf}
{\sabs{\svar{x}}{\svar{A}}
  {\sabs{\svar{f}}
      {\sarrow{\svar{A}}{\svar{A}}}
      {{\sapp{\svar{f}}{\spar{\sapp{\spar{\sidf{\svar{A}}}}{\svar{x}}}}}}}}
\newcommand{\sappidarg}[1]
{\sabs{\svar{f}}
      {\sarrow{\svar{A}}{\svar{A}}}
      {{\sapp{\svar{f}}{\spar{\sapp{\spar{\sidf{\svar{A}}}}{#1}}}}}}
\newcommand{\sappidargs}[2]
{\sapp{#2}{\spar{\sapp{\spar{\sidf{\svar{A}}}}{#1}}}}


\newcommand{\sappidtyp}
  {\sall{\svar{A}}
      {\sarrow{\svar{A}}{\sarrow{\spar{\sarrow{\svar{A}}{\svar{A}}}}{\svar{A}}}}}
\newcommand{\sappidtypf}
  {\sarrow{\svar{A}}{\sarrow{\spar{\sarrow{\svar{A}}{\svar{A}}}}{\svar{A}}}}
\newcommand{\sappidtypff}
  {\sarrow{\spar{\sarrow{\svar{A}}{\svar{A}}}}{\svar{A}}}

%% TARGET
\newcommand{\tfapp}{\tabs{\tvar{f}}{\tabs{\tvar{x}}{\sapp{\tvar{f}}{\tvar{x}}}}}
\newcommand{\tfapparg}[1]{\tabs{\tvar{x}}{\sapp{#1}{\tvar{x}}}}
\newcommand{\tfappargs}[2]{\sapp{#1}{#2}}
\newcommand{\tinc}{\tabs{\tvar{x}}{\target{x+1}}}
\newcommand{\tid}{\tabs{\tvar{x}}{\tvar{x}}}


%% Logical relation
\newcommand{\rexp}[3]{#2 \source{\lesssim}_{#1} #3}
\newcommand{\rval}[3]{#2 \source{\leq}_{#1} #3}
\newcommand{\renv}[2]{\source{\mathcal{G}\llbracket} #1 \source{\rrbracket}_{#2}}
\newcommand{\inrenv}[4]{(#1,#2) \in \renv{#3}{#4}}
\newcommand{\rtypenv}[1]{\source{\mathcal{D}\llbracket} #1 \source{\rrbracket}}
\newcommand{\inrtypenv}[2]{#1 \in \rtypenv{#2}}

\newcommand{\rtexp}[3]{#2 \target{\lesssim}_{#1} #3}
\newcommand{\rtval}[3]{#2 \target{\leq}_{#1} #3}
\newcommand{\rtenv}[2]{\target{\mathcal{G}{\llbracket}} #1 \target{\rrbracket}_{#2}}
\newcommand{\inrtenv}[4]{(#1,#2) \in \rtenv{#3}{#4}}


\newcommand{\setst}[3]{ \{ (#1, #2) ~|~ #3 \}}
\newcommand{\setsingl}[1]{ \{\}}
\newcommand{\setsto}[2]{ \{ #1 ~|~ #2 \}}

\newcommand\defeq{\mathrel{\overset{\makebox[0pt]{\mbox{\normalfont\tiny\sffamily def}}}{=}}}

\newcommand{\sforce}[1]{\source{\mathit{force}}(#1)}
\newcommand{\tforce}[1]{\target{\mathit{force}(#1)}}


\section{Experimental evaluation}
\label{sec:experiments}

We now present an experimental evaluation of \metafstar.
% from various angles. -- only 2
%
% First, we evaluate raw performance of our native plugins by proving
% simple, yet intensive, program equivalence in a small assembly-like
% language.
%
First, we provide benchmarks comparing our reflective canonicalizer
from \autoref{sec:canon_semiring} to calling the SMT solver directly
without any canonicalization.
%
% \ch{All this is gone, right?
% ``We then provide more detail on our encoding of separation logic and
% its companion tactic, along with a pattern matching tactic
% combinator.''}
%
Then, we return to the parsers and serializers from
\autoref{sec:minilowparse} and show
how, for VCs that arise, a domain-specific tactic is much more
tractable than a SMT-only proof.
%
% Finally, we comment on two other interesting \metafstar case studies:
% a pattern matcher and a tactic for separation logic that, following
% our hybrid approach, limits metaprogramming to inferring heap frames,
% leaving the rest to the SMT solver.





% focusing primarily on a quantitative study of its three execution
% strategies. Our high-level conclusion is that, as expected, NbE is a
% significant improvement over \fstar's existing interpreter. Further,
% native plugins significantly boost the performance of both \fstar
% normalizers. However, as explained in \autoref{sec:evaluation}, native
% plugins cannot handle reduction of open terms or of metaprograms that
% use dynamic quotation: to benefit from native plugins in such
% scenarios requires leveraging parametricity through a little
% ingenuity.

\iffalse
\subsection{Program Equivalence by Partial Evaluation and SMT}
\label{sec:registers}

Our first quantitative case study involves reasoning about programs in
a small, imperative deeply embedded language in \fstar.
%
The language includes reads and writes to a fixed number of mutable
global integer-typed variables (aka registers), sequencing,
conditionals, and basic arithmetic operations over integers.
%
The listing below shows its syntax: the type of register names is
is \ls$reg_t$ and \ls$regmap$ is a map, whose representation we vary as below,
storing the values of all the
registers.
%
%% Our experiments vary the choice of representation of \ls$regmap$, as
%% discussed shortly.

\begin{lstlisting}
type reg_t  = x:nat{x<=10}
type regmap = R.map int
type inst = | Add : reg_t -> reg_t -> reg_t -> inst $\ldots$ | Const : int -> reg_t -> inst
          | If0 : reg_t -> list inst -> list inst -> inst | Seq : list inst -> inst
\end{lstlisting}

\noindent The semantics of \ls{inst} is given
by a straightforward recursive function,
%
\ls$eval : inst -> regmap -> regmap$.
%%  whose implementation is
%% entirely straightforward.

Our goal is to analyze programs in this language for equivalence:
%% according to the following notion of equivalence:
instruction sequences \ls$p1$ and \ls$p2$ are equivalent if evaluating them in
any initial state \ls$rm$ produces equivalent final states.

\begin{lstlisting}
let equiv p1 p2 = forall (rm:regmap) (r:reg_t). R.sel (eval (Seq p1) rm) r == R.sel (eval (Seq p2) rm) r
\end{lstlisting}

Our strategy for determining equivalence of a given pair of programs
is to partially evaluate them in a symbolic initial state, producing
an open term denoting their final states and aiming to decide
equivalence of these final states using SMT.
%
Without partial evaluation, feeding equivalence checks directly to an
SMT solver does not scale beyond a few instructions, since
this would require the SMT solver to itself compute with the
definition of \ls$eval$. \fstar's encoding to SMT makes it possible
for an SMT solver to compute with recursive definitions, but this is
not particularly efficient beyond unrolling just a few recursive calls
(usually 2-4, but sometimes around 20).
%
By delegating partial evaluation to \fstar's normalizers, we leave to
SMT just the problem of deciding equivalence of arithmetic expressions.
%
Although very simple, this strategy is intended as a standalone model
of a more substantial case study involving the use of partial
evaluation to generating efficient SMT-friendly verification
conditions for assembly languages deeply embedded
in \fstar.\footnote{\url{https://github.com/project-everest/vale}}

\paragraph{Choices of representation of register maps}
Partially evaluating \ls$eval$ over an arbitrary initial state \ls$rm$
involves reducing open terms. Yet, to make register lookups efficient,
we would like to compile \ls$regmap$ operations to native code and use
them as a plugin. Our key observation is that operations to select and
update registers are polymorphic in the type of register values
$\alpha$, as shown below:

\begin{lstlisting}
val R.const: 'a -> regmap 'a
val R.sel : regmap 'a -> reg_t -> 'a
val R.upd : regmap 'a -> reg_t -> 'a -> regmap 'a
\end{lstlisting}

This opens the door to computing natively on symbolic register maps,
since, as explained in \autoref{sec:parametricity}, we can interface with
native plugins at polymorphic types simply by passing, possibly
open, arguments as-is.
%
In the experiments reported here, we consider two representations
for \ls$regmap 'a$; (1) a list of pairs together with a default
value \ls@list (int * 'a) * 'a@; and (2) a function \ls$int -> 'a$.

To make this work we introduce a revised, normalization-friendly
notion (\ls$equiv_norm$) equivalence easily proven equivalent
to \ls$equiv$. 
%
\begin{lstlisting}
let all_equiv (rm1 rm2:regmap) =
    let rec aux (r:reg_t) = R.sel rm1 r == R.sel rm2 r /\ (if r = 0 then True else aux (r $-$ 1)) in
    aux 10
let equiv_norm p1 p2 = forall rm. let rm = R.eta_map 10 rm in all_equiv (eval (Seq p1) rm) (eval (Seq p2) rm))
\end{lstlisting}
%
There are two differences. First, \ls$all_equiv$ reduces the universal
quantification over registers to a finite conjunction. More
technically, to enable reduction of \ls$eval$ on a symbolic machine
state, we use an eta-expanded form of \ls$rm$. For the list
representation, this turns \ls$rm$ into an open term of the form
\ls$[(0, R.sel' rm 0), ..., (10, R.sel' rm 10)], snd rm$, where
the skeleton of the list structure of \ls$rm$ is unfolded and the
values in the range are uncompiled open terms \ls$R.sel' rm 10$. This
expanded form unblocks reductions that pattern match on \ls$rm$. For
function representation, \ls{eta_map} is an identity.
\aseem{What is \ls{sel'}?
How is it different from \ls{sel}?}
%
Putting the pieces together, to decide the equivalence of two
programs \ls$p1, p2$, one can write:
%
\ls$assert (equiv_norm p1 p2) by (norm [delta; zeta; iota]; smt())$

The table alongside shows the time taken (in seconds) by the
normalization step above evaluated on programs \ls$p1, p2$ that
contain on the order of $10 \times 2^k$ instructions, for $k \in \{6,
7, 8, 9\}$.
%
\setlength\intextsep{-0.5pt}
\begin{wrapfigure}{r}{0.67\textwidth}
\begin{small}
\begin{tabular}{|l|l|l|l|l|}
\hline
\# k            &  6            & 7           &       8         & 9 \\\hline
%% \hline\multicolumn{5}{|l|}{Functions} \\\hline
functions               & 0.204 / 0.092         & 0.413 / 0.177       & 0.820 / 0.341           & 1.579  / 0.713 \\
%% \hline\multicolumn{5}{|l|}{Functions, native plugins} \\\hline
%%                &   0.092  &  0.177  &  0.341  &  0.713 \\
%% \hline\multicolumn{5}{|l|}{Lists} \\\hline
lists             &   0.504 / 0.097  &  1.017 / 0.191  &  2.028 / 0.360  &  3.976 / 0.724 \\\hline
%% \hline\multicolumn{5}{|l|}{Lists, native plugins} \\\hline
%%              &   0.097  &  0.192  &  0.360  &  0.724 \\
%% \hline
\end{tabular}
\end{small}
\end{wrapfigure}
%
We report the performance of only our NbE based normalizer in four
configurations---the interpreter is far slower, e.g., failing to
complete the benchmark on lists after several minutes for $k>6$.
%
The first row shows the performance using NbE on functional
register maps; the second on maps represented as lists.
%
The two numbers in each column show the execution time (in seconds) for default NbE
and NbE accelerated with native plugins for registers.
%
In all configurations, we observed a consistent gain in performance of
around a factor of two.


\ifdraft
\begin{verbatim}
2. Canon.Monoid/Canon.Semiring etc.
   a. Intepretation

   b. Native compilation, possible, but at the cost of refactoring
      Guido / Victor working on this?      

   c. Nbe
      Zoe / Guido working on this
      GM: As of now, we don't have NbE tactics running well

3. Parsers/Printers: Tahina
   a. Monolithic syntax building
   b. Incremental, custom type-checking

4. Brief coverage of other examples

   Not sure we really want/have the time to benchmark all of these
   too.

     * sl and slauto
     * transforming monads to build effectful printers
     * bvtac
     * ... 
     ?
\end{verbatim}
\fi

\nik{If we do all the above, this might consume the entire section.}
\fi

\subsection{A Reflective Tactic for Partial Canonicalization}
\label{sec:canonicalization}

% Motivated primarily by the poor support of non-linear arithmetic in SMT
% solvers, and especially when combined with other theories, we set out to develop
% canonicalization tactics that remove the need for non-linear reasoning
% by the SMT solver.
%
%

\newcommand{\canonfig}{
  \begin{wrapfigure}{r}{0.51\textwidth}
      \small
  % \begin{tabular}{|ccccc|}
  % \hline
  %             & Suc. Rate & Total                   & SMT                   & Tactic \\
  % \hline
  %  smt1x      & 5\%       & 2.55   $\pm$      0     & 2.045   $\pm$   0     & 0                   \\
  %  smt2x      & 38\%      & 4.484  $\pm$      1.495 & 4.007   $\pm$   1.492 & 0                   \\
  %  smt3x      & 33\%      & 3.908  $\pm$      1.725 & 3.448   $\pm$   1.714 & 0                   \\
  %  interp     & 100\%     & 3.869  $\pm$      0.494 & 1.736   $\pm$   0.288 & 1.544    $\pm$ 0.196 \\
  %  native     & 100\%     & 2.52   $\pm$      0.347 & 1.704   $\pm$   0.221 & 0.294    $\pm$ 0.048 \\
  % \hline
  % \end{tabular}
      \vspace{-8mm}
  \begin{tabular}{|c||c|c|c|c|}
  \hline
          & Rate     & Queries              & Tactic & Total \\
  \hline
   smt1x   & 0.5\%   & 0.216  $\pm$ 0.001   & --     & 2.937 \\
   smt2x   & 2\%     & 0.265  $\pm$ 0.003   & --     & 2.958 \\
   smt3x   & 4\%     & 0.304  $\pm$ 0.004   & --     & 3.022 \\
   smt6x   & 10\%    & 0.401  $\pm$ 0.008   & --     & 3.155 \\
   smt12x  & 12.5\%  & 0.596  $\pm$ 0.031   & --     & 3.321 \\
   smt25x  & 16.5\%  & 1.063  $\pm$ 0.079   & --     & 3.790 \\
   smt50x  & 22\%    & 2.319  $\pm$ 0.230   & --     & 5.030 \\
   smt100x & 24\%    & 5.831  $\pm$ 0.776   & --     & 8.550 \\
   interp  & 100\%   & 0.141  $\pm$ 0.001   & 1.156  & 4.003 \\
   native  & 100\%   & 0.139  $\pm$ 0.001   & 0.212  & 3.071 \\
  \hline
  \end{tabular}
      \vspace{-6mm}
  \end{wrapfigure}
}

In \autoref{sec:canon_semiring}, we have described the
\ls$canon_semiring$ tactic that rewrites semiring expressions into sums
of products.
%
% It follows the well-known style of \emph{proofs by
% reflection}~\cite{GregoireM05, gonthier2008formal,ChaiebN08,Besson06},
% reducing the amount of proving that needs to be performed to a few key
% points, and leaving bulk of the work to mechanical computation.
%
% Concretely, we first develop a tactic for canonicalizing expressions in
% commutative monoids \ls$canon_monoid$, and then ``stack'' it on top of
% itself to work over commutative semirings.
%
% Differently from other existing tactics, it does not fully canonicalize
% expressions. Instead, it turns them into a shape that is close
% enough to their normal form for the SMT solver to be able to conclude
% their equality.
%
% In general, most of our tactics prove equalities or formulas ``modulo''
% the well-behaved SMT fragment, reducing the metaprogramming task to
% only the key steps.
%
% The extra metaprograming work to fully-canonicalize an expression
% is modest, but in other cases it can make a substantial difference
% (c.f. \autoref{sec:seplogic}).
%
We find that this tactic significantly improves proof robustness.
%
The table below compares the success rates and times for the
\ls$poly_multiply$ lemma from \autoref{sec:poly_multiply}.
%
To test the robustness of each alternative, we run the tests 200 times
while varying the SMT solver's random seed.
%
The \texttt{smt$i$x} rows represent asking the solver to prove the lemma
without any help from tactics, where $i$ represents the resource limit
(\texttt{rlimit}) multiplier given to the solver.
%
This \texttt{rlimit} is memory-allocation based and independent of the
particular system or current load.
%
For the \ls$interp$ and \ls$native$ rows, the
\ls$canon_semiring$ tactic is used, running it using \fstar's KAM normalizer and
as a native plugin respectively---both with an \texttt{rlimit}
of $1$.
%
For each setup, we display the success rate of verification, the average
(CPU) time taken for the SMT queries (not counting the time for parsing/processing
the theory)
with its standard deviation, and
the average total time (its standard deviation coincides with that of
the queries).
%
When applicable, the time for tactic execution (which is independent of
the seed) is displayed.
%
The \texttt{smt} rows show very poor success rates: even when upping the
\texttt{rlimit} to a whopping 100x, over three quarters of the attempts
fail.
%
\canonfig
%
Note how the (relative) standard deviation increases with the
\texttt{rlimit}: this is due to successful runs taking rather random
times, and failing ones exhausting their resources in similar times.
%
The setups using the tactic show a clear increase in
robustness: canonicalizing the assertion causes this proof to always
succeed, even at the default \texttt{rlimit}.
%
% \guido{Here, say something about that single data point. are tactic
% timings better on success? on failure? always?}
%
We recall that the
tactic variants still leave goals for SMT solving, namely, 
the skeleton for the original VC
and the canonicalized equality left by the tactic, easily dischargeable
by the SMT solver through much more well-behaved linear reasoning.
%
The last column shows that native compilation speeds up this tactic's
execution by about 5x.


\subsection{Combining SMT and Tactics for the Parser Generator}
\label{sec:eval-minilowparse}

% Tactics: (i: parser impl; ': additional measure)
%% 11760(4) 26738(4i)
%% 11858(4') 25897(4'i) 
%% 34767(7) 79286(7i)
%% 40234(7') 76638(7'i)
%% 74663(10) 148672(10i)
%% 81012(10') 134515(10'i)

In \autoref{sec:minilowparse}, we presented a library of combinators
and a metaprogramming approach to automate the construction of
verified, mutually inverse, low-level parsers and serializers from type
descriptions.
%
% Metaprogramming helps inasmuch as it relieves the programmer from
% directly programming against our combinator library.
%
Beyond generating the code, tactics are used to process and discharge
proof obligations that arise when using the combinators.

We present three strategies for discharging these obligations,
including those of bijectivity that arise when constructing parsers
and serializers for enumerated types. First, we used \fstar's default
strategy to present all of these proofs directly to the SMT
solver. Second, we programmed a $\sim$100 line tactic to discharge
these proofs without relying on the SMT solver at all. Finally, we
used a hybrid approach where a simple, 5-line tactic is used to prune
the context of the proof removing redundant facts before presenting
the resulting goals to the SMT solver.

%% The parser generator library we presented
%% in \autoref{sec:minilowparse} promotes a style in which a user relies
%% on SMT and dependent types to craft and prove correct a high- and
%% low-level libraries of combinators for parsing and serialization.  The
%% library internalizes proofs of correctness, mutual inversion, and
%% memory safety. However, rather than use the library directly, we
%% showed how to generate high-lve

%% high-level parsers/serializers can be
%% automatically transformed to equivalent low-level ones; in contrast,
%% constructing the low-level code directly, even while using our
%% combinator library, would require a significantly larger effort:
%% first, as shown earlier, simply assembling combinators for even simple
%% types is tedious; second, for some types, like enumerations (or, more
%% generally, tagged unions), the usage of the combinators imposes
%% additional proof obligations (e.g., bijectivity).


%% The only proofs\ch{why are these bijectivity proofs more proofs
%% than the ``other mundane aspects'' below? aren't those things proved as well?
%%   are those things proved generically in the library? (it's not clear below)}
%% remaining on the low-level code show
%% bijectivity\guido{!!!} of the pure functions that describe the parsing/serialization
%% formats---other important but mundane aspects like memory safety
%% are handled within the library and the
%% transformations implemented by our metaprograms. 
%% The parser generator library presented
%% in \autoref{sec:minilowparse} promotes a style in which one relies
%% on SMT and dependent types to craft and prove correct high-level
%% combinators for parsing and serialization. Then, using
%% metaprograms, high-level parsers/serializers can be automatically
%% transformed to equivalent low-level ones, while verifying the low-level code
%% directly would require a significantly larger effort.


%% The only proofs\ch{why are these bijectivity proofs more proofs
%% than the ``other mundane aspects'' below? aren't those things proved as well?
%%   are those things proved generically in the library? (it's not clear below)}
%% remaining on the low-level code show
%% bijectivity\guido{!!!} of the pure functions that describe the parsing/serialization
%% formats---other important but mundane aspects like memory safety
%% are handled within the library and the
%% transformations implemented by our metaprograms. Here, we report on a
%% domain-specific tactic for discharging the arising bijectivity
%% obligations, and compare it to \fstar's default procedure for discharging
%% them, including simplification stages, and ultimately SMT solving.


\begin{wrapfigure}{r}{0.42\textwidth}
    \small
    \vspace{-7mm}
    \begin{tabular}{|c||c|c|c|}
        \hline
                       Size           & SMT only & Tactic only & Hybrid \\
        \hline
                         4            & 178      & 17.3        & 6.6    \\
                         7            & 468      & 38.3        & 9.8    \\
                        10            & 690      & 63.0        & 19.4   \\ \hline
    \end{tabular}
    \vspace{-8mm}
\end{wrapfigure}

The table alongside shows the total time in seconds for verifying
metaprogrammed low-level parsers and serializers for enumerations of
different sizes.
%
In short, the hybrid approach scales the best; the
tactic-only approach is somewhat slower; while the SMT-only approach
scales poorly and is an order of magnitude slower. Our hybrid approach
is very simple. With some more work, a more sophisticated hybrid
strategy could be more performant still, relying on tactic-based
normalization proofs for fragments of the VC best handled
computationally (where the SMT solver spends most of its time), while
using SMT only for integer arithmetic, congruence closure
etc. However, with \metafstar's ability to manipulate proof contexts
programmatically, our simple context-pruning tactic provides a big
payoff at a small cost.

\iffalse
subsection{Other}
\label{sec:other}

We now go over some other interesting metaprograms we have developed,
to convey an idea of the variety of areas where \metafstar finds
applications.

\guido{should be explicit that metaprograms can inspect environments}
\paragraph{Concise pattern matching}
%
We have briefly touched on the basic mechanism by which \metafstar
exposes terms to metaprograms (\autoref{sec:quoting})\guido{fix}.
%
However, this mechanism is often somewhat inconvenient for metaprogrammers,
e.g., when implementing the logic for searching for a specific pattern in the context
or a goal.
%
While entirely possible, the process is exhausting as each pattern of
interest must be written as matches over the \ls$term_view$.
%
Via a careful use of quotation and reflection, we have implemented
a library offering a much more effective and intuitive alternative.
%
The snippet below contains the implementation of a simple tautology prover
built exploiting this library:
\begin{lstlisting}
let rec tauto () : Tac unit = repeat #unit (fun () -> first_match
    with_pat (fun (g: goal (squash True)) -> trivial ());
    with_pat (fun a (h: hyp a) (g: goal a) -> exact (as_term h));
    with_pat (fun a b (g: goal (squash (a /\ b))) -> split ());
    with_pat (fun a b (g: goal (squash (a ==> b))) -> mintro ());
    with_pat (fun a b (g: goal (squash (a \/ b))) -> (fun () -> left (); tauto ()) `or_else` (fun () -> right (); tauto ()))]); 
   qed ()
\end{lstlisting}
\guido{inspecting runtime syntax, raise eyebrow? can we do this pattern matching
with the saner quote?}
The \ls$with_pat$ metaprogram inspects the arguments provided
by the user and interprets them as patterns.
%
Essentially, unlabeled binders are variables, \ls$goal p$
is a pattern that should match \ls$p$ over the goal type, and likewise
\ls$hyp p$ searchs for \ls$p$ in the local context.
%
The \ls$with_pat$ metaprogram also takes a continuation 
argument---this will be run when a match is found, receiving the 
instantiations of the variables.
%
The patterns are in fact \emph{compiled} by \ls$with_pat$ to speed-up
posterior queries.
%
This metaprogram is a fairly complete
implementation of the type of backtracking pattern-matching primitive in
Coq's Ltac~\cite{Delahaye00} or in Lean~\cite{EbnerURAM17},
but built instead as a library.

%(*TODO: allow `e1;e2` even when `e1` doesn't return unit; implement \ls$mintro$; rename \ls$with_pat, first_match$ etc.)
% \cpc{here's a suggestion:\guido{I liked it}}

%
\iffalse
The \ls$with_pat$ metaprogram, of type \ls$#a:Type -> a -> Tac unit$
is has two main
components: a DSL to express pattern-matching problems, and an
interpreter to produce an assignment matching it.
%
Take \ls$with_pat (fun a b (h: hyp (a ==> b)) (_: goal (squash b)) -> mapply h)$
as a concrete example.
%
Using dynamic quoting to inspect its argument, \ls$with_pat$ obtains an
\fstar AST and inspects it
in search of the \ls$goal$ and \ls$hyp$ tags
which indicate what each pattern applies to.
%
From this AST, a pattern-matching problem is built, containing
variables to be instantiated (here \ls$a$ and \ls$b$),
hypothesis patterns (here, \ls$PApp (PQn "==>") [PVar a; PVar b]$), goal
patterns (here, \ls$PApp (PQn "squash") [PVar b]$), and a continuation
to run when a match is found (here, \ls$fun a b h _ => mapply h$).
%
A recursive backtracking unification procedure matches the pattern
against the current goal and hypotheses, producing a lazy stream of
satisfying assignments and running the continuation with each of them,
until the continuation succeeds.
\fi


% The metaprogram \ls$with_pat : #a:Type -> a -> Tac unit$ takes an
% argument (of any type), dynamically quotes it, and then compiles its
% syntax into (1) patterns for the goal and hypotheses of the context,
% and (2) a tactic \ls$k$ of type \ls$matching_solution -> Tac t$ to be
% run on the pattern matching results (for some type \ls$t$ determined
% by the context, \ls$unit$ in the example above).
% %
% It then tries to find a solution to the patterns and then applies the
% solution, if any, to the tactic \ls$k$.

In effect, \ls$with_pat$ hijacks the syntax
\guido{and type-inference: how?}
of \fstar to allow the programmer to build concise pattern-matching
clauses, and then interprets the syntax according to the DSL that it
provides---a testament to the flexibility of the quotation and syntax
inspection mechanisms that \metafstar provides.
\fi



\iffalse
The program below is a stateful recursive function (\ls$assignL$) that
copies the contents of a list \ls$l$ into a mutable array $b$. Stateful
computations in \fstar have computation types of the form
\ls$ST a (requires pre) (ensures post)$ (resembling Hoare
type theory~\cite{nmb08htt}), describing computations which when run
in an initial state \ls$h_0$ satisfying \ls$pre h_0$, return
results \ls$v:a$ in a final state \ls$h_1$ satisfying
%
\ls$post h_0 v h_1$.

\begin{lstlisting}
let rec assignL #a (l: list a) (b: array a): ST unit
  (requires (fun h0 -> live h0 b /\ length b = L.length l))
  (ensures (fun h0 _ h1 -> live h1 b /\ modifies {b} h0 h1 /\ contents b h1 == l))
= match l with
  | [] -> ()
  | hd :: tl -> b.(0ul) := hd; assignL tl (offset b 1ul)
\end{lstlisting}

The precondition of the \ls$assignL$ enforces temporal memory safety
(i.e., \ls$b$ must be live) and spatial safety (the array must be
large enough to hold the list). Its postcondition states that the
array is still live; that only the array is modified; and functional
correctness (i.e., the contents of the array in the final state match
the list).

Code like this is generic and convenient to verify using SMT only,
relying on carefully crafted \fstar libraries with E-matching
triggers. But, clients who want to initialize large arrays sometimes
prefer a long straight line of writes to this recursive,
case-analyzing function, e.g. for efficiency in the case of \fstar's C
backend. Of course, directly
verifying long sequences 
of ad hoc, heap mutating code is tedious boilerplate and, worse, as
discussed in depth in \autoref{sec:registers}, can perform poorly when
using SMT based proofs.

Its easy to get the best of both worlds with \metafstar since
metaprograms can transform verified code. When such transformations
are based on the partial evaluation supported by \fstar's logic, they
incur no addition verification cost, as shown below.

\paragraph{Unrolling a loop with partial evaluation}
The function \ls$initialize10$ initializes an array with the first 10
unsigned integers. Its body is a constructed by a metaprogram which
when evaluated during type-checking, fills the hole in the body
of \ls$initialize10$ with a fully unrolled version of \ls$assignL$.
The metaprogram works by using the \ls$solve_then$ combinator from
the \metafstar library to (1) construct the syntax of
%
\ls$assignL [0ul;...;9ul]$ by quoting it and assigning it
as the solution to the inference problem corresponding to the hole,
thereby using the compact specification of \ls{assignL} for
verifying the postcondition of \ls{initialize10};
%
(2) then invoking \fstar's normalizer to evaluate the partial
application \ls$assignL [0ul;...;9ul]$, which fully unrolls the
recursive function; and finally (3) fill the hole with the desired
straightline, effectful function. The resulting \lowstar code can then
be extracted to C using KreMLin, the compiler from \fstar to C.

\begin{lstlisting}
let initialize10(b:array uint32) : ST ... (ensures (fun h0 _ h1 -> ... contents b == [0ul;...;9ul]))
    (_ by solve_then (fun () -> exact (quote (assignL [0ul;...;9ul])))
                     (fun () -> norm [delta_only `%assignL; iota; zeta]))
\end{lstlisting}

\fi

\section{Related Work}
\label{sec:related}

\ch{Should also cite this very recent F-IDE 2018 paper
\url{https://arxiv.org/abs/1811.10814}}

\ch{Another tool we should probably compare with is KeY, which also
  has a notion of interactive proofs
  \url{https://pdfs.semanticscholar.org/2f5d/6602fbaa97fef0a9df1216451e536c33d617.pdf}.
  With a lot of work going into a usable GUI for proofs:
  \url{https://ieeexplore.ieee.org/document/7582776}.}

% \ch{Positioning: like in the intro trying to position ourselves
%   as competition to Dafny not Coq}

% Metaprogramming has a long and rich history, starting from
% LISP~\cite{McCarthy:1962:LPM:1096473}. We focus primarily on the use
% of metaprogramming in program verification tools and proof assistants.

Many SMT-based program verifiers~\cite{BarnettCDJL05,
BurdyCCEKLLP05, BarnettDFJLSV05, FlanaganLLNSS13, LeinoN98},
  % traditionally provide little support for interactive proofs, and rely instead
  % on SMT solvers
%
% They instead rely on SMT solvers and other first-order logic theorem
% provers
% for discharging VCs automatically.
%
% When SMT automation fails,
%% lacking a way to dynamically inspect the proof state,
rely on user hints, in the form of assertions and lemmas, to complete proofs.
% is left trying to add extra assertions or stating inductive lemmas,
% in the hope
% that these provide sufficient guidance to the solver.
This is the predominant style of proving
used in tools like Dafny~\cite{leino10dafny}, Liquid Haskell~\cite{VazouTCSNWJ18}, Why3~\cite{filliatre13why3}, and \fstar
itself~\cite{mumon}.
%
%When SMT automation fails, however, the user has to step in and narrow
%down the problem by adding extra assertions that can then also guide
%the SMT solver towards finding a proof.
%%
%Dafny~\cite{leino10dafny, Amin13, AminLR14} and \fstar~\cite{mumon}
%additionally allow users to prove inductive lemmas by writing pure
%recursive functions that return a trivial result (unit), but that can
%be called for the logical side-effect of their post-condition.
%%
%Similar ideas have recently been used in Liquid
%Haskell~\cite{VazouLP17, VazouTCSNWJ18}.
%%
%While this can be very useful for overcoming the SMT solver's
%limitations, it can be very difficult to figure out what lemmas to
%prove without a way to inspect the current proof goal.
%
% Also, this proof style can be very tedious in practice without
% metaprogramming support: for instance the non-linear arithmetic
% proofs we do automatically in \autoref{sec:semirings} used to take
% several dozens of manual invocations to lemmas such as associativity,
% commutativity, and distributivity.
%
However, there is a growing trend to augment this style of semi-automated proof with interactive proofs.
For example, systems like Why3~\cite{filliatre13why3} allow
VCs to be discharged
%not just using automatic theorem provers, but also
using ITPs such as Coq, Isabelle/HOL, and PVS, but this requires
an additional embedding of VCs into the logic of the ITP in question.
%
In recent concurrent work, support for {\em effectful} reflection
proofs was added to Why3~\cite{why3reflection}, and it would be
interesting to investigate if this could also be done in \metafstar.
%
\citet{GrovT16} present Tacny, a tactic framework for Dafny,
which is, however, limited in that it only transforms source code,
with the program verifier unchanged.
% a scripting layer on top of Dafny and separate from it that can
% be used to generate fragments of Dafny code to be verified by Dafny's
% existing SMT-based tool chain. While Tacny can be used to automate
% some verification tasks in Dafny, it is still heavily reliant on SMT
% solving as its only engine for proofs.
%
In contrast, \metafstar{} combines the benefits of an SMT-based
program verifier and those of tactic proofs within a single language.
% as an embedded DSL within a
% general purpose programming and verification-oriented host
% language.

Moving away from SMT-based verifiers, ITPs have long
relied on separate languages for proof scripting,
% Tactics and metaprogramming are the workhorse of interactive theorem
% provers,
starting with Edinburgh
LCF~\cite{gordon-edinburghLCF:1979} and ML,
%as its metaprogramming language.
%
and continuing with HOL, Isabelle and Coq, which are either extensible
via ML, or have dedicated tactic languages~\cite{isar, Delahaye00,
  AnandBCST18, StampoulisS10}.
%such as Isar for Isabelle~\cite{isar} and Ltac for Coq~\cite{Delahaye00}.
%
%
%HOL, Isabelle, and Coq still use ML even today
%to various extents for some of their metaprogramming tasks.
%%
%Most Isabelle proofs are, however, written in Isar, a specialized
%language for structured proofs~\cite{isar}.
%%
%Analogously, most Coq proofs are written in Ltac~\cite{Delahaye00,
%  ltac}, a small untyped functional language with primitive support
%for pattern-matching and back-tracking.
%
% Ltac: \url{http://coqhott.gforge.inria.fr/blog/coq-tactic-engine/}
%
% Various replacements for Ltac have been proposed~\cite{}.
% -- whatever, none of them really used
%
% Closest related work: Mtac, Idris, Lean
%
\metafstar builds instead on a recent idea in the space of dependently
typed ITPs \cite{ZilianiDKNV15, Mtac2, ChristiansenB16,
  EbnerURAM17} of reusing the %purely-functional programming
object-language as the meta-language.
%
This idea first appeared in Mtac, a Coq-based tactics framework
for Coq%
%using an abstract monad with primitive actions such as pattern matching
~\cite{ZilianiDKNV15, Mtac2},
% \ch{Not making a distinction between Mtac1 and 2, should we?}
%
and has many generic benefits including reusing the standard library,
IDE support, and type checker of the proof assistant.
%
Mtac can additionally check the partial
correctness of tactics, % using dependent types.
% by a shallow embedding %of terms into Coq
which is also sometimes possible in \metafstar but still rather limited
(\autoref{sec:specifying-metaprograms}).
% CH: replaced
% cannot yet model such properties; doing
% so remains for future work.
% . We choose instead to not prove any deep properties about our
% tactics, which is perhaps more suitable for a general-purpose metaprogramming
% framework that goes beyond tactics.
%
%While proving tactics correct can be sometimes appealing (as proofs by
%reflection show), requiring that all tactics are statically verified
%to either solve the goal or
%fail restricts the kind of metaprogramming tasks one can solve with Mtac.
%%
%Also proving tactics correct can be very challenging (if at all
%possible), which means that key tactics in Mtac are still only typed
%with a very uninformative type (basically an existential package).
%
\metafstar's design is instead more closely inspired by the
metaprogramming frameworks of Idris~\cite{ChristiansenB16} and
Lean~\cite{EbnerURAM17}, which provide a deep embedding of terms that
metaprograms can inspect and construct at will without dependent types
getting in the way.
%% %
%% This is very powerful and Lean allows most aspects of the proof
%% assistant to be metaprogrammed using monadic programs in do notation.
%% %
%% \fstar's support for effects allows such metaprograms to be written in ML-style.
%% %familiar since Milner's Edinburgh LCF~\cite{gordon-edinburghLCF:1979}
%% % which we believe still provides a good trade-off between
%% % expressiveness and static guarantees for metaprograms.
%% %
However, \fstar's effects, its weakest precondition calculus, and its
use of SMT solvers distinguish \metafstar from these other frameworks,
presenting both challenges and opportunities, as discussed in this
paper.


% \cpc{Do we want to mention SMT tactics somewhere?
% \url{https://z3prover.github.io/api/html/ml/Z3.Tactic.html}}
% \ch{maybe; does any program verification tool use them?
%   or what would be the connection to our work?}
% \cpc{Not entirely sure how often it's used.  Sounds like an
%   alternative to munging the VCs on the F* side.}

Some SMT solvers also include tactic engines~\cite{smttactics},
which allow to process queries in custom ways.
%
However, using SMT tactics from a program verifier is not very
practical.
%
To do so effectively, users must become familiar not only with the
solver's language and tactic engine, but also with the translation from
the program verifier to the solver.
%
Instead, in \metafstar, everything happens within a single language.
%
Also, to our knowledge, these tactics are usually coarsely-grained,
and we do not expect them to enable developments such as
\autoref{sec:seplogic}.
%
Plus, SMT tactics do not enable metaprogramming.


% Other verification frameworks~\cite{ynot-icfp08, nmb08htt,
%   Krebbers0BJDB17} %value metaprogramming and tactics so much that they
% build directly on top of Coq, even if until recently~\cite{CzajkaK17}
% this meant having no SMT automation whatsoever.

Finally, ITPs are seeing increasing use of ``hammers'' such as
Sledgehammer~\cite{BlanchetteBP13, PaulsonB10, Blanchette013} in
Isabelle/HOL, and similar tools for HOL Light and
HOL4~\cite{KaliszykU14}, and Mizar~\cite{KaliszykU15a}, to interface
with ATPs.
%
This technique is similar to \metafstar, which, given its support for
a dependently typed logic is especially related to a recent hammer for
Coq~\cite{CzajkaK17}. Unlike these hammers, \metafstar does not aim to
reconstruct SMT proofs, gaining efficiency at the cost of trusting the
SMT solver. Further, whereas hammers run in the background, lightening
the load on a user otherwise tasked with completing the entire
proof, \metafstar relies more heavily on the SMT solver as an end-game
tactic in nearly all proofs.

% \fstar{}'s SMT encoding is targeted at soundness and at efficient,
% scalable, and reproducible SMT solver behavior, while hammers often
% give up many of these and then regain them by proof reconstruction.


% Other tactic frameworks:
% Ltac, SSReflect, Isabelle Isar, etc.
% Tactics in Isabelle: \cite{MatichukWM14} -- not the main thing!
% http://ts.data61.csiro.au/publications/nicta_full_text/7847.pdf

% Auto-active: Leino and Moskal~\citep{LeinoM10} categorize program
% verifiers into two classes, based on where in this process user input
% is supplied. In their words:

\section{Conclusions}
\label{sec:conclusions}

A key challenge in program verification is to balance
automation and expressiveness. Whereas tactic-based ITPs
support highly expressive logics, the tactic author is responsible for all
the automation.
%
Conversely, SMT-based program verifiers
provide good, scalable automation for comparatively weaker logics, but offer
little recourse when verification fails. A design that allows picking the
right tool, at the granularity of each verification sub-task, is a
worthy area of research.
%
\metafstar presents a new point in this space: by using hand-written
tactics alongside SMT-automation, we have written
proofs that were previously impractical in \fstar, and
(to the best of our
knowledge) in other SMT-based program verifiers.

% Tactics are just one flavor of metaprogramming. Using
% \metafstar, we are now able to script the behavior of \fstar in \fstar in a type-safe
% manner, e.g. implementing type classes as a user program. Tactics also enable
% synthesis of sophisticated \fstar programs, e.g., memory-safe and correct-by-construction
% low-level parsers and serializers, that can be extracted in a semantics-preservering manner
% to C.

\iffull
This flurry of new use-cases is backed by an efficient native
compilation scheme.
% , which we capture via a formal model of multi-language interoperability.
%
Natively evaluating metaprograms allows to dynamically extend
\fstar with type-safe, semantically correct (by the safety of
\ls$TAC$), user-defined custom behavior without a performance
penalty---witnessing the close interoperability between \fstar and its
metalanguage.
\fi


% Hacks
\newcommand\grantsponsor[3]{#2}
\newcommand\grantnum[2]{#1-#2}

%% Acknowledgments
\ifcamera
\paragraph{Acknowledgements}
  We thank Leonardo de Moura
  and the Project Everest team for many useful discussions.
  The work of
    Guido Mart\'inez,
    Nick Giannarakis,
    Monal Narasimhamurthy,
  and Zoe Paraskevopoulou
  was done, in part, while interning at Microsoft Research.
  Cl\'ement Pit-Claudel's work was in part done
  during an internship at Inria Paris.
  The work of Danel Ahman, Victor Dumitrescu, and C\u{a}t\u{a}lin Hri\c{t}cu
  is supported by the MSR-Inria Joint Centre and the
  \grantsponsor{1}{European Research Council}{https://erc.europa.eu/}
  under ERC Starting Grant SECOMP (\grantnum{1}{715753}).
\fi

%% Bibliography
%\bibliography{bibfile}


%% Appendix
\iffull
\newpage
\appendix
\begin{figure}
\begin{lstlisting}
Lemma poly_multiply
  (n p r h r0 r1 h0 h1 h2 s1 d0 d1 d2 hh : Z)
  (H: p > 0 /\ r1 >= 0 /\ n > 0 /\ 4 * (n * n) = p + 5 /\ r = r1 * n + r0 /\
            h = h2 * (n * n) + h1 * n + h0 /\ s1 = r1 + (r1 / 4) /\ r1 mod 4 = 0 /\
            d0 = h0 * r0 + h1 * s1 /\ d1 = h0 * r1 + h1 * r0 + h2 * s1 /\
            d2 = h2 * r0 /\ hh = d2 * (n * n) + d1 * n + d0)
  : ((h * r) mod p = hh mod p).
Proof.
  repeat match goal with H : ?A /\ ?B |- _ => destruct H end.
  subst.
  match goal with |- ?A mod p = ?B mod p => change (eqm p A B) end.
  match goal with K : r1 mod 4 = 0 |- _ =>
   apply Z_div_exact_full_2 in K; try omega;
   revert K;
   generalize (r1 / 4);
   intro u;
   intros; subst
  end.
  pose (b := (h2 * n + h1) * u).
  generalize (Zopp_eqm p).
  generalize (Zplus_eqm p).
  generalize (Zminus_eqm p).
  generalize (Zmult_eqm p).
  generalize (eqm_setoid p).
  set (e := eqm p).
  intros.
  match goal with |- e _ ?R =>
    generalize (modulo_addition_lemma R p b)
  end.
  fold e. intro K. setoid_rewrite <- K. clear K.
  assert (p = 4 * (n * n) - 5) as J by omega.
  replace (b * p) with (b * (4 * (n * n) - 5)) by congruence.
  unfold b.
  apply eq_implies_eqm.
  ring.
Qed.
\end{lstlisting}
\caption{Coq proof of \ls$poly_multiply$}
\label{fig:coqpolymultiply}
\end{figure}

\section{Metaprogramming verified parsers and serializers}
\label{sec:parsers-detail}

We consider parser and serializer specifications as ghost code that
``operates'' on finite sequences of bytes of unbounded lengths.
For a given type $t$, a serializer for $t$ marshals any data of type
$t$ into a finite sequence of bytes; and a parser for $t$ reads from a
sequence of bytes, checks whether it corresponds to a valid data of
type $t$, and if so, returns that parsed data and the number of bytes
consumed.

The serializer specification always succeeds; the parser specification
succeeds if and only if the input data is valid with respect to the
parser. On those specifications, we aim to prove that the parser and
the serializer are partial inverses of each other. We encode these
requirements as refinements on the types of parser and serializer
specifications:
%
\begin{lstlisting}
type byte = FStar.UInt8.t
type parser t = (input: seq byte) -> GTot (option (t * (x: nat { x < length input } )))
let serialize_then_parse_eq_id #t (p: parser t) (s: (t -> GTot (seq byte))) =
  (forall (x: t) . let y = s x in p y == Some (x, length y))
let parse_then_serialize_eq_id #t (p: parser t) (s: (t -> GTot (seq byte))) =
  (forall (x: seq byte) . match p x with | Some (y, len) -> s y == slice x 0 len | _ -> True) })
type serializer #t (p: parser t) = (s: (t -> GTot (seq byte))
  { serialize_then_parse_eq_id p s /\ parse_then_serialize_eq_id p s })
\end{lstlisting}

Whereas those specifications are ghost code, we would like to generate
implementations that can be extracted to C. To this end, we generate
\emph{stateful} implementations written in the \lowstar set of \fstar~\cite{lowstar},
operating on buffers, which are \lowstar mutable data structures
representing C arrays. There, instead of a sequence of bytes, the
parser implementation is given an input buffer and its length; and the
serializer implementation is given an output buffer (along with its
length) onto which it is to serialize the data. This means that, given
a piece of data to serialize and a destination buffer, a serializer
implementation can succeed only if the serialized data fits into the
destination buffer; if so, then the serializer will return the number
of bytes written.

We bake the correctness of the implementations with respect to their
specifications at the level of their types:

\begin{lstlisting}
type buffer8 = LowStar.Buffer.buffer FStar.UInt8.t
type u32 = FStar.UInt32.t

let parser32_postcond #t (p: parser t) (input: buffer8) (l: u32 {l == len input})
  (h: mem) (res: option (t * u32)) (h': mem) =
  modifies loc_none h h' /\ (* memory safety *)
  match p (as_seq h input), res with (* functional correctness *)
  | Some (x, ln), Some (x', ln') -> x' == x /\ FStar.UInt32.v ln' == ln
  | None, None -> True   | _ -> False

type parser32 #t (p: parser t) = (input: buffer8) -> (l: u32 { l == len input } )
                                              -> Stack (option (t * u32))
  (requires (fun h -> live h input)) (ensures (parser32_postcond p input))

let serializer32_postcond #t (p: parser t) (s: serializer p) (output: buffer8)
  (l: u32 {l == len output}) (x: t) (h: mem) (res: option u32) (h': mem) =
  live h output /\ live h' output /\ modifies (loc_buffer b) h h' /\ (* memory safety *)
  let ln = length (s x) in
  match res with (* functional correctness *)
  | Some ln' ->
    FStar.UInt32.v ln' == ln /\ ln <= FStar.UInt32.v l /\
    let b' = sub b 0ul ln' in modifies (loc_buffer b') h h' /\ as_seq h' b' == s x
  | None -> ln > FStar.UInt32.v l

type serializer32 #t (#p: parser t) (s: serializer p) =
  (output: buffer8) -> (l: u32 { l == len output } ) -> (x: t) -> Stack (option u32)
  (requires (fun h -> live h output)) (ensures (serializer32_postcond s output l x))
\end{lstlisting}

\tahina{TODO: gen\_parser can be implemented using typeclasses, triggering gen\_enum\_parser.}

\paragraph{\bf Generating the specification from a \fstar datatype}

Instead of mandating the user to write the parser and serializer
specification themselves, we write a \metafstar
tactic, \verb+gen_specs+, to generate both
the parser and the serializer specifications directly from the \fstar
type $t$ given by the user. These tactics operate by syntax inspection
on $t$ itself, and generates the parser and serializer specifications
using our combinators.

Contrary to Coq's Ltac, \metafstar tactics can also inspect
definitions of inductive types. From there, we made \verb+gen_specs+
also generate a parser and serializer specification out of an
enumeration type defined as a \fstar inductive type whose constructors
have no arguments: for such a type with $n \leq 256$ constructors, the
corresponding parser will associate a unique 8-bit integer from $0$ to
$n-1$ to each constructor. For instance, if $t$ is defined as a \fstar
enumeration type:

As before, by virtue of the correctness of serializers being baked in
their type, the correctness of the serializer generated
by \verb+gen_specs+ does not appear as such as a verification
condition when \fstar is type-checking the generated term. In the
example above, the only proof obligations generated are those to
correctly type-check the rewriting functions themselves passed
to \verb+serialize_synth+ (e.g. enum constructors are rewritten to
integers less than 4), and the precondition of \verb+serialize_synth+
(i.e. the fact that the two rewriting functions are inverse of each
other.) These can be discharged by the SMT solver, but they can also
be discharged automatically by our tactics directly by case analysis
on the inductive type, or by bounded integer enumeration.

\paragraph{\bf Binders}

For instance, we specified \verb+seq_p+, a combinator for dependent
parsing:\protect{\footnote{This parsing combinator is based on yet
unpublished work by Tej Chajed. Building serializers for parsers
derived from \ls$seq_p$ is explicitly out of the scope of this
paper.}} \tahina{I mentioned Tej in the footnote. This part of the
footnote should be removed if Tej becomes a coauthor of the paper.}

\begin{lstlisting}
let seq_p #t1 (p1: parser t1) #t2 (p2: (t1 -> Tot (parser t2))) : Tot (parser t2) =
  fun input -> match p1 input with | None -> None | Some (x1, cons1) ->
    let input2 = slice input cons1 (length input) in
    match p2 x1 input2 with
    | None -> None | Some (x2, cons2) -> Some (x2, cons1 + cons2)
\end{lstlisting}

Then, we implemented it in \lowstar (\verb+seq_pi+), and
manually proved it correct. From there, the user can write a simple
parser specification for a parser that parses an integer encoded
in one or two bytes depending on its value (where \verb+parse_ret x+
is a parser that returns \verb+x+ without reading its byte input:)

\begin{lstlisting}
let example : parser_spec FStar.Int16.t =
  seq_p parse_u8
           (fun lo -> if lo <^ 128uy then parse_ret (cast_u8_to_i16 lo)
                   else parse_synth parse_u8 (fun hi -> cast_u8_to_i16 (lo %^ 128uy)
                                                 +^ cast_u8_to_i16 hi))
\end{lstlisting}
%
Then, if the user writes:
%
\begin{lstlisting}
let example_impl : parser_impl example = _ by gen_parser_impl
\end{lstlisting}
%
then \verb+gen_parser32+ inspects the shape of the \emph{goal},
which is \verb+parser32 example+, and so generates the following \lowstar implementation:
%
\begin{lstlisting}
seq_pi parse32_u8 (fun lo ->
  parse32_ifthenelse (lo <^ 128uy) (fun _ -> parse32_ret (cast_u8_to_i16 lo)) (fun _ -> parse32_synth parse32_u8 (fun hi -> cast_u8_to_i16 (lo %^ 128uy) +^ cast_u8_to_i16 hi)))
\end{lstlisting}

from where KreMLin \cite{lowstar} then inlines the corresponding implementation combinators to produce C code.

First, \verb+gen_parser32+ inspects the shape of the goal to determine
the type of the data to be parsed and the parser specification; then,
it calls a recursive tactic \verb+gen_parser32'+ that inspects the
shape of the parser specification and actually builds the
implementation.

\begin{lstlisting}
module T = FStar.Tactics
let rec gen_parser32' (env: T.env) (t: T.term) (p: T.term) : T.Tac T.term =
  let (hd, tl) = app_head_tail p in
  if hd `T.term_eq` (`(parse_ret)) then T.mk_app (`(parse32_ret)) tl else
  if hd `T.term_eq` (`(parse_u8)) then (`(parse32_u8)) else
  if hd `T.term_eq` (`(seq_p)) then match tl with
  | [(t, _); (p, _); (t', _); (p', _)] ->
    begin match T.inspect p' with
    | T.Tv_Abs bx body ->
      let p32 = gen_parser32' env k t p in
      let env' = T.push_binder env bx in
      let body' = gen_parser32' env' k' t' body in
      let p32' = T.pack (T.Tv_Abs bx body') in
      (`(seq_pi #(`#t) #(`#p) (`#p32) (`#t') (`#p') (`#p32')))
    | _ -> ...
    end
  | _ -> ...
  else
  ...

let gen_parser32 () : T.Tac unit =
  let (hd, tl) = app_head_tail (T.cur_goal ()) in
  if hd `T.term_eq` (`parser32)
  then match tl with
  | [(t, _); (p, _)] ->
    let env = T.cur_env () in
    let p32 = gen_parser32' env t p in
    T.exact_guard p32;
  | _ -> ...
\end{lstlisting}

When \verb+gen_parser32+ (resp. \verb+gen_serializer32+) builds the
parser (resp. serializer) implementation, it may generate verification
conditions due to specific preconditions required by certain
combinators (such as the precondition of the \verb+serialize_synth+
rewriting serializer combinator: the rewriting functions being inverse
of each other). Contrary to Coq, there are no proof objects, so those
preconditions need to be proven again even if they had been proven
earlier by the user at the specification level. Nevertheless, those
preconditions can be automatically solved by our tactics, or they can
directly be sent to the SMT solver.

However, by virtue of the correctness of the implementation
combinators being baked in their type, typechecking the resulting term
generated by \verb+gen_parser32+ triggers no verification condition
other than those preconditions required by the specific combinators,
and unification constraints, the latter solved by reflexivity. In
particular, the memory safety and correctness of the implementations
generated by \verb+gen_parser32+ do not appear as such as verification
conditions when \fstar is type-checking the terms generated
by \verb+gen_parser32+. \tahina{Is it already important to mention
here the kinds of proof obligations generated by the tactics, and that
they can be solved either by SMT or tactics? Or should we move to
section 5?}

\subsection{Verified meta-programming of program transformations}

\ls$compile_bind$ expects, in \ls$bind f1 f2$, that \ls$f2$ be an
abstraction, and so needs to manipulate binders and recursively
compile \ls$f1$ and the body of \ls$f2$. This will, in the latter
case, make some bound variables appear free in terms being compiled.

If the head term $t_0$ is a free variable not corresponding to any
combinator, then \ls$compile_fvar$ unfolds it, compiles the unfolded
term, and inserts a coercion.\footnote{We defined an explicit coercion
combinator, \ls$coerce_sz$, such that \ls$coerce_sz f1 f1_sz f2$ is
well-typed and is a \ls$m_sz f2$ as soon as \ls$f1 () == f2 ()$ as $m$
specifications.} In most cases, a free variable to be unfolded will
correspond to a function definition, so that unfolding will yield a
$\beta$-redex, which \ls$compile_fvar$ reduces by calling a
normalization tactic that we provide as part of \metafstar. This needs
an \emph{environment}
\ls{(e: env)} tracking the bound variables encountered by nested calls
to \ls$compile_bind$.

We show below the implementation of \ls$compile_bind$, which is
representative of most features of \metafstar that we use for
our \ls$compile$ tactic.

\begin{lstlisting}
 (* extract the head symbol and the arguments of an application *)
let rec app_head_rev_tail (t: term) : Tac (term * list argv) = 
  match inspect t with | Tv_App u v -> let (x, l) = app_head_rev_tail u in (x, v :: l) | _ -> (t, [])
let app_head_tail (t: term) : Tac (term * list argv) = let (x, l) = app_head_rev_tail t in (x, rev l)

let compile_bind (e: env) (ty: term) (t: term) (compile: env -> term -> term -> Tac term) : Tac term =
  let (t_0, ar) = app_head_tail t in
  guard (t_0 = quote bind); (* fail if false *)
  match ar with
  | [ ($\tau_1$, _); ($\tau_2$, _); (f1, _); (f2, _); ] ->
    begin match inspect f2 with
    | Tv_Abs v f2_body ->
      let f1' = compile e $\tau_1$ f1 in
      let e' = push_binder e v in
      let f2_body' = compile e' $\tau_2$ f2_body in
      let f2' = pack (Tv_Abs v f2_body') in
      (`(bind_sz #(`#tau_1) #(`#tau_2) (`#f1) (`#f1') (`#f2) (`#f2')))
    | _ -> fail "compile_bind: Not an abstraction"
    end
  | _ -> fail "compile_bind: 4 arguments expected"
\end{lstlisting}

\paragraph{\bf Putting everything together}
Finally, the following wrapper combinator $\Phi$ puts everything
together, taking a specification $f$ and its size and buffer
implementations, and yielding a stateful function, fitting in
the \lowstar subset of \fstar amenable to extraction to C, that first
computes the expected size of the output, then allocates one single
buffer using the C \ls$malloc$, then uses it to write the output of
$f$:
\begin{lstlisting}
let $\Phi$ $\tau$ (f: m $\tau$) (f_sz: m_sz f) (f_st: m_st f) : ST (option ($\tau$ * buffer U8.t)) (requires (fun _ -> True))
  (ensures (fun h res h' -> let (r, log) = f () in
    if length log > $2^{31} - 1$ then res == None else
    Some? res /\ (let (Some (r', b)) = res in r == b /\ fresh b h h' /\  as_seq h b == log)
  )) = match f_sz with | None -> None | Some (_, sz') ->
       let b = rcreate_mm root 42uy sz in (* then allocate a fresh buffer *)
       let (r, _) = f_st b in (* then write the output into the buffer *)
       Some (r, b)
\end{lstlisting}

%%% Tahina: this paragraph removed because it should rather go to a
%%% parser-specific paper.
%%% CH: Added this back, Nada is on the ESOP %%% PC!
\paragraph{\bf Related Work}
Our approach of generating verified low-level formatters is related to
Amin and Rompf's~\cite{Amin:2017:LAW:3009837.3009867} work on
LMS-Verify. They use metaprogramming facilities in Scala to generate C
code together with proof annotations to be checked by
Frama-C~\citep{CuoqKKPSY12}, an SMT-based C verifier. In contrast, we
generate correct-by-construction imperative \lowstar code, which is
verified by \fstar{} before being translated to C by the KreMLin
compiler.  LMS-Verify has also been used to generate efficient and
safe HTTP parsers, a larger-scale effort than the network message
formatters than we have currently done.

%%%%% Appendix for Section 4

\newpage

\section{Proof of \autoref{thm:split}}
\label{app:correctness}
This proof is based on the formal definition of \emf, a recent formalization
of an \fstar{} subset~\cite{dm4free}.
%x
We use the same notation and rule names from there, but the proof is
nevertheless quite direct with just minimum familiarity with \emf.
%
We use $E$ to represent \fstar contexts:

\[
\begin{array}{rclcl}
    E &=&    \cdot  & & \\
      & \mid & t E &                 \mid & E t \\
      & \mid & \lambda(x:t). E &     \mid & \lambda(x:E). t \\
      & \mid & (x:t) \rightarrow E & \mid & (x:E) \rightarrow t \\
      & \mid & x:t\{E\} &            \mid & x:E\{t\} \\
      & ~ & \ldots
\end{array}
\]

Contexts present a set of bound variables to their hole. We denote
those variables by $\gamma(E)$:

\[
\begin{array}{lcl}
    \gamma(\cdot) &=& \emptyset \\
    \gamma(\lambda(x:t). E) &=& \{x\} \cup \gamma(E) \\
    \gamma(\lambda(x:E). t) &=& \gamma(E) \\
    \ldots
\end{array}
\]
    % γ(∘) = {}
    % γ(λ(x:t). E) = (x:t) ++ γ(E)
    % γ(λ(x:E). t) = γ(E)
    % ...

We say a context $E$ has type $t_1 \Rightarrow t_2$ for $\Gamma$ (noted
by $\Gamma \vdash E : t_1 \Rightarrow t_2$) when for all $e$ such that
$\Gamma, \gamma(E) \vdash e : t_1$ we also have $\Gamma \vdash E[e] : t_2$.

\subsection{Soundness of splitting the proof obligation}

The proof follows mainly from the following lemma:

\begin{lemma}
    \[
    \infer{\Gamma \vDash E[t_1] = E[t_2]}
          {\Gamma, \gamma(E) \vDash t_1 = t_2}
    \]
\end{lemma}
\begin{proof}
    The proof follows from applying functional extensionality and using
    a carefully crafted abstraction.

    Let $S = \lambda f. E[f (\gamma(E))]$ (where $f (\gamma(E))$
    represents $f$ applied to every variable in $\gamma(E)$,
    sequentially). Also, $\lambda \gamma(E). t$ represents
    abstracting $E$ over the variables in $\gamma(E)$, and
    similarly for $\forall \gamma(E). t$.

    Then,
    \[
        \infer[\mbox{Reduction}]{\Gamma \vDash E[t_1] = E[t_2]}
        {\infer[\mbox{V-EqP}]{\Gamma \vDash S (\lambda \gamma(E). t_1) = S (\lambda \gamma(E). t_2)}
        {\infer[\mbox{V-Ext}]{\Gamma \vDash (\lambda \gamma(E). t_1) = (\lambda \gamma(E). t_2)}
        {\infer[\mbox{V-$\forall$i}]{\Gamma \vDash \forall \gamma(E). t_1 = t_2}
                             {\Gamma, \gamma(E) \vDash t_1 = t_2}}}}
    \]

    Note that the particular set of contexts E does not influence
    the proof.



        % Γ, γ(E) |= t1 = t2
        % ------------------------- [V-∀i, repeatedly]
        % Γ |= forall γ(E). t1 = t2
        % ------------------------------ [V-Ext, repeatedly]
        % Γ |= (λγ(E). t1) = (λγ(E). t2)
        % --------------------------------- [V-EqP, with P = S]
        % Γ |= S (λy(E). t1) = S (γ(E). t2)
        % --------------------------------- [Red]
        % Γ |= E[t1] = E[t2]
\end{proof}

Having such lemma, the theorem is proven as follows:
\[
    \infer{\Gamma \vDash E[\phi]}
         {\Gamma \vDash E[\top]
         & \infer[\mbox{Lemma}]{\Gamma \vDash E[\phi] = E[\top]}{
         \infer[\mbox{PropExt}]{\Gamma, \gamma(E) \vDash \phi = \top}{
         \infer[\mbox{Trivial}]{\Gamma, \gamma(E) \vDash \phi \iff \top}
                               {\Gamma, \gamma(E) \vDash \phi}}}
         }
\]


\subsection{Partial completeness of splitting the proof obligation}

While the previous theorem allows to soundly split any
subformulae within a VC, we have seen that some of them
constraint the system.
%
Here we prove that for a particular set of contexts, splitting
does not make the judgments stronger, and via Theorem 1, then they are
equivalent.

We define a particular shape of ``positive'' contexts $P$.
We don't attempt to follow all of EMF*'s syntax, just the parts
we commonly see as VCs in practice.

\[
\begin{array}{rclcl}
    P &=&    \cdot  & & \\
       & \mid & \phi \land P            &\mid& P \land \phi \\
       & \mid & \forall (x:t). P \\
\end{array}
\]

In \emf, an implication $a \Rightarrow b$ is just sugar for
$\forall (\_:a). b$, so they are considered as well.

% This is sufficient to catch most VCs we encounter. Most likely, a reason
% is that disjunction/existentials can never appear in strictly-positive
% positions due to the inherent conjunctivity of WPs, but we haven't
% proved it. In any case, this property is not crucial in any way.

\begin{theorem}
    Take $\Gamma$, $E$, and $\phi$ such that:
    (1) $\phi$ is squashed,
    (2) $\Gamma \vdash E : prop \rightarrow prop$
    and (3) $\Gamma \vdash \phi : prop$.
    Then the following holds

    \[
        \infer{\Gamma, \gamma(P) \vDash \phi \qquad \Gamma \vDash P[\top]}
              {\Gamma \vDash P[\phi]}
    \]
\end{theorem}
\begin{proof}
    By induction on the structure of $P$, keeping $\Gamma$ and $\phi$
    universally quantified.
    \begin{itemize}
        \item For $P = \cdot$, trivial
        \item For $P = \psi \land P'$, our hypothesis is
              $\Gamma \vDash \psi \land P'[\phi]$.
              %
              Hence,
              $\Gamma \vDash \psi$ and
              $\Gamma \vDash P'[\phi]$.
              %
              By the IH, we get
              $\Gamma, \gamma(P') \vDash \phi$
              (note that $\gamma(P) = \gamma(P')$)
              and
              $\Gamma \vDash P'[\top]$.
              %
              By weakening and conjunction, 
              we can prove $\Gamma, \gamma(P') \vDash \psi \land \phi$
              and $\Gamma \vDash \psi \land P'[\top]$, and conclude.

        \item For $P = P' \land \psi$, the reasoning is analogous.

        \item For $P = \forall (x:t). P'$, our hypothesis is
              $\Gamma \vDash \forall (x:t). P'[\phi]$. Hence,
              we get $\Gamma, x:t \vDash P'[\phi]$ by
              eliminating the quantifier.
              %
              From the IH, we get
              $\Gamma, x:t, \gamma(P') \vDash \phi$
              and
              $\Gamma, x:t \vDash E[\top]$.
              %
              From this last one, we can use V-$\forall$i to
              obtain $\Gamma \vDash \forall (x:t). E[\top]$, and conclude
              (noting that $\gamma(P) = x:t, \gamma(P')$).
    \end{itemize}
\end{proof}


\newpage

\section{Modelling native plugins}
\label{app:native}

This appendix contains the formal definitions of our multi-language
interoperability model. We formalize the semantics of source and target language
as well as the translation between the two. We convey the essential ideas 
in the text below, and present full details in 
Figures \ref{fig:sourcelanglong}--\ref{fig:translationlong} that follow the discussion below.

\subsection*{Part 1: Modeling Native Plugins with Simple Types}

%We now present our design of native plugins using a small
%multi-language calculus, a model of interactions between the
%normalizers and native plugins, with appropriate conversions
%between the term representations. We intend to only convey the essential ideas
%of our implementation. A full formalization is out of scope, although more details are in the supplementary material.

\iffalse
\begin{figure}[h]
  \centering
\[\begin{array}{@{}l@{~~}l@{~}r@{}l@{}}
\mbox{Types} & \stau & \bnfdef & \stunit% \bnfalt \svar{A}
      \bnfalt \sarrow{\stau_{\source{1}}}{\stau_{\source{2}}}
      \bnfalt \sprod{\stau_{\source{1}}}{\stau_{\source{2}}} \bnfalt \ssum{\stau_{\source{1}}}{\stau_{\source{2}}} \\
& & & \\

\mbox{Expressions} & \sexp & \bnfdef & \sunit \bnfalt \svar{x} \bnfalt \sabs{\svar{x}}{\stau}{\sexp} \bnfalt \sapp{\sexp_{\source{1}}}{\sexp_{\source{2}}} \\
%& & \bnfalt & \stapp{\sexp}{\stau}\\
& & \bnfalt & \spair{\sexp_{\source{1}}}{\sexp_{\source{2}}} \bnfalt \sfst{\sexp} \bnfalt \ssnd{\sexp} \\
& & \bnfalt & \sinl{\sexp} \bnfalt \sinr{\sexp} \\
& & \bnfalt & \scasex{\sexp}{\stau_{\source{1}}}{\sexp_{\source{1}}}{\stau_{\source{2}}}{\sexp_{\source{2}}} \\
& & \bnfalt & \salien{\texp}{\stau}{\sdelta} \\
% \bnfalt \sopaque{\texp}
\end{array}\]
  \caption{Syntax of the source language}
  \label{fig:syntax-source}
\end{figure}

\begin{figure}[h]
  \centering
\[\begin{array}{@{}l@{~~}l@{~}r@{}l@{}}
\mbox{Expressions} & \texp & \bnfdef & \tunit \bnfalt \tvar{x} \bnfalt \tabs{\tvar{x}}{\texp} \bnfalt \tapp{\texp_{\target{1}}}{\texp_{\target{2}}} \\
& & \bnfalt & \tpair{\texp_{\target{1}}}{\texp_{\target{2}}} \bnfalt \tfst{\texp} \bnfalt \tsnd{\texp} \\
& & \bnfalt & \tinl{\texp} \bnfalt \tinr{\texp} \\
& & \bnfalt & \tcasex{\texp}{\texp_{\target{1}}}{\texp_{\target{2}}} \\
& & \bnfalt & \talien{\sexp}{\stau}{\sdelta} \bnfalt
%\topaque{\sexp}\\

\end{array}\]
  \caption{Syntax of the target language}
  \label{fig:syntax-target}
\end{figure}
\fi

Reflecting our requirement that plugins are only supported at
ML-typeable interfaces, we start with a
source language that is an intrinsically typed (Church
style), standard, simply typed lambda calculus with pairs and
sums. Later in the section, we
will add ML-style rank 1 polymorphism to it.
%
Conversely, reflecting
the type-less native representation of compiled OCaml code, our target
language is the untyped lamdba calculus.
For clarity, we markup the the syntax using colors, using
$\source{blue}$ for the source language and $\target{red}$ for the target.
Aside from the standard forms, we
have two new expression forms to account for the ``alien''
expressions of one language appearing in the other:
$\talien{\sexp}{\stau}{}$, which \emph{embeds} a
$\stau$-typed source language expression into the target language, and 
$\salien{\texp}{\stau}{}$,
which \emph{unembeds} a target language expression to the source language at
type $\stau$.

%% $\salien{\texp}{\stau}{}$,
%% which \emph{unembeds}
%% a target computation in the source language at type $\stau$, and
%% $\talien{\sexp}{\stau}{}$, which
%% \emph{embeds} a $\stau$-typed source computation in the target language.

%% To model native plugins, consider a simply-typed lambda calculus
%% source term $\source{E(\sexp : \stau)}$, a context $\source{E}$
%% surrounding a term $\sexp : \stau$.
%% %
%% % JP: I think these two were salien as opposed to talien?
%% We translate $\sexp$ to an untyped lambda term $\target{\texp}$ (by
%% erasing type information) and unembed it in the original source
%% context as $\source{E (\salien{\target{\texp}}{\stau}{})}$.
%% %
%% We study the interactions of such hybrid terms, aiming to ensure that
%% $\source{E (\salien{\target{\texp}}{\stau}{})}$ is observationally
%% equivalent to $\source{E(\sexp : \stau)}$.

We explain the reductions through a small example. Consider a source language term
$\sapp{\source{(}\sabs{\svar{x}}{\mathbb{\source{Z}}}{\svar{x}}\source{)}}{\source{0}}$, where we want to compile the
identity function to native code and then perform the application. The
first step is to translate the function to the target, by systematically
erasing the types, and unembed it in the source. In our example, the
source language term then becomes
$\sapp{\salien{\tabs{\tvar{x}}{\tvar{x}}}{\mathbb{\source{Z}} \source{->} \mathbb{\source{Z}}}{}}{\source{0}}$.


%% The main characteristic of the rules is that whenever a target term blocks
%% reduction in the source language, we $\sforce{}$ the term into a source value of
%% the appropriate type, using the unembedding function $\scoerce\stau\relax\texp$.
%% When $\sforce{\svar}$ is applied at function type, this involves passing a
%% source term as an argument to a target language function,
%% after \emph{unembedding} said function as a source term, i.e., although not
%% present initially, unembeddings of target terms in the source arise as during
%% reductions.
%% \guido{First use of $\scoerce\stau\relax\texp$, confused. Last sentence
%% seems explained backwards to me, hard to follow.}

%% \jonathan{I'm presenting both reduction semantics in this paragraph, then
%% jumping to the description of the two coercion functions.}

%% \nik{I suggest presenting embeddings first, since
%% this is the direction with which things ``start'' in the plugin
%% scenario.}

%% \zp{What about this structure:
%% \begin{itemize}
%% \item Embeddings : syntax + semantics (reduction and coersions)
%% \item Unembeddings : ditto
%% \item Other (non-standard) rules of the semantics + force
%% \end{itemize}
%% }

% \nik{remove $\delta$ at first and the protection function at first.}
%
%% \begin{mathpar}
%%  \inferrule*[right=Alien]
%%             {\sexp \sevals \sval}
%%             {\talien{\sexp}{\stau}{\relax} \teval \tcoerce{\stau}{\relax}{\sval}}
%% \end{mathpar}

\paragraph{\bf Source semantics and unembedding (\autoref{fig:sourceseman})}
At a beta-reduction step (e.g. rule {\sc{S-App}} below), the source
semantics uses a meta-function
$\sforce{\source{v}}$ to examine the head of the application $\svar{v}$.
In addition to the usual values, the term
$\salien{\tval}{\stau}{\relax}$ is a source value \emph{iff}
$\tval$ is a target value. If the head is such a
value, $\sforce{}$ invokes the unembedding coercion
$\scoerce\stau\relax\tval$ to coerce it to a suitable source
value. For structural types, this amounts to lazily coercing the head
constructor of the term (see below). For function types, a source-level closure is
allocated, which first \emph{embeds} its argument in the target,
reduces a target application (using rule {\sc{S-Alien}}), and
then \emph{unembeds} the result back in the source.
%
%
%% $\sforce{}$ examines its argument and invokes the the embedding meta-function
%% $\scoerce\stau\relax\tval$ whenever it hits an alien value. The embedding
%% meta-function $\scoerce\stau\relax\tval$ takes a target \emph{value} (the target
%% is CBV) and coerces it suitably as a source \emph{value}.
%% \zp{The result of the coercion will always be a value in the simply typed world.
%% It can be an expression when we add poly types.}
%
%
%
%% To define the semantics of the two languages we need a way to interpret alien terms
%% in both contexts. 
%%  The source reduces an alien term $\salien{\texp}{\stau}{\relax}$
%% by invoking the semantics of the target language to obtain a target value (rule {\sc Alien}).
%% Hence, the term $\salien{\tval}{\stau}{\relax}$ is a value in the source \emph{iff}
%% $\tval$ is a value in the target language. Whenever a reduction step is blocked
%% because of an alien value, the semantics will coerce this value to a source value by a
%% call to the meta-function $\sforce{}$.
%
\begin{figure}
\begin{mathpar}
  \inferrule*[right=S-Alien]
             {\texp \tevals \tval}
             {\salien{\texp}{\stau}{} \seval \salien{\tval}{\stau}{}}
  %
%  \and
%  \inferrule*[right=Proj]
%             {\sforce{\sval}=\spair{\sexp_{\source{1}}}{\sexp_{\source{2}}}}
%             {\sval.\source{i} \seval \sexp_{\source{i}}}
 %
  \and
  %
  \inferrule*[right=S-App]
             {\sforce{\sval}=\sabs{\svar{x}}{\stau}{\sexp}}
             {\sapp{\sval}{\sexpii} \seval \sexp \ssubst{\svar{x}}{\sexpii}}
 %
 \and
  \inferrule*{}{
    \sforce{\sval} =
    \begin{cases}
      \scoerce{\stau}{\relax}{\tval}
      & \text{if }\sval = \salien{\tval}{\stau}{\relax} \cr
        \sval
      & \text{otherwise}
    \end{cases}
  }
\end{mathpar}
%
\vspace{0.05cm}
\[
\begin{array}{r l l}
  \scoerce{\stunit}{}{\tunit} & = & \sunit

\\[0.3em]

  \scoerce{\sprod{\staui}{\stauii}}{\relax}{\tpair{\tvali}{\tvalii}} & = &
  \spair{\salien{\tvali}{\staui}{\relax}}{\salien{\tvalii}{\stauii}{\relax}}

\\[0.3em]

  \scoerce{\sarrow{\staui}{\stauii}}{\relax}{\texp} & = &
  \sabs{\svar{x}}{\staui}{\salien{\tapp{\texp}{\talien{\svar{x}}\staui\relax}}{\stauii}{\relax}}
\end{array}
\]
    \caption{Source semantics}
    \label{fig:sourceseman}
\end{figure}

In our example, reduction proceeds by rule {\sc{S-App}} by first
applying the unembedding coercion
$\scoerce{\source{\mathbb{Z} -> \mathbb{Z}}}{}{}$ to $\tabs{\tvar{x}}{\tvar{x}}$ to get
$\sabs{\svar{x}}{\mathbb{\source{Z}}}{\salien{\tapp{\target{(}\tabs{\tvar{x}}{\tvar{x}}\target{)}}{\talien{\svar{x}}{\mathbb{\source{Z}}}{}}}{\mathbb{\source{Z}}}{}}$,
and then applying the substitution to get
$\salien{\tapp{\target{(}\tabs{\tvar{x}}{\tvar{x}}\target{)}}{\talien{\source{0}}{\mathbb{\source{Z}}}{}}}{\mathbb{\source{Z}}}{}$.

%
%\nik{describe embeddings for pairs and then for functions, showing how the rule requires unembedding source terms and passing it to the target language.}
%
%% \jonathan{Probably need to change the rule above to account for the descrepancy
%% between the lambda and the f}
%% As such, when a function $\svar{f}$ is compiled as a plugin, one of
%% \fstar's normalizers reduces
%% $\sapp{\salien{\tvar{f}}{\sarrow\staui\stauii}\relax}{\sexp}$. The source
%% semantics \emph{force} the head of the application by applying the arrow
%% embedding above, followed by rule \textsc{App}, which gives 
%% $
%% \salien{\tapp{\tvar{f}}{\talien\sexp\staui\relax}}{\stauii}{\relax}
%% $.
%% \jonathan{Not sure about this bit.}
%% \zp{From the formalization prospective seems accurate. That about implementation?
%% Does this reflect what happens?}
%% %
%% By virtue of being call by value, the \textsc{Alien} rule for the target
%% language then gives
%% $
%% \salien{\tapp{\tvar{f}}{\tcoerce\staui\relax{\svar{v}}}}{\stauii}{\relax}
%% $.
%% In practice, reducing source terms from within the target language via rule
%% \textsc{Alien} is implemented via a callback to \fstar's normalizer.

\paragraph{\bf Target semantics and embeddings (\autoref{fig:targetseman})}
%% Although not present initially at the top-level, unembeddings of source terms in
%%  the target arise during reductions of embdeddings. Hence, we need to coerce
%%  them in the target.
The target semantics is standard
 CBV.
 %% Embeddings of source terms are not values in the target language. As such,
 %% embedded source terms are fully evaluated when passed to the target language.
 The only new rule is \textsc{T-Alien}, which first reduces the foreign
 source term using  the source semantics, and then proceeds
 to \emph{embed} the resulting value into the target
 language, using the embedding coercion $\tcoerce\stau\relax\sval$,
 as shown below. In our implementation, this rule is realized via a callback
 to the \fstar's normalizer.

 \begin{figure}
\begin{mathpar}
\inferrule*[right=T-Alien]
            {\sexp \sevals \sval}
            {\talien{\sexp}{\stau}{} \teval \tcoerce{\stau}{}{\sval}}
\and
\inferrule*[right=T-App]
            {\quad}
            {\tapp{\target{(}\tabs{\tvar{x}}{\texp}\target{)}}{\tval} \teval \texp\tsubst{\tvar{x}}{\tval}}
\end{mathpar}
%% \jonathan{I'm not sure about having expressions in these rules below, because I
%% thought that the target Alien rule forced reduction of the source expression as
%% a value? Is this meant to imply that we only need to reduce to weak head normal
%% form?}
%% \zp{Yes, we only need whnf to apply the coercions. I'm also using it to indicate that the semantics of the source
%% (i.e. the nornalizer) are CBN and therefore values are in weak head normal form. The rules below look good.
%% Only changed the last one to stress that the inner red term will be a value. }
%Coercions are defined in a similar manner.
\[
\begin{array}{r l l}
  \tcoerce{\sarrow{\staui}{\stauii}}{\relax}\sexp  & = & 
  \tabs{\tvar{x}}{\talien{\sapp\sexp{\salien{\tvar{x}}{\staui}{\relax}}}{\stauii}{\relax}}
  \\[0.3em]
  \tcoerce{\sprod{\staui}{\stauii}}{\relax}{\spair{\sexpi}{\sexpii}}  & = & 
  \tpair{\talien{\sexpi}{\staui}{\relax}}{\talien{\sexpii}{\stauii}{\relax}}
  %
  \\[0.3em] 
  \tcoerce{\staup}{\relax}{\salien{\tval}{\stau}{\relax}}  & = &  \tval
  \text{ if } \stau \equiv \staup
\end{array}
\]
     \caption{Target semantics}
     \label{fig:targetseman}
 \end{figure}

The last case in the definition of the coercion cancels
 superfluous embeddings, in case the term inside is an
 unembedding. Continuing with our example,
$\salien{\tapp{\target{(}\tabs{\tvar{x}}{\tvar{x}}\target{)}}{\talien{\source{0}}{\mathbb{\source{Z}}}{}}}{\mathbb{\source{Z}}}{}$
 reduces using {\sc{S-Alien}}, which reduces the application in its
 premise by first reducing
 $\talien{\source{0}}{\mathbb{\source{Z}}}{}$
 to
$\target{0}$ (using rule {\sc{T-Alien}}, with
 $\tcoerce{\mathbb{Z}}{\relax}{\source{0}} = \target{0}$), and
then using rule {\sc{T-App}} to get
 $\target{0}$. The rule {\sc{S-Alien}} then returns
 $\salien{\target{0}}{\mathbb{\source{Z}}}{}$
 to the \fstar normalizer, which applies the unembedding coercion to
 get $\source{0}$.
 
\iffull

\paragraph{\bf Translation from source to target}
Fig.~\ref{fig:translation} shows a few
selected cases in the translation of source to target terms.
The key rule is \textsc{Abs}, which translates a source
abstraction to a target abstraction. Rather than maintaining an explicit mapping of
source variables to target variables, the translation \emph{embeds} the target
variable $\tvar x$
in the source body of the abstraction. Later on, the \textsc{Box} rule finds the
embedding, and removes it, leaving only the intended variable in the translated
target program. As such, any closed source term is translated to a closed target
term, without any embeddings or unembeddings in it.

Rule \textsc{Var} is therefore only triggered for open variables that have not
been unembedded by \textsc{Abs} already. The formalization accounts
for this, and remains valid in the presence of open
terms. From an implementation perspective, however, we cannot in general compile
open terms to OCaml. Therefore, we restrict ourselves to compiling top-level
definitions, in which the only open variables are other top-level definitions.
Such top-level definitions have a name, which we know how to compile to OCaml.
Because of this restriction, plugin extraction works at the level of top-level
definitions, which are the only AST nodes that can be marked with the
\ls$plugin$ attribute.

\begin{figure}[h]
  \centering
\begin{mathpar}
  \inferrule*[right=Var]
             {\\}
             {\svar{x} \trans \talien{\svar{x}}{\senv\spar{\svar{x}}}{\cdot}}
  %
  \and
  %
  \inferrule*[right=Box]
             {\\}
             {\salien{\texp}{\stau}{\sdelta} \trans \texp}
  %
  \and
  %
  \inferrule*[right=Pair]
             {\source{e_i} \trans \target{e_i}}
             { \spair{\source{e_1}}{\source{e_2}} \trans \tpair{\target{e_1}}{\target{e_2}}}
  %
  \and
  %
  \inferrule*[right=Abs]
             {\sexp\ssubst{\svar{x}}{\salien{\tvar{x}}{\source{\tau}}{\cdot}} \trans \texp}
             {\sabs{\svar{x}}{\stau}{\sexp} \trans \tabs{\tvar{x}}{\texp}}
  %
  \and
  %
  \inferrule*[right=App]
             {\source{e_i} \trans \target{e_i}}
             {\sapp{\source{e_1}}{\source{e_2}} \trans \tapp{\target{e_1}}{\target{e_2}}}
  %
           \end{mathpar}
           \caption{Translation from source to target}
           \label{fig:translation}
\end{figure}
\fi

%% which is then the whole expression reduces to , which finally
%% reduces to  using rule .\guido{\textsc{S-Alien} only seems to
%% return unembedded terms? missing a coercion maybe?}
%% %% \jonathan{Re the rule for arrow types: changed it to no longer do a built-in
%% substitution. (Simpler?)}

%% The base and structural cases (e.g. pairs) are straightforward. The arrow case
%% is symmetrical to $\scoerce{\sarrow\staui\stauii}\relax\relax$. At run-time,
%% this amount to allocating a closure in the target language (OCaml) that calls
%% back into the normalizer, then embeds the result back in the target.
%% \jonathan{I feel like we should say something else in this paragraph, because
%% looking at the definition of force, it seems like we're always applying the
%% cancellation rule (may confuse the reader).}

% ----- Aseem: Begin: Text before starting tilde experiment -----

%% We start with a source language that is an intrinsically typed (Church
%% style), standard, simply-typed lambda calculus with pairs and sums and
%% a target language that is the untyped lambda calculus.
%% %
%% For clarity, we markup the the syntax using colours, using
%% \source{blue} for the source language and \target{red} for the target.
%% %
%% Two constructs are new: $\salien{\texp}{\stau}{}$, which \emph{unembeds}
%% a target computation in the source language at type $\stau$, and
%% $\talien{\sexp}{\stau}{}$, which
%% \emph{embeds} a $\stau$-typed source computation in the target language.

%% To model native plugins, consider a simply-typed lambda calculus
%% source term $\source{E(\sexp : \stau)}$, a context $\source{E}$
%% surrounding a term $\sexp : \stau$.
%% %
%% % JP: I think these two were salien as opposed to talien?
%% We translate $\sexp$ to an untyped lambda term $\target{\texp}$ (by
%% erasing type information) and unembed it in the original source
%% context as $\source{E (\salien{\target{\texp}}{\stau}{})}$.
%% %
%% We study the interactions of such hybrid terms, aiming to ensure that
%% $\source{E (\salien{\target{\texp}}{\stau}{})}$ is observationally
%% equivalent to $\source{E(\sexp : \stau)}$.

%% \iffull

%% Figure \ref{fig:translation} shows a few
%% selected cases in the translation of source to target terms.
%% The key rule is \textsc{Abs}, which translates a source
%% abstraction to a target abstraction. Rather than maintain an explicit mapping of
%% source variables to target variables, the translation \emph{embeds} the target
%% variable $\tvar x$
%% in the source body of the abstraction. Later on, the \textsc{Box} rule finds the
%% embedding, and removes it, leaving only the intended variable in the translated
%% target program. As such, any closed source term is translated to a closed target
%% term, without any embeddings or unembeddings in it.

%% Rule \textsc{Var} is therefore only triggered for open variables that have not
%% been unembedded by \textsc{Abs} already. The formalization accounts
%% for this, and remains valid in the presence of open
%% terms. From an implementation perspective, however, we cannot in general compile
%% open terms to OCaml. Therefore, we restrict ourselves to compiling top-level
%% definitions, in which the only open variables are other top-level definitions.
%% Such top-level definitions have a name, which we know how to compile to OCaml.
%% Because of this restriction, plugin extraction works at the level of top-level
%% definitions, which are the only AST nodes that can be marked with the
%% \ls$plugin$ attribute.

%% \begin{figure}[h]
%%   \centering
%% \begin{mathpar}
%%   \inferrule*[right=Var]
%%              {\\}
%%              {\svar{x} \trans \talien{\svar{x}}{\senv\spar{\svar{x}}}{\cdot}}
%%   %
%%   \and
%%   %
%%   \inferrule*[right=Box]
%%              {\\}
%%              {\salien{\texp}{\stau}{\sdelta} \trans \texp}
%%   %
%%   \and
%%   %
%%   \inferrule*[right=Pair]
%%              {\source{e_i} \trans \target{e_i}}
%%              { \spair{\source{e_1}}{\source{e_2}} \trans \tpair{\target{e_1}}{\target{e_2}}}
%%   %
%%   \and
%%   %
%%   \inferrule*[right=Abs]
%%              {\sexp\ssubst{\svar{x}}{\salien{\tvar{x}}{\source{\tau}}{\cdot}} \trans \texp}
%%              {\sabs{\svar{x}}{\stau}{\sexp} \trans \tabs{\tvar{x}}{\texp}}
%%   %
%%   \and
%%   %
%%   \inferrule*[right=App]
%%              {\source{e_i} \trans \target{e_i}}
%%              {\sapp{\source{e_1}}{\source{e_2}} \trans \tapp{\target{e_1}}{\target{e_2}}}
%%   %
%%            \end{mathpar}
%%            \caption{Translation from source to target}
%%            \label{fig:translation}
%% \end{figure}

%% \fi

%% %% The main characteristic of the rules is that whenever a target term blocks
%% %% reduction in the source language, we $\sforce{}$ the term into a source value of
%% %% the appropriate type, using the unembedding function $\scoerce\stau\relax\texp$.
%% %% When $\sforce{\svar}$ is applied at function type, this involves passing a
%% %% source term as an argument to a target language function,
%% %% after \emph{unembedding} said function as a source term, i.e., although not
%% %% present initially, unembeddings of target terms in the source arise as during
%% %% reductions.
%% %% \guido{First use of $\scoerce\stau\relax\texp$, confused. Last sentence
%% %% seems explained backwards to me, hard to follow.}

%% %% \jonathan{I'm presenting both reduction semantics in this paragraph, then
%% %% jumping to the description of the two coercion functions.}

%% %% \nik{I suggest presenting embeddings first, since
%% %% this is the direction with which things ``start'' in the plugin
%% %% scenario.}

%% %% \zp{What about this structure:
%% %% \begin{itemize}
%% %% \item Embeddings : syntax + semantics (reduction and coersions)
%% %% \item Unembeddings : ditto
%% %% \item Other (non-standard) rules of the semantics + force
%% %% \end{itemize}
%% %% }

%% % \nik{remove $\delta$ at first and the protection function at first.}
%% %
%% %% \begin{mathpar}
%% %%  \inferrule*[right=Alien]
%% %%             {\sexp \sevals \sval}
%% %%             {\talien{\sexp}{\stau}{\relax} \teval \tcoerce{\stau}{\relax}{\sval}}
%% %% \end{mathpar}

%% \paragraph{Embeddings}
%% To define the semantics of the two languages we need a way to interpret alien terms
%% in both contexts. 
%%  The source reduces an alien term $\salien{\texp}{\stau}{\relax}$
%% by invoking the semantics of the target language to obtain a target value (rule {\sc Alien}).
%% Hence, the term $\salien{\tval}{\stau}{\relax}$ is a value in the source \emph{iff}
%% $\tval$ is a value in the target language. Whenever a reduction step is blocked
%% because of an alien value, the semantics will coerce this value to a source value by a
%% call to the meta-function $\sforce{}$.
%% %
%% \begin{mathpar}
%%   \inferrule*[right=Alien]
%%              {\texp \tevals \tval}
%%              {\salien{\texp}{\stau}{\sdelta} \seval \scoerce{\stau}{\sdelta}{\tval}}
%%   %
%% %  \and
%% %  \inferrule*[right=Proj]
%% %             {\sforce{\sval}=\spair{\sexp_{\source{1}}}{\sexp_{\source{2}}}}
%% %             {\sval.\source{i} \seval \sexp_{\source{i}}}
%%  %
%%   \and
%%   %
%%   \inferrule*[right=App]
%%              {\sforce{\sval}=\sabs{\svar{x}}{\stau}{\sexp}}
%%              {\sapp{\sval}{\sexpii} \seval \sexp \ssubst{\svar{x}}{\sexpii}}
%%  %
%%  \and
%%   \inferrule*{}{
%%     \sforce{\sval} =
%%     \begin{cases}
%%       \scoerce{\stau}{\relax}{\tval}
%%       & \text{if }\sval = \salien{\tval}{\stau}{\relax} \cr
%%         \sval
%%       & \text{otherwise}
%%     \end{cases}
%%   }
%% \end{mathpar}
%% $\sforce{}$ examines its argument and invokes the the embedding meta-function
%% $\scoerce\stau\relax\tval$ whenever it hits an alien value. The embedding
%% meta-function $\scoerce\stau\relax\tval$ takes a target \emph{value} (the target
%% is CBV) and coerces it suitably as a source \emph{value}.
%% \zp{The result of the coercion will always be a value in the simply typed world.
%% It can be an expression when we add poly types.}
%% For structural types, this amounts to lazily coercing the head constructor of
%% the term (below). For function types, a source-level closure is allocated, which
%% first \emph{embeds} its argument in the target, reduces a target application,
%% then \emph{unembeds} the result back in the source.
%% %
%% \[
%% \begin{array}{c c}
%%   \scoerce{\sprod{\staui}{\stauii}}{\relax}{\tpair{\tvali}{\tvalii}} =
%%   \spair{\salien{\tvali}{\staui}{\relax}}{\salien{\tvalii}{\stauii}{\relax}}

%%   \qquad

%% &

%%   \qquad

%%   \scoerce{\sarrow{\staui}{\stauii}}{\relax}{\tabs{\tvar{x}}{\texp}} =
%%   \sabs{\svar{x}}{\staui}{\salien{\tapp{(\tabs{\tvar{x}}{\texp})}{\talien{\svar{x}}\staui\relax}}{\stauii}{\relax}}
%% \end{array}
%% \]
%% %
%% %\nik{describe embeddings for pairs and then for functions, showing how the rule requires unembedding source terms and passing it to the target language.}
%% %
%% \jonathan{Probably need to change the rule above to account for the descrepancy
%% between the lambda and the f}
%% As such, when a function $\svar{f}$ is compiled as a plugin, one of
%% \fstar's normalizers reduces
%% $\sapp{\salien{\tvar{f}}{\sarrow\staui\stauii}\relax}{\sexp}$. The source
%% semantics \emph{force} the head of the application by applying the arrow
%% embedding above, followed by rule \textsc{App}, which gives 
%% $
%% \salien{\tapp{\tvar{f}}{\talien\sexp\staui\relax}}{\stauii}{\relax}
%% $.
%% \jonathan{Not sure about this bit.}
%% \zp{From the formalization prospective seems accurate. That about implementation?
%% Does this reflect what happens?}
%% %
%% By virtue of being call by value, the \textsc{Alien} rule for the target
%% language then gives
%% $
%% \salien{\tapp{\tvar{f}}{\tcoerce\staui\relax{\svar{v}}}}{\stauii}{\relax}
%% $.

%% \paragraph{Unembeddings}
%% Although not present initially at the top-level, unembeddings of source terms in
%%  the target arise during reductions of embdeddings. Hence, we need to coerce
%%  them in the target. The reduction semantics for the target language is standard
%%  CBV. Embeddings of source terms are not values in the target language. As such,
%%  embedded source terms are fully evaluated when passed to the target language.
%%  The only new rule is \textsc{Alien}, which first reduces the foreign term using
%%  the source language semantics, then proceeds to \emph{embed} it into the target
%%  language, using the target \emph{unembedding} $\tcoerce\stau\relax\sval$.
%% \begin{mathpar}
%% \inferrule*[right=Alien]
%%             {\sexp \sevals \sval}
%%             {\talien{\sexp}{\stau}{\sdelta} \teval \tcoerce{\stau}{\sdelta}{\sval}}
%% \end{mathpar}
%% %% \jonathan{I'm not sure about having expressions in these rules below, because I
%% %% thought that the target Alien rule forced reduction of the source expression as
%% %% a value? Is this meant to imply that we only need to reduce to weak head normal
%% %% form?}
%% %% \zp{Yes, we only need whnf to apply the coercions. I'm also using it to indicate that the semantics of the source
%% %% (i.e. the nornalizer) are CBN and therefore values are in weak head normal form. The rules below look good.
%% %% Only changed the last one to stress that the inner red term will be a value. }
%% Coercions are defined in a similar manner.
%% \[
%% \begin{array}{c c c}
%%   \tcoerce{\sarrow{\staui}{\stauii}}{\relax}\sexp  = 
%%   \tabs{\tvar{x}}{\talien{\sapp\sexp{\salien{\tvar{x}}{\staui}{\relax}}}{\stauii}{\relax}}
%%   & 
%%   \quad
%%   \tcoerce{\sprod{\staui}{\stauii}}{\relax}{\spair{\sexpi}{\sexpii}}  = 
%%   \tpair{\talien{\sexpi}{\staui}{\relax}}{\talien{\sexpii}{\stauii}{\relax}}
%%   %
%%   \quad
%%   & 
%%   \tcoerce{\staup}{\relax}{\salien{\tval}{\stau}{\relax}}  =  \tval
%%   \text{ if } \stau \equiv \staup
%% \end{array}
%% \]
%% \jonathan{Re the rule for arrow types: changed it to no longer do a built-in
%% substitution. (Simpler?)}

%% The base and structural cases (e.g. pairs) are straightforward. The arrow case
%% is symmetrical to $\scoerce{\sarrow\staui\stauii}\relax\relax$. At run-time,
%% this amount to allocating a closure in the target language (OCaml) that calls
%% back into the normalizer, then embeds the result back in the target. The last
%% case covers the cancellation of embeddings and un-embeddings, and avoids
%% generating superfluous unembeddings. 
%% \jonathan{I feel like we should say something else in this paragraph, because
%% looking at the definition of force, it seems like we're always applying the
%% cancellation rule (may confuse the reader).}

% ----- Aseem: End: Text before starting tilde experiment -----

\subsection*{Part 2: Adding ML-style Polymorphism} %\nik{Budget: 1/2 page}}
\label{sec:parametricity}

%% We extend the source language with ML-style polymorphism and show how
%% we use it to support a limited but useful form of open reduction.
%%  This
%% facilitates the use of polymorphic libraries (e.g., classic library
%% functions on lists) as native plugins.

The key observation in adding ML-style polymorphism is as
follows. Since we only consider target terms
compiled from well-typed source terms, by virtue of
parametricity, a compiled polymorphic function must treat its
argument abstractly. As a result, the normalizers can pass such
arguments \emph{as-is} without applying any embedding coercions---passing
an open term is simply a subcase without further
difficulty. Thus, we can leverage polymorphism to support such limited
but useful open reductions.

In the formal model, we add an \emph{opaque} construct to the target
language, written as
$\topaque\sexp$, which denotes an embedding of $\sexp$ at some
abstract type $\source{A}$ (as mentioned above, coercion to opaque is
a no-op in the implementation\guido{not mentioned}).
%
The coercion functions are extended
with rules to introduce and eliminate opaque terms (the $\sdelta$
superscript carries type substitutions as a technical device), e.g.:
\[
\begin{array}{c c c}
  \tcoerce{\svar{A}}{\sdelta}{\sexp}  = \topaque{\sexp}
  \quad & \quad
  \scoerce{\svar{A}}{\sdelta}{\topaque{\sexp}} = \sexp
  \quad
  &
  \quad
  \tcoerce{\sprod{\staui}{\stauii}}{\sdelta}{\spair{\sexpi}{\sexpii}} =
  \tpair{\prot{\source{e_1}}{\source{\tau_1}}{\sdelta}}{\prot{\source{e_2}}{\source{\tau_2}}{\sdelta}}
\end{array}
\]
where $\prot{\sexp}{\stau}{\sdelta}$ is $\topaque{\sexp}$ when
$\stau = \source{A}$, or a usual type-directed embedding otherwise. The source
semantics is 
extended with a rule for type applications, while target
semantics, being untyped, remains as-is.

Consider an application of the polymorphic identity function to a
variable $\svar{y}$:
$\sapp{\stapp{(\stabs{\svar{A}}{\sabs{\svar{x}}{\svar{A}}{\svar{x}}})}{\source{\mathbb{Z}}}}{\svar{y}}$.
As before, we would like to compile the function to native code and
then reduce the applications---except this time it is an open term.
%
Translation of the function and its embedding in the source yields
$\sapp{\stapp{\salien{\tabs{\tvar{x}}{\tvar{x}}}{\sall{\svar{A}}{\sarrow{\svar{A}}{\svar{A}}}}{\sesubstemp}}{\source{\mathbb{Z}}}}{\svar{y}}$.
%
We now apply the unembedding coercion
$\scoerce{\sall{\svar{A}}{\sarrow{\svar{A}}{\svar{A}}}}{\sesubstemp}{\tabs{\tvar{x}}{\tvar{x}}}$
which returns \linebreak a source type abstraction
$\stabs{\svar{A}}{\scoerce{\sarrow{\svar{A}}{\svar{A}}}{\sesubstemp}{\tabs{\tvar{x}}{\tvar{x}}}}$.
%
Reducing the arrow coercion next, we get
$\stabs{\svar{A}}{\sabs{\svar{x}}{\svar{A}}{\salien{\tapp{(\tabs{\tvar{x}}{\tvar{x}})}{\prot{\svar{x}}{\svar{A}}{\sesubstemp}}}{\svar{A}}{\sesubstemp}}}$,
%
where $\prot{\svar{x}}{\svar{A}}{\sesubstemp} = \topaque{\svar{x}}$, as mentioned above.
This is the key idea---when the original application reaches the beta-reduction with
argument $\svar{y}$ (after the type application is reduced), $\svar{y}$ is simply substituted
in this opaque construct, which is then returned back to the \fstar normalizer as
$\salien{\topaque{\svar{y}}}{\svar{A}}{\sesubstsingl{\svar{A}}{\source{\mathbb{Z}}}}$.
%
The normalizer can then apply the identity unembedding and get $\svar{y}$.

This ability to reduce open terms relying on polymorphism allows us to
evaluate expressions like
$\sapp{\sapp{\stapp{\stapp{\salien{\target{\text{\ls{List.map}}}}{\sall{\svar{A}}{\sall{\svar{B}}{\sarrow{\sarrow{(\sarrow{\svar{A}}{\svar{B}})}{\source{\text{\ls{list}}} \svar{A}}}{\source{\text{\ls{list}}} \svar{B}}}}}{\sesubstemp}}{\source{\mathbb{Z}}}}{\source{\mathbb{Z}}}}{\source{\text{\ls{f}}}}}{\source{\text{\ls{[0; ...; 10}$^6$\ls{]}}}}$,
with \ls$f$ a free variable,
using native datastructure implementations (\ls{List} in this case), 
which is much faster than using the normalizers.
% (see \autoref{sec:registers}).

\subsection*{Part 3: Full Details on our Multi-language Interoperability Model}

%% \aseem{TODO: add an example of applying the polymorphic identity
%% function to a variable. Add the property we hope for.}

%% prevents further reduction or coercion (it is a value in
%% the target language). An opaque construct preserves a source term at an abstract type.
%% %
%% We define the abstraction-protecting embedding $\prot\sexp\stau\sdelta$ to be
%% either a regular $\talien\sexp\stau\sdelta$ as before, or an opaque
%% $\topaque\sexp$ if $\sdelta(\stau) = \svar{A}$. Rules that previously generated
%% $\talien\sexp\stau\sdelta$ now generate $\prot\sexp\stau\sdelta$, e.g.
%% $
%%   \tcoerce{\sarrow{\staui}{\stauii}}{\sdelta}\sexp  = 
%%   \tabs{\tvar{x}}{\prot{\sapp\sexp{\salien{\tvar{x}}{\staui}{\sdelta}}}{\stauii}{\sdelta}}
%% $


%% In the formal model, we add a new \em

%% In a polymorphic setting, abstract type variables appear in the language.
%% To ensure soundness, coercions in the target language must stop at such abstract
%% type variables. (Coercing underneath an abstraction would reveal the underlying
%% type to the source language, and return an ill-typed term.) We add a type
%% substitution $\sdelta$ in:
%% $\tcoerce\stau\sdelta\sval$, 
%% $\scoerce\stau\sdelta\tval$, 
%% $\talien\sval\stau\sdelta$, 
%% $\salien\tval\stau\sdelta$. Upon entering a type scheme
%% $\sall{\svar{A}}{\ssig}$, $\sdelta$ is extended and becomes
%% ${\sesubstext{\sdelta}{\svar{A}}{\stau}}$.

%% We extend the target language with an \emph{opaque} construct, denoted
%% $\topaque\sexp$, which prevents further reduction or coercion (it is a value in
%% the target language). An opaque construct preserves a source term at an abstract type.
%% %
%% We define the abstraction-protecting embedding $\prot\sexp\stau\sdelta$ to be
%% either a regular $\talien\sexp\stau\sdelta$ as before, or an opaque
%% $\topaque\sexp$ if $\sdelta(\stau) = \svar{A}$. Rules that previously generated
%% $\talien\sexp\stau\sdelta$ now generate $\prot\sexp\stau\sdelta$, e.g.
%% $
%%   \tcoerce{\sarrow{\staui}{\stauii}}{\sdelta}\sexp  = 
%%   \tabs{\tvar{x}}{\prot{\sapp\sexp{\salien{\tvar{x}}{\staui}{\sdelta}}}{\stauii}{\sdelta}}
%% $

%% Crucially, the coercion functions are extended with rules to deal with opaque
%% terms.
%% \[
%% \begin{array}{c c}
%%   \tcoerce{\svar{A}}{\sdelta}{\sexp}  = \topaque{\sexp}
%%   \quad & \quad
%%   \scoerce{\svar{A}}{\sdelta}{\topaque{\sexp}} = \sexp
%% \end{array}
%% \]
%% In essence, these rules embody the fact that a source function that is
%% parametric over type $\svar{A}$ cannot examine a value at this abstract type,
%% meaning that there actually is no need to unembed or embed such at value at the
%% language boundary: the target program cannot do anything with it, except return
%% it.

%% One consequence of this is that by using polymorphic functions, the programmer
%% can devise a library of natively-compiled plugins that are capable of dealing
%% with open terms. The polymorphic versions are compiled natively, but are used
%% from the source at concrete types that may contain open variables. No embeddings
%% are generated, meaning that the target code is oblivious that the values it is
%% passing around are opaque, open source terms.

%% \jonathan{Is this what we want to say?}
%% Naturally, \fstar goes beyond ML-style polymorphism. For the time being, we
%% decline to compile plugins that, at extraction-time, cannot be re-typed with ML
%% polymorphism.
%% \aseem{We mentioned it before in the limitations of native plugins.}


%\newcommand\path{interop_lazy/}%

\begin{wrapfigure}[htb]{l}{\textwidth}%
    %% \centering
    %% \begin{minipage}{0.54\textwidth}\label{fig:srcexp}
      \centering
      \begin{framed}
        \[\begin{array}{@{}l@{~~}l@{~}r@{}l@{}}
\mbox{Types} & \stau & \bnfdef & \stunit \bnfalt \svar{A}
      \bnfalt \sarrow{\stau_{\source{1}}}{\stau_{\source{2}}}
      \bnfalt \sprod{\stau_{\source{1}}}{\stau_{\source{2}}} \bnfalt \ssum{\stau_{\source{1}}}{\stau_{\source{2}}} \\
& & & \\

\mbox{Type Schemes} & \ssig & \bnfdef & \sall{\svar{A}}{\ssig} \bnfalt \stau \\
& & & \\


\mbox{Environments} & \senv & \bnfdef & \senvemp \bnfalt \senvext{\senv}{\svar{x}}{\stau} \\
& \stenv & \bnfdef & \senvemp \bnfalt \senvtext{\stenv}{\svar{A}} \\
& & & \\
\mbox{Substitutions} & \sesubst & \bnfdef & \sesubstemp \bnfalt \sesubstext{\sesubst}{\svar{A}}{\stau} \\
& & & \\
\mbox{Expressions} & \sexp & \bnfdef & \sunit \bnfalt \svar{x} \bnfalt \sabs{\svar{x}}{\stau}{\sexp} \bnfalt \sapp{\sexp_{\source{1}}}{\sexp_{\source{2}}} \\
& & \bnfalt & \stabs{A}{e} \bnfalt \stapp{\sexp}{\stau} \\
& & \bnfalt & \spair{\sexp_{\source{1}}}{\sexp_{\source{2}}} \bnfalt \sfst{\sexp} \bnfalt \ssnd{\sexp} \\
& & \bnfalt & \sinl{\sexp} \bnfalt \sinr{\sexp} \\
& & \bnfalt & \scasex{\sexp}{\stau_{\source{1}}}{\sexp_{\source{1}}}{\stau_{\source{2}}}{\sexp_{\source{2}}} \\
& & \bnfalt & \salien{\texp}{\stau}{\sdelta} \\
% \bnfalt \sopaque{\texp}
\end{array}\]

      \end{framed}
      \caption{Source Language}
      \label{fig:sourcelanglong}
    %% \end{minipage}\hfill
    %% \begin{minipage}{0.44\textwidth}\label{fig:trgexp}
\end{wrapfigure}


\begin{wrapfigure}[htb]{l}{\textwidth}%
      \centering
      \begin{framed}
        \[\begin{array}{@{}l@{~~}l@{~}r@{}l@{}}
\mbox{Expressions} & \texp & \bnfdef & \tunit \bnfalt \tvar{x} \bnfalt \tabs{\tvar{x}}{\texp} \bnfalt \tapp{\texp_{\target{1}}}{\texp_{\target{2}}} \\
& & \bnfalt & \tpair{\texp_{\target{1}}}{\texp_{\target{2}}} \bnfalt \tfst{\texp} \bnfalt \tsnd{\texp} \\
& & \bnfalt & \tinl{\texp} \bnfalt \tinr{\texp} \\
& & \bnfalt & \tcasex{\texp}{\texp_{\target{1}}}{\texp_{\target{2}}} \\
& & \bnfalt & \talien{\sexp}{\stau}{\sdelta} \bnfalt
\topaque{\sexp}\\
& & & \\
\mbox{Environments} & \tenv & \bnfdef & \tenvemp \bnfalt \tenvext{\tenv}{\tvar{x}} \\

\end{array}\]

      \end{framed}
      \caption{Target Language}
%    \end{minipage}
\end{wrapfigure}

\begin{wrapfigure}[htb]{l}{\textwidth}%
  \label{fig:srctyp}
  \fbox{$\iswf{\stenv}{\stau}$}
  \begin{framed}
    \begin{mathpar}
  \inferrule*[right=TyVar]
  {\\}
  {\iswf{\senvtext{\stenv}{\svar{A}}}{\svar{A}}}
  %
  \and
  %
  \inferrule*[right=Scheme]
  {\iswf{\senvtext{\stenv}{\svar{A}}}{\sexp}}
  {\iswf{\stenv}{\sall{\svar{A}}{\stau}}}
  %
  \and
  %
  \inferrule*[right=Env]
  {\forall ~x,~\iswf{\stenv}{\senv(x)}}
  {\iswf{\stenv}{\senv}}
  %% \inferrule*[right=Opaque]
  %% {\\}
  %% {\hastyp{\senv}{\sopaque{\texp}}{\sttop}}
\end{mathpar}

  \end{framed}
  \caption{Type Well-formedness}
\end{wrapfigure}

\begin{wrapfigure}[htb]{l}{\textwidth}%
  \label{fig:srctyp}
  \fbox{$\iswf{\stenv}{\stau}$}
  \begin{framed}
    \begin{mathpar}
  \inferrule*[right=Empty]
  {\\}
  {\iswf{\cdot}{\cdot}}
  %
  \and
  %
  \inferrule*[right=Cons]
             {
               \iswf{\stenv}{\sdelta} \\
               \iswf{\stenv}{\stau} \\
             }
             {
               \iswf{\senvtext{\stenv}{\svar{A}}}
                    {\sesubstext{\sdelta}{\svar{A}}{\stau}}
             }
\end{mathpar}

  \end{framed}
  \caption{Type Substitution Well-formedness}
\end{wrapfigure}


\begin{wrapfigure}[htb]{l}{\textwidth}%
  \label{fig:srctyp}
  \fbox{$\hastyp{\stenv}{\senv}{\sexp}{\stau}$}
  \begin{framed}
    \begin{mathpar}
\inferrule*[right=Alien]
  {\isok{\stenv}{\senv}{\texp}}
  {\hastyp{\stenv}{\senv}{\salien{\texp}{\stau}{\sesubst}}{\sesubst \stau}}
  %% %
  %% \and
  %% %
  %% \inferrule*[right=Opaque]
  %% {\\}
  %% {\hastyp{\senv}{\sopaque{\texp}}{\sttop}}
\end{mathpar}

  \end{framed}
  \caption{Source Typing}
\end{wrapfigure}

%% \begin{wrapfigure}[htb]{l}{\textwidth}%
%%   \label{fig:srctyp}
%%   \fbox{$\isok{\stenv}{\senv}{\texp}$ (This might not be needed)}
%%   \begin{framed}
%%     \input{trg_typing}
%%   \end{framed}
%%   \caption{Target well-formedness}
%% \end{wrapfigure}

\begin{wrapfigure}[htb]{l}{\textwidth}%
  \label{fig:srceval}
  \fbox{$\sexp \seval \sexpp$}
  \begin{framed}
    \begin{mathpar}
  \inferrule*[right=Ctx]
             {\sexp \seval \sexpp}
             {\scon{\sexp} \seval \scon{\sexpp}}
  %
  \and
  \inferrule*[right=App]
             {\sforce{\sval}=\sabs{\svar{x}}{\stau}{\sexp}}
             {\sapp{\sval}{\sexpii} \seval \sexp \ssubst{\svar{x}}{\sexpii}}

  %
  \and
  %
  \inferrule*[right=TApp]
             {\sforce{\sval}=\stabs{\svar{A}}{\sexpp}}
             { \stapp{\sval}{\stau} \seval
               \sexpp \ssubst{\svar{A}}{\stau}}
  %
  \and
  %
  \inferrule*[right=CaseL]
             {\sforce{\sval} = \sinl{\sexpp}}
             {\scasex{\sval}{\stau_{\source{1}}}{\sexp_{\source{1}}}{\stau_{\source{2}}}{\sexp_{\source{2}}} \seval \sexp_{\source{1}} \ssubst{\svar{x}}{\sexpp}}
  %
  \and
  %
    \inferrule*[right=CaseR]
             {\sforce{\sval} = \sinr{\sexpp}}
             {\scasex{\sval}{\stau_{\source{1}}}{\sexp_{\source{1}}}{\stau_{\source{2}}}{\sexp_{\source{2}}} \seval \sexp_{\source{2}} \ssubst{\svar{x}}{\sexpp}}
  %
  \and
  %
  \inferrule*[right=Proj]
             {\sforce{\sval}=\spair{\sexp_{\source{1}}}{\sexp_{\source{2}}}}
             {\sval.\source{i} \seval \sexp_{\source{i}}}
 %
  \and
  %
  \inferrule*[right=Alien]
             {\texp \tevals \tval}
             {\salien{\texp}{\stau}{\sdelta} \seval \salien{\sval}{\stau}{\sdelta}}
   %
  \and
  %
  %% \inferrule*[right=ProjAlien]
  %%            {\\}
  %%            {\salien{\tpair{\texp_{\target{1}}}{\texp_{\target{2}}}}{\sprod{\stau_{\source{1}}}{\stau_{\source{2}}}}{\sdelta}
  %%              \seval
  %%              \salien{\texp_{\target{1}}}{\stau_{\source{1}}}{\sdelta} }
  %% %
  %% \and
  %% %
  %% \inferrule*[right=CaseLAlien]
  %%            {\\}
  %%            {\scasex{\salien{\tinl{\texp}}{\ssum{\stau_{\source{1}}}{\stau_{\source{2}}}}{\sdelta}}{\stau_{\source{1}}}{\sexp_{\source{1}}}{\stau_{\source{2}}}{\sexp_{\source{2}}}
  %%              \seval
  %%              \sexp_{\source{1}} \ssubst{\svar{x}}{\salien{\texp}{\stau_{\source{1}}}{\sdelta}}}
  %% %
  %% \and
  %% %
  %% \inferrule*[right=AppAlien]
  %%            {\\}
  %%            {\sapp{\salien{\tabs{\tvar{x}}{\texp}}{\sarrow{\stau_{\source{1}}}{\stau_{\source{2}}}}{\sdelta}}
  %%              {\sexp}
  %%              \seval
  %%              %% \sabs{\svar{x}}{\sdelta\stau_{\source{1}}}
  %%                   {\kap{\stau_{\source{1}}}{\sdelta}{\sexp}
  %%                     {\mfun{\texpp}{\salien{\texp\ssubst{\tvar{x}}{\texpp}}{\stau_{\source{2}}}{\sdelta}}}}
  %%            }
  %% %
  %% \and
  %% %
  %% \inferrule*[right=TAppAlien]
  %%            {\\}
  %%            {\stapp{\salien{\texp}{\sall{\svar{A}}{\stau}}{\sdelta}}{\staup}
  %%              \seval
  %%              \salien{\texp}{\stau}{\sesubstext{\sdelta}{\svar{A}}{\staup}}
  %%            }
  %% %
  %% \and
  %% %
  %% \inferrule*[right=AlienOpaque]
  %%            {\stau \not = \sttop}
  %%            {\salien{\topaque{\sexp}}{\stau}{\sdelta}
  %%              \seval
  %%              \sexp
  %%            }
  %% %
  %% \and
  %% %
  %% \inferrule*[right=AlienCancel]
  %%            {\\}
  %%            {\salien{\talien{\sexp}{\staup}{\sdeltap}}{\stau}{\sdelta}
  %%              \seval
  %%              \sexp}
\end{mathpar}

    where
    \[\begin{array}{r l l}
  \sforce{\sval} =
  \begin{cases}
    \scoerce{\stau}{\sdelta}{\tval}
    & \text{if }\sexp = \salien{\tval}{\stau}{\sdelta} \text{ and }
      \sexp \sevals \sval
\\
      \sval
    & \text{otherwise}
  \end{cases}
\end{array}\]

  \end{framed}
  \caption{Evaluation in the Source Language}
\end{wrapfigure}

\begin{wrapfigure}[htb]{l}{\textwidth}%
  \label{fig:trgeval}
  \fbox{$\texp \teval \texpp$}
  \begin{framed}
    \begin{mathpar}
    \inferrule*[right=Ctx]
             {\texp \teval \texpp}
             {\tcon{\texp} \teval \tcon{\texpp}}
  %
  \and
  %
  \inferrule*[right=App]
             {\\}
             {\tapp{(\tabs{\tvar{x}}{\texp})}{\tval} \teval \texp \tsubst{\tvar{x}}{\tval}}

  %
  \and
  %
  \inferrule*[right=CaseL]
             {\\}
             {\tcasex{\tinl{\tval}}{\texp_{\target{1}}}{\texp_{\target{2}}} \teval \texp_{\target{1}} \tsubst{\tvar{x}}{\tval}}
  %
  \and
  %
  \inferrule*[right=CaseL]
             {\\}
             {\tcasex{\tinr{\tval}}{\texp_{\target{1}}}{\texp_{\target{2}}} \teval \texp_{\target{2}} \tsubst{\tvar{x}}{\tval}}
  %
  \and
  %
  \inferrule*[right=Proj]
             {\\}
             {\tpair{\tval_{\target{1}}}{\tval_{\target{2}}}\target{.i} \teval \tval_{\target{i}}}
 %
 \and
 %
 \inferrule*[right=Alien]
            {\sexp \sevals \sval}
            {\talien{\sexp}{\stau}{\sdelta} \teval \tcoerce{\stau}{\sdelta}{\sval}}
 %
 \and
 %%  %
 %%  \inferrule*[right=ProjAlien]
 %%             {\\}
 %%             {\salien{\tpair{\texp_{\target{1}}}{\texp_{\target{2}}}}{\sprod{\stau_{\source{1}}}{\stau_{\source{2}}}}{\sdelta}
 %%               \seval
 %%               \salien{\texp_{\target{1}}}{\stau_{\source{1}}}{\sdelta} }
 %%  %
 %%  \and
 %%  %
 %%  \inferrule*[right=CaseLAlien]
 %%             {\\}
 %%             {\scasex{\salien{\tinl{\texp}}{\ssum{\stau_{\source{1}}}{\stau_{\source{2}}}}{\sdelta}}{\stau_{\source{1}}}{\sexp_{\source{1}}}{\stau_{\source{2}}}{\sexp_{\source{2}}}
 %%               \seval
 %%               \sexp_{\source{1}} \ssubst{\svar{x}}{\salien{\texp}{\stau_{\source{1}}}{\sdelta}}}
 %%  %
 %%  \and
 %%  %
 %%  \inferrule*[right=AppAlien]
 %%             {\\}
 %%             {\sapp{\salien{\tabs{\tvar{x}}{\texp}}{\sarrow{\stau_{\source{1}}}{\stau_{\source{2}}}}{\sdelta}}
 %%               {\sexp}
 %%               \seval
 %%               \sabs{\svar{x}}{\sdelta\stau_{\source{1}}}
 %%                    {\kap{\stau_{\source{1}}}{\sdelta}{\svar{x}}
 %%                      {\mfun{\texpp}{\salien{\texp\ssubst{\tvar{x}}{\texpp}}{\stau_{\source{2}}}{\sdelta}}}}
 %%             }
 %%  %
 %%  \and
 %%  %
 %%  \inferrule*[right=TAppAlien]
 %%             {\\}
 %%             {\stapp{\salien{\texp}{\sall{\svar{A}}{\stau}}{\sdelta}}{\staup}
 %%               \seval
 %%               \salien{\texp}{\stau}{\sesubstext{\sdelta}{\svar{A}}{\staup}}
 %%             }
 %%  %
 %%  \and
 %%  %
 %%    \inferrule*[right=AlienOpaque]
 %%             {\\}
 %%             {\salien{\topaque{\sexp}}{\stau}{\sdelta}
 %%               \seval
 %%               \sexp
 %%             }
 %%               %
 %%  \and
 %%  %
  %% \inferrule*[right=AlienCancel]
  %%            {\\}
  %%            {\talien{\salien{\texp}{\staup}{\sdeltap}}{\stau}{\sdelta}
  %%              \teval
  %%              \texp}
\end{mathpar}

  \end{framed}
  \caption{Evaluation in the Target Language}
\end{wrapfigure}

\medskip

\begin{wrapfigure}[htb]{l}{\textwidth}%
    \centering
    \begin{minipage}[t]{0.45\textwidth}\label{fig:stypsubst}
      \centering
      \begin{framed}
        \[
\begin{array}{r l l}
  \sunit \ssubst{\svar{A}}{\stau} & = & \sunit
  \\[0.3em]
  %
  \sabs{\svar{x}}{\staup}{\sexp} \ssubst{\svar{A}}{\stau} & = &
  \sabs{\svar{x}}{\staup\ssubst{\svar{A}}{\stau}}{\sexp\ssubst{\svar{A}}{\stau}}
  \\[0.3em]
  %
  \vdots 
  \\[0.3em]
  \salien{\texp}{\staup}{\sdelta}\ssubst{\svar{A}}{\stau}  & = &
  \salien{\texp\ssubst{\svar{A}}{\stau}}{\staup}
 {\sesubstext{\sdelta}{\svar{A}}{\stau}}
  \\[0.3em]
  %
  %
  %
  %% \scoerce{\sttop}{\sdelta}{\salien{\tval}{\staup}{\sdeltap}} & = &
  %% \tval
  %% \\[0.3em]
\end{array}
\]

      \end{framed}
      \caption{Type Substitution in Source Terms}
    \end{minipage}\hfill
    \begin{minipage}[t]{0.45\textwidth}\label{fig:ttypsubst}
      \centering
      \begin{framed}
       \[
\begin{array}{r l l}
  \tunit \ssubst{\svar{A}}{\stau} & = & \tunit
  \\[0.3em]
  %
  \tabs{\tvar{x}}{\texp} \ssubst{\svar{A}}{\stau} & = &
  \tabs{\tvar{x}}{\texp\ssubst{\svar{A}}{\stau}}
  \\[0.3em]
  %
  \vdots 
  \\[0.3em]
  \talien{\sexp}{\staup}{\sdelta}\ssubst{\svar{A}}{\stau}  & = &
  \talien{\sexp\ssubst{\svar{A}}{\stau}}{\staup}{\sesubstext{\sdelta}{\svar{A}}{\stau}}
  \\[0.3em]
  \topaque{\sval}  & = & 
  \topaque{\sval \ssubst{\svar{A}}{\stau}}
  \\[0.3em]
\end{array}
\]

      \end{framed}
      \caption{Type Substitution in Target Terms}
    \end{minipage}
\end{wrapfigure}


\medskip

\begin{wrapfigure}[htb]{L}{\textwidth}%
  \label{fig:trans}
  \fbox{$\sexp \trans \texp$} \\
  $\senv$ is a global parameter 
  \begin{framed}
    \begin{mathpar}
  \inferrule*[right=Unit]
             {\\}
             {\sunit \trans \tunit}
  %
  \and
  \inferrule*[right=Var]
             {\\}
             {\svar{x} \trans \talien{\svar{x}}{\senv\spar{\svar{x}}}{\cdot}}

  %
  \and
  %
  \inferrule*[right=Box]
             {\\}
             {\salien{\texp}{\stau}{\sdelta} \trans \texp}
  %
  \and
  %
  \inferrule*[right=Pair]
             {\source{e_i} \trans \target{e_i}}
             { \spair{\source{e_1}}{\source{e_2}} \trans \tpair{\target{e_1}}{\target{e_2}}}
  %
  \and
  %
  \inferrule*[right=Proj]
             {\sexp \trans \texp}
             {\sproji{\sexp} \trans \tproji{\texp}}
  %
  \and
  %
  \inferrule*[right=Inl]
             {\sexp \trans \texp}
             {\sinl{\sexp} \trans \tinl{\texp}}
  %
  \and
  %
  \inferrule*[right=Inr]
             {\sexp \trans \texp}
             {\sinr{\sexp} \trans \tinr{\texp}}
  %
  \and
  %
  \inferrule*[right=Case]
             {\sexp \trans \texp \\
              \source{e_i}\ssubst{\svar{x}}{\salien{\tvar{x}}{\source{\tau_i}}{\cdot}} \trans \target{e_i}
             }
             {\scasex{\sexp} {\staui}{\sexpi} {\stauii}{\sexpii} \trans \tcasex{\texp} {\texpi} {\texpii}}
  %
  \and
  %
  \inferrule*[right=Abs]
             {\sexp\ssubst{\svar{x}}{\salien{\tvar{x}}{\source{\tau}}{\cdot}} \trans \texp}
             {\sabs{\svar{x}}{\stau}{\sexp} \trans \tabs{\tvar{x}}{\texp}}
  %
  \and
  %
  \inferrule*[right=App]
             {\source{e_i} \trans \target{e_i}}
             {\sapp{\source{e_1}}{\source{e_2}} \trans \tapp{\target{e_1}}{\target{e_2}}}
  %
  \and
  %
  \inferrule*[right=TAbs]
             {\sexp \trans \texp}
             {\stabs{\svar{A}}{\sexp} \trans \texp}
  %
  \and
  %
  \inferrule*[right=TApp]
             {\sexp \trans \texp}
             {\stapp{\sexp}{\stau} \trans \texp}

 %% %
 %%  \and
 %%  %
 %%  \inferrule*[right=Box]
 %%             {\texp \tevals \tval}
 %%             {\salien{\texp}{\stau}{\sdelta} \seval \salien{\tval}{\stau}{\sdelta}}
 %%   %
 %%  \and
 %%  %
 %%  \inferrule*[right=ProjAlien]
 %%             {\\}
 %%             {\salien{\tpair{\texp_{\target{1}}}{\texp_{\target{2}}}}{\sprod{\stau_{\source{1}}}{\stau_{\source{2}}}}{\sdelta}
 %%               \seval
 %%               \salien{\texp_{\target{1}}}{\stau_{\source{1}}}{\sdelta} }
 %%  %
 %%  \and
 %%  %
 %%  \inferrule*[right=CaseLAlien]
 %%             {\\}
 %%             {\scasex{\salien{\tinl{\texp}}{\ssum{\stau_{\source{1}}}{\stau_{\source{2}}}}{\sdelta}}{\stau_{\source{1}}}{\sexp_{\source{1}}}{\stau_{\source{2}}}{\sexp_{\source{2}}}
 %%               \seval
 %%               \sexp_{\source{1}} \ssubst{\svar{x}}{\salien{\texp}{\stau_{\source{1}}}{\sdelta}}}
 %%  %
 %%  \and
 %%  %
 %%  \inferrule*[right=AppAlien]
 %%             {\\}
 %%             {\sapp{\salien{\tabs{\tvar{x}}{\texp}}{\sarrow{\stau_{\source{1}}}{\stau_{\source{2}}}}{\sdelta}}
 %%               {\sexp}
 %%               \seval
 %%               \sabs{\svar{x}}{\sdelta\stau_{\source{1}}}
 %%                    {\kap{\stau_{\source{1}}}{\sdelta}{\svar{x}}
 %%                      {\mfun{\texpp}{\salien{\texp\ssubst{\tvar{x}}{\texpp}}{\stau_{\source{2}}}{\sdelta}}}}
 %%             }
 %%  %
 %%  \and
 %%  %
 %%  \inferrule*[right=TAppAlien]
 %%             {\\}
 %%             {\stapp{\salien{\texp}{\sall{\svar{A}}{\stau}}{\sdelta}}{\staup}
 %%               \seval
 %%               \salien{\texp}{\stau}{\sesubstext{\sdelta}{\svar{A}}{\staup}}
 %%             }
 %%  %
 %%  \and
 %%  %
 %%  \inferrule*[right=AlienOpaque]
 %%             {\stau \not = \sttop}
 %%             {\salien{\topaque{\sexp}}{\stau}{\sdelta}
 %%               \seval
 %%               \sexp
 %%             }
 %%  %
 %%  \and
 %%  %
 %%  \inferrule*[right=AlienCancel]
 %%             {\\}
 %%             {\salien{\talien{\sexp}{\staup}{\sdeltap}}{\stau}{\sdelta}
 %%               \seval
 %%               \sexp}
\end{mathpar}

  \end{framed}
  \caption{Translation}
\end{wrapfigure}

\medskip

\begin{wrapfigure}[htb]{L}{\textwidth}%
    \centering
      \centering
      \begin{framed}
        \[
\begin{array}{r l l}
  \scoerce{\stunit}{\sdelta}{\tunit} & = & \sunit
  \\[0.3em]
  %
  \scoerce{\sarrow{\staui}{\stauii}}{\sdelta}{\tabs{\tvar{x}}{\texp}} & = &
  \sabs{\svar{x}}{\sdelta\staui}{\salien{\texp\tsubst{\tvar{x}}{\prot{\svar{x}}{\staui}{\sdelta}}} {\stauii} {\sdelta}}
  \\[0.3em]
  %
  \scoerce{\sprod{\staui}{\stauii}}{\sdelta}{\tpair{\tvali}{\tvalii}} & = &
  \spair{\salien{\tvali}{\staui}{\sdelta}}{\salien{\tvalii}{\stauii}{\sdelta}}
  \\[0.3em]
  %
  \scoerce{\ssum{\staui}{\stauii}}{\sdelta}{\tinl{\tval}} & = &
  \sinl{\salien{\tval}{\staui}{\sdelta}}
  \\[0.3em]
  %
  \scoerce{\ssum{\staui}{\stauii}}{\sdelta}{\tinr{\tval}} & = &
  \sinr{\salien{\tval}{\stauii}{\sdelta}}
  \\[0.3em]
  % 
  \scoerce{\svar{A}}{\sdelta}{\topaque{\sexp}} & = &
  \sexp \\[0.3em]
  & & \\
  %
  \scoerce{\sall{\svar{A}}{\ssig}}{\sdelta}{\tval} & = &
  \stabs{\svar{A}}{\scoerce{\ssig}{\sdelta}{\tval}}
  \\[0.3em]

  

  %
  %
  %% \scoerce{\sttop}{\sdelta}{\salien{\tval}{\staup}{\sdeltap}} & = &
  %% \tval
  %% \\[0.3em]
\end{array}
\]

      \end{framed}
      \caption{Target to Source Coercion}
      \medskip\medskip\medskip\medskip\medskip
      \centering
      \begin{framed}
        \[
\begin{array}{r l l}
  \tcoerce{\stunit}{\sdelta}{\sunit} & = & \tunit
  \\[0.3em]
  %
  \tcoerce{\sarrow{\staui}{\stauii}}{\sdelta}{\sabs{\svar{x}}{\stau}{\sexp}} & = &
  \tabs{\tvar{x}}{\prot{\sexp\tsubst{\svar{x}}{\salien{\tvar{x}}{\staui}{\sdelta}}} {\stauii} {\sdelta}}
  \\[0.3em]
  %
  \tcoerce{\sprod{\staui}{\stauii}}{\sdelta}{\spair{\sexpi}{\sexpii}} & = &
  \tpair{\prot{\sexpi}{\staui}{\sdelta}}{\prot{\sexpii}{\stauii}{\sdelta}}
  \\[0.3em]
  %
  \tcoerce{\ssum{\staui}{\stauii}}{\sdelta}{\sinl{\sexp}} & = &
  \tinl{\prot{\sexp}{\staui}{\sdelta}}
  \\[0.3em]
  %
  \tcoerce{\ssum{\staui}{\stauii}}{\sdelta}{\sinr{\sexp}} & = &
  \tinr{\prot{\sexp}{\stauii}{\sdelta}}
  \\[0.3em]
  %
  %% \tcoerce{\sall{\svar{A}}{\stau}}{\sdelta}{\stabs{\svar{A}}{\sexp}} & = &
  %% \prot{\sexp\ssubst{\svar{A}}{\sttop}}{\stau\ssubst{\svar{A}}{\sttop} }{\sdelta}
  %% \\[0.3em]
  %% %
  \tcoerce{\svar{A}}{\sdelta}{\sexp} & = & \topaque{\sexp}
  \\[0.3em]
  \tcoerce{\staup}{\sdelta}{\salien{\texp}{\stau}{\sdelta}} & = & \texp
  \hfill \text{if } \stau \equiv \staup
  \\[0.3em]
\end{array}
\]

      \end{framed}
      \caption{Source to Target Coercion}
\end{wrapfigure}

\medskip

\begin{wrapfigure}[htb]{l}{\textwidth}%
  \label{fig:protect}
  %\fbox{$\prot{\sexp}{\stau}{\sdelta}$}
  \begin{framed}
    \[
\prot{\sexp}{\stau}{\sdelta} =
\begin{cases}
  \topaque{\sexp} & \text{if }\stau = \svar{A} \\
  \talien{\sexp}{\stau}{\sdelta} & \text{otherwise}
\end{cases}
\]

  \end{framed}
  \caption{Protect meta-function}
\end{wrapfigure}

\medskip

\begin{wrapfigure}[htb]{l}{\textwidth}%
  \label{fig:typeq}
  %\fbox{$\staui \equiv \stauii$}
  \begin{framed}
    \[
\begin{array}{r l l}
  \stunit \equiv \stunit & \eqdef & \top \\
  %
  \source{A} \equiv \source{B} & \eqdef & \top \\
  %
  \spair{\staui}{\stauii} \equiv
  \spair{\source{\tau_1'}}{\source{\tau_2'}}
  & \eqdef & \staui \equiv \source{\tau_1'} \wedge \stauii \equiv \source{\tau_2'}  \\
  %
  \sarrow{\staui}{\stauii} \equiv \sarrow{\source{\tau_1'}}{\source{\tau_2'}}
  & \eqdef & \staui \equiv \source{\tau_1'} \wedge \stauii \equiv \source{\tau_2'}  \\
  %
  \ssum{\staui}{\stauii} \equiv \ssum{\source{\tau_1'}}{\source{\tau_2'}}
  & \eqdef & \staui \equiv \source{\tau_1'} \wedge \stauii \equiv \source{\tau_2'}  \\
  %
  \stau \equiv \stau' 
  & \eqdef & \bot \hfill \text{otherwise}  \\

\end{array}
\]

  \end{framed}
  \caption{Type equivalence}
\end{wrapfigure}


\medskip

\begin{wrapfigure}[htb]{l}{\textwidth}%
  \label{fig:trans}
  \fbox{$\sexp \trans \texp$} \\
  %$\senv$ is a global parameter 
  \begin{framed}
    \begin{mathpar}
  \inferrule*[right=Unit]
             {\\}
             {\sunit \trans \tunit}
  %
  \and
  \inferrule*[right=Var]
             {\\}
             {\svar{x} \trans \talien{\svar{x}}{\senv\spar{\svar{x}}}{\cdot}}

  %
  \and
  %
  \inferrule*[right=Box]
             {\\}
             {\salien{\texp}{\stau}{\sdelta} \trans \texp}
  %
  \and
  %
  \inferrule*[right=Pair]
             {\source{e_i} \trans \target{e_i}}
             { \spair{\source{e_1}}{\source{e_2}} \trans \tpair{\target{e_1}}{\target{e_2}}}
  %
  \and
  %
  \inferrule*[right=Proj]
             {\sexp \trans \texp}
             {\sproji{\sexp} \trans \tproji{\texp}}
  %
  \and
  %
  \inferrule*[right=Inl]
             {\sexp \trans \texp}
             {\sinl{\sexp} \trans \tinl{\texp}}
  %
  \and
  %
  \inferrule*[right=Inr]
             {\sexp \trans \texp}
             {\sinr{\sexp} \trans \tinr{\texp}}
  %
  \and
  %
  \inferrule*[right=Case]
             {\sexp \trans \texp \\
              \source{e_i}\ssubst{\svar{x}}{\salien{\tvar{x}}{\source{\tau_i}}{\cdot}} \trans \target{e_i}
             }
             {\scasex{\sexp} {\staui}{\sexpi} {\stauii}{\sexpii} \trans \tcasex{\texp} {\texpi} {\texpii}}
  %
  \and
  %
  \inferrule*[right=Abs]
             {\sexp\ssubst{\svar{x}}{\salien{\tvar{x}}{\source{\tau}}{\cdot}} \trans \texp}
             {\sabs{\svar{x}}{\stau}{\sexp} \trans \tabs{\tvar{x}}{\texp}}
  %
  \and
  %
  \inferrule*[right=App]
             {\source{e_i} \trans \target{e_i}}
             {\sapp{\source{e_1}}{\source{e_2}} \trans \tapp{\target{e_1}}{\target{e_2}}}
  %
  \and
  %
  \inferrule*[right=TAbs]
             {\sexp \trans \texp}
             {\stabs{\svar{A}}{\sexp} \trans \texp}
  %
  \and
  %
  \inferrule*[right=TApp]
             {\sexp \trans \texp}
             {\stapp{\sexp}{\stau} \trans \texp}

 %% %
 %%  \and
 %%  %
 %%  \inferrule*[right=Box]
 %%             {\texp \tevals \tval}
 %%             {\salien{\texp}{\stau}{\sdelta} \seval \salien{\tval}{\stau}{\sdelta}}
 %%   %
 %%  \and
 %%  %
 %%  \inferrule*[right=ProjAlien]
 %%             {\\}
 %%             {\salien{\tpair{\texp_{\target{1}}}{\texp_{\target{2}}}}{\sprod{\stau_{\source{1}}}{\stau_{\source{2}}}}{\sdelta}
 %%               \seval
 %%               \salien{\texp_{\target{1}}}{\stau_{\source{1}}}{\sdelta} }
 %%  %
 %%  \and
 %%  %
 %%  \inferrule*[right=CaseLAlien]
 %%             {\\}
 %%             {\scasex{\salien{\tinl{\texp}}{\ssum{\stau_{\source{1}}}{\stau_{\source{2}}}}{\sdelta}}{\stau_{\source{1}}}{\sexp_{\source{1}}}{\stau_{\source{2}}}{\sexp_{\source{2}}}
 %%               \seval
 %%               \sexp_{\source{1}} \ssubst{\svar{x}}{\salien{\texp}{\stau_{\source{1}}}{\sdelta}}}
 %%  %
 %%  \and
 %%  %
 %%  \inferrule*[right=AppAlien]
 %%             {\\}
 %%             {\sapp{\salien{\tabs{\tvar{x}}{\texp}}{\sarrow{\stau_{\source{1}}}{\stau_{\source{2}}}}{\sdelta}}
 %%               {\sexp}
 %%               \seval
 %%               \sabs{\svar{x}}{\sdelta\stau_{\source{1}}}
 %%                    {\kap{\stau_{\source{1}}}{\sdelta}{\svar{x}}
 %%                      {\mfun{\texpp}{\salien{\texp\ssubst{\tvar{x}}{\texpp}}{\stau_{\source{2}}}{\sdelta}}}}
 %%             }
 %%  %
 %%  \and
 %%  %
 %%  \inferrule*[right=TAppAlien]
 %%             {\\}
 %%             {\stapp{\salien{\texp}{\sall{\svar{A}}{\stau}}{\sdelta}}{\staup}
 %%               \seval
 %%               \salien{\texp}{\stau}{\sesubstext{\sdelta}{\svar{A}}{\staup}}
 %%             }
 %%  %
 %%  \and
 %%  %
 %%  \inferrule*[right=AlienOpaque]
 %%             {\stau \not = \sttop}
 %%             {\salien{\topaque{\sexp}}{\stau}{\sdelta}
 %%               \seval
 %%               \sexp
 %%             }
 %%  %
 %%  \and
 %%  %
 %%  \inferrule*[right=AlienCancel]
 %%             {\\}
 %%             {\salien{\talien{\sexp}{\staup}{\sdeltap}}{\stau}{\sdelta}
 %%               \seval
 %%               \sexp}
\end{mathpar}

  \end{framed}
  \caption{Translation}
  \label{fig:translationlong}
\end{wrapfigure}

\fi

\bibliographystyle{abbrvnaturl}
\bibliography{fstar}

\end{document}
